% Lesson 29 — Grammar: Direct and indirect speech — Anglais Tle A — Cours signé OnBuch+
% Programme officiel MINESEC. Compiler avec : tectonic EN29-grammar-direct-and-indirect-speech.tex
\newcommand{\DOCMATIERE}{Anglais}
\newcommand{\DOCNIVEAU}{Tle A}
\newcommand{\DOCMODULE}{Chapter 3 : Environment, Well-being and Health}
\newcommand{\DOCLECON}{EN29}
\newcommand{\DOCTITRE}{Lesson 29 — Grammar: Direct and indirect speech}
% OnBuch+ — préambule « cours signés » ENRICHI (compilé avec tectonic / XeLaTeX)
% Système visuel « Pop » d'OnBuch+ : fond crème (jamais blanc), bordures encre
% épaisses, ombres « dures » décalées, titres Archivo Black en capitales.
% Version MATHÉMATIQUES : graphes (pgfplots/tikz), boîtes pédagogiques
% (définition, propriété, méthode, attention, le savais-tu, exercice résolu,
% exercice, TP, bilan), page de garde et signature OnBuch+ ancrée en bas.
%
% Macros attendues dans main.tex AVANT \input{preamble} :
%   \DOCMATIERE  \DOCNIVEAU  \DOCMODULE  \DOCLECON (numéro)  \DOCTITRE
\documentclass[11pt]{article}

\usepackage{fontspec}
\usepackage[english,french]{babel} % français = langue principale ; l'anglais via \eng{..} / textebox
\usepackage[a4paper,margin=2cm,top=3.4cm,bottom=2.8cm,headheight=1.4cm,footskip=1.3cm]{geometry}
\usepackage{amsmath,amssymb,amsfonts}
\usepackage{mathtools} % \xmapsto, \coloneqq... souvent utilisés par les modèles
\usepackage{cancel} % simplifications barrées (bilans de forces)
% \abs{z} (module, valeur absolue) et \norm{u} (norme), utilisés par les modèles
\providecommand{\abs}{}\let\abs\relax\DeclarePairedDelimiter{\abs}{\lvert}{\rvert}
\providecommand{\norm}{}\let\norm\relax\DeclarePairedDelimiter{\norm}{\lVert}{\rVert}
\usepackage[table]{xcolor} % [table] : fournit aussi \rowcolors (alternance de lignes)
\usepackage{pagecolor}
\usepackage{tikz}
\usetikzlibrary{arrows.meta,positioning,calc,shapes.geometric,shapes.misc,shapes.arrows,decorations.pathmorphing,decorations.pathreplacing,decorations.markings,patterns,shadows,mindmap,trees,fit,backgrounds,matrix,intersections,angles,quotes}
\usepackage{pgfplots}
\usepackage[version=4]{mhchem} % équations nucléaires \ce{^{235}_{92}U}
\usepackage[european,siunitx]{circuitikz} % schémas électriques
\pgfplotsset{compat=1.18}
\usepgfplotslibrary{fillbetween,groupplots} % aires entre courbes (calcul intégral)
\usepackage{enumitem}
\usepackage{fancyhdr}
% [most] charge toutes les bibliothèques tcolorbox (listings, minted, raster,
% documentation...) et gonfle inutilement la mémoire TeX sur un document long
% (~300 boîtes) : on ne charge que ce qui est réellement utilisé.
\usepackage[skins,breakable]{tcolorbox}
\usepackage{graphicx}
\usepackage{colortbl}
\usepackage{booktabs}
\usepackage{tabularx}
\usepackage{diagbox} % case d'en-tête diagonale des tableaux à double entrée
% Colonnes extensibles centrée / gauche / droite (C, L, R), souvent utilisées par les modèles
\newcolumntype{C}{>{\centering\arraybackslash}X}
\newcolumntype{L}{>{\raggedright\arraybackslash}X}
\newcolumntype{R}{>{\raggedleft\arraybackslash}X}
\usepackage{array}
\usepackage{multicol}
\usepackage{multirow}
\usepackage{siunitx}
% parse-numbers=false : accepte aussi \SI{2,5 \times 10^{-3}}{...}
\sisetup{locale=FR,output-decimal-marker={,},group-minimum-digits=4,parse-numbers=false}
\DeclareSIUnit\ohm{\ensuremath{\symup{\Omega}}} % U+2126 absent de la police
\DeclareSIUnit\division{div} % réglages d'oscilloscope (ms/div, V/div)
\DeclareSIUnit\tr{tr} % tours (vitesse de rotation, ex. \SI{1500}{\tr\per\minute})
\DeclareSIUnit\molar{M} % concentration molaire (ex. \SI{3}{\molar}, \SI{1}{\micro\molar})
\DeclareSIUnit\an{an} % année (le modèle confond parfois avec \year, un compteur TeX) — ex. \SI{5}{\kilo\meter\per\an}
\DeclareSIUnit\FCFA{FCFA} % franc CFA (ex. \SI{500}{\FCFA\per\kilo\gram})
\DeclareSIUnit\degreeBrix{\textdegree Brix} % teneur en sucre d'un sirop/jus (réfractométrie)
\DeclareSIUnit\month{mois} % le modèle utilise parfois \month, un compteur TeX, comme unité de durée
\usepackage{qrcode}
% \degree : commande fréquemment utilisée par les modèles pour un angle ou une
% température en dehors de \SI{}{} (non fournie par les paquets chargés).
\providecommand{\degree}{^{\circ}}
% \qed (fin de démonstration), utilisé par les modèles sans amsthm
\providecommand{\qed}{\hfill$\square$}
% \barre : double barre des valeurs interdites dans un tableau de variations (tkz-tab)
\providecommand{\barre}{\ensuremath{\|}}
% environnement proof (amsthm non chargé) utilisé par les modèles
\ifdefined\proof\else\newenvironment{proof}[1][Démonstration]{\par\noindent\textit{#1.}\ }{\qed\par}\fi
\usepackage{lastpage}
\usepackage{hyperref}
\hypersetup{hidelinks}

% ============================================================
% Palette « Pop » OnBuch+ — mêmes hex que la classe P (plus_ui.dart)
% ============================================================
\definecolor{poporange}{HTML}{F59321}
\definecolor{popdark}{HTML}{E07A0C}
\definecolor{popcream}{HTML}{FFF6E8}
\definecolor{popink}{HTML}{1C1714}
\definecolor{poppurple}{HTML}{7A5AE0}
\definecolor{popgreen}{HTML}{2E9E6A}
\definecolor{poppink}{HTML}{E0457B}
\definecolor{popblue}{HTML}{2F7DE1}
\definecolor{popgold}{HTML}{C98A12}
\definecolor{popmuted}{HTML}{8A7F76}
\definecolor{popline}{HTML}{EADFD2}
\definecolor{popgray}{HTML}{8A7F76}
\colorlet{popgrey}{popgray}
% Alias tolérants (le modèle nomme parfois les couleurs différemment)
\colorlet{popwhite}{white}
\colorlet{popblack}{popink}
\colorlet{popred}{poppink}
\colorlet{popyellow}{popgold}
% Teintes claires (fonds de tableaux/encadrés) — le modèle nomme parfois
% les couleurs avec un suffixe L (ex. popblueL) sans les avoir définies.
\colorlet{poporangeL}{poporange!20}
\colorlet{popdarkL}{popdark!20}
\colorlet{poppurpleL}{poppurple!20}
\colorlet{popgreenL}{popgreen!20}
\colorlet{poppinkL}{poppink!20}
\colorlet{popblueL}{popblue!20}
\colorlet{popgoldL}{popgold!20}
\colorlet{popredL}{popred!20}
\colorlet{popgrayL}{popgray!20}
% Teintes claires pour fonds de tableaux / zones
\colorlet{popblueL}{popblue!10!white}
\colorlet{poporangeL}{poporange!14!white}
\colorlet{popgreenL}{popgreen!12!white}
\colorlet{poppurpleL}{poppurple!12!white}
\colorlet{poppinkL}{poppink!12!white}
\colorlet{popgoldL}{popgold!14!white}

\pagecolor{popcream}

% ============================================================
% Typographie
% ============================================================
\defaultfontfeatures{Ligatures=TeX}
\setmainfont{PlusJakartaSans-Medium.ttf}[Path=../../../fonts/,BoldFont=PlusJakartaSans-Bold.ttf]
% Maths en sans-serif (Fira Math) avec chiffres et lettres droites de Plus
% Jakarta, pour que les formules se fondent dans le texte.
\usepackage[math-style=ISO,bold-style=ISO]{unicode-math}
\setmathfont{FiraMath-Regular.otf}[Path=../../../fonts/]
\setmathfont{PlusJakartaSans-Medium.ttf}[Path=../../../fonts/,range={up/num,up/{Latin,latin},\mathup}]
\usepackage{icomma} % virgule décimale sans espace en mode math
% Caractères mathématiques Unicode absents des polices de texte (Plus Jakarta,
% Archivo Black) : rendus par leur équivalent mathématique, en texte comme en
% formule — sinon ils s'affichent en carré vide (ex. « Arithmétique dans ℤ »).
\usepackage{newunicodechar}
\newunicodechar{ℝ}{\ensuremath{\mathbb{R}}}
\newunicodechar{ℤ}{\ensuremath{\mathbb{Z}}}
\newunicodechar{ℂ}{\ensuremath{\mathbb{C}}}
\newunicodechar{ℕ}{\ensuremath{\mathbb{N}}}
\newunicodechar{ℚ}{\ensuremath{\mathbb{Q}}}
\newunicodechar{𝔻}{\ensuremath{\mathbb{D}}}
\newunicodechar{α}{\ensuremath{\alpha}}
\newunicodechar{β}{\ensuremath{\beta}}
\newunicodechar{γ}{\ensuremath{\gamma}}
\newunicodechar{δ}{\ensuremath{\delta}}
\newunicodechar{ε}{\ensuremath{\varepsilon}}
\newunicodechar{θ}{\ensuremath{\theta}}
\newunicodechar{λ}{\ensuremath{\lambda}}
\newunicodechar{μ}{\ensuremath{\mu}}
\newunicodechar{σ}{\ensuremath{\sigma}}
\newunicodechar{τ}{\ensuremath{\tau}}
\newunicodechar{φ}{\ensuremath{\varphi}}
\newunicodechar{ψ}{\ensuremath{\psi}}
\newunicodechar{ω}{\ensuremath{\omega}}
\newunicodechar{Σ}{\ensuremath{\Sigma}}
% majuscules produites par \MakeUppercase dans les titres (α-aminés -> Α)
\newunicodechar{Α}{\ensuremath{\alpha}}
\newunicodechar{Β}{\ensuremath{\beta}}
\newunicodechar{Γ}{\ensuremath{\Gamma}}
\newunicodechar{Δ}{\ensuremath{\Delta}}
\newunicodechar{Ε}{\ensuremath{\varepsilon}}
\newunicodechar{Θ}{\ensuremath{\Theta}}
\newunicodechar{Λ}{\ensuremath{\Lambda}}
\newunicodechar{Μ}{\ensuremath{\mu}}
\newunicodechar{Τ}{\ensuremath{\tau}}
\newunicodechar{Φ}{\ensuremath{\Phi}}
\newunicodechar{Ψ}{\ensuremath{\Psi}}
\newunicodechar{Ω}{\ensuremath{\Omega}}
\newunicodechar{↦}{\ensuremath{\mapsto}}
\newunicodechar{∈}{\ensuremath{\in}}
\newunicodechar{∉}{\ensuremath{\notin}}
\newunicodechar{∧}{\ensuremath{\wedge}}
\newunicodechar{⇔}{\ensuremath{\Leftrightarrow}}
\newunicodechar{⟺}{\ensuremath{\Longleftrightarrow}}
\newunicodechar{⇒}{\ensuremath{\Rightarrow}}
\newunicodechar{⟹}{\ensuremath{\Longrightarrow}}
\newunicodechar{∘}{\ensuremath{\circ}}
\newunicodechar{∪}{\ensuremath{\cup}}
\newunicodechar{∩}{\ensuremath{\cap}}
\newunicodechar{≡}{\ensuremath{\equiv}}
\newunicodechar{⊂}{\ensuremath{\subset}}
\newunicodechar{⊥}{\ensuremath{\perp}}
\newunicodechar{∃}{\ensuremath{\exists}}
\newunicodechar{∀}{\ensuremath{\forall}}
\newunicodechar{∛}{\ensuremath{\sqrt[3]{\,}}}
\newunicodechar{∜}{\ensuremath{\sqrt[4]{\,}}}
\newunicodechar{⃗}{\ensuremath{{}^{\to}}}
\newunicodechar{ⁿ}{\ifmmode^{n}\else\textsuperscript{\ensuremath{n}}\fi}
\newunicodechar{ˣ}{\ifmmode^{x}\else\textsuperscript{\ensuremath{x}}\fi}
\newunicodechar{ᵇ}{\ifmmode^{b}\else\textsuperscript{\ensuremath{b}}\fi}
\newunicodechar{ʳ}{\ifmmode^{r}\else\textsuperscript{\ensuremath{r}}\fi}
\newunicodechar{ᵘ}{\ifmmode^{u}\else\textsuperscript{\ensuremath{u}}\fi}
\newunicodechar{ʸ}{\ifmmode^{y}\else\textsuperscript{\ensuremath{y}}\fi}
\newunicodechar{ᵃ}{\ifmmode^{a}\else\textsuperscript{\ensuremath{a}}\fi}
\newunicodechar{ᵉ}{\ifmmode^{e}\else\textsuperscript{\ensuremath{e}}\fi}
\newunicodechar{ᵏ}{\ifmmode^{k}\else\textsuperscript{\ensuremath{k}}\fi}
\newunicodechar{ᵖ}{\ifmmode^{p}\else\textsuperscript{\ensuremath{p}}\fi}
\newunicodechar{ᴺ}{\ifmmode^{N}\else\textsuperscript{\ensuremath{N}}\fi}
\newunicodechar{ᶜ}{\ifmmode^{c}\else\textsuperscript{\ensuremath{c}}\fi}
\newunicodechar{ʰ}{\ifmmode^{h}\else\textsuperscript{\ensuremath{h}}\fi}
\newunicodechar{ᵠ}{\ifmmode^{q}\else\textsuperscript{\ensuremath{q}}\fi}
\newunicodechar{ₙ}{\ifmmode_{n}\else\textsubscript{\ensuremath{n}}\fi}
\newunicodechar{ₐ}{\ifmmode_{a}\else\textsubscript{\ensuremath{a}}\fi}
\newunicodechar{ᵢ}{\ifmmode_{i}\else\textsubscript{\ensuremath{i}}\fi}
\newunicodechar{ᵣ}{\ifmmode_{r}\else\textsubscript{\ensuremath{r}}\fi}
\newunicodechar{ₖ}{\ifmmode_{k}\else\textsubscript{\ensuremath{k}}\fi}
\newunicodechar{ᵦ}{\ifmmode_{\beta}\else\textsubscript{\ensuremath{\beta}}\fi}
\newunicodechar{ₓ}{\ifmmode_{x}\else\textsubscript{\ensuremath{x}}\fi}
\newfontfamily\popdisplay{ArchivoBlack-Regular.ttf}[Path=../../../fonts/]
\newfontfamily\poplogo{SpaceGrotesk-Bold.ttf}[Path=../../../fonts/]
\newfontfamily\popsemi{PlusJakartaSans-SemiBold.ttf}[Path=../../../fonts/]
\newfontfamily\popxbold{PlusJakartaSans-ExtraBold.ttf}[Path=../../../fonts/]
\newfontfamily\popxboldls{PlusJakartaSans-ExtraBold.ttf}[Path=../../../fonts/,LetterSpace=3.0]

\setlist[enumerate]{leftmargin=1.7em,itemsep=5pt,topsep=4pt}
\setlist[itemize]{leftmargin=1.7em,itemsep=4pt,topsep=4pt,label={\color{poporange}\textbullet}}
\setlength{\parindent}{0pt}
\setlength{\parskip}{5pt}
\renewcommand{\arraystretch}{1.25}

% pgfplots : style Pop par défaut
\pgfplotsset{
  every axis/.append style={axis line style={popink,line width=1pt},tick style={popink},
    grid style={popline,line width=0.5pt},label style={font=\small},tick label style={font=\footnotesize}},
  popcurve/.style={poporange,line width=1.8pt,no markers,smooth},
}
% Même style disponible dans un \draw TikZ simple (hors pgfplots)
\tikzset{popcurve/.style={poporange,line width=1.8pt}}
\tikzset{popgrid/.style={popline,very thin}}
\colorlet{popcurve}{poporange} % le nom du style est parfois utilisé comme couleur

% ============================================================
% Pill / squiggle
% ============================================================
\newcommand{\poppill}[3]{%
  \tikz[baseline=-0.55ex]\node[draw=popink,line width=1.2pt,rounded corners=2.6pt,%
    fill=#2,inner xsep=6.5pt,inner ysep=3pt]{%
    \popxboldls\fontsize{7}{7}\selectfont\color{#3}\MakeUppercase{#1}};%
}
\newcommand{\popsquiggle}[1]{%
  \tikz[baseline=-0.4ex]{
    \draw[#1,line width=2.6pt,line cap=round]
      plot [smooth] coordinates {(0,0.09) (0.34,0.01) (0.68,0.21) (1.02,0.01) (1.36,0.10)};
  }%
}
% Mot-clé surligné (marqueur orange sous le texte)
\newcommand{\cle}[1]{{\popxbold\color{popdark}#1}}

% ============================================================
% En-tête / pied
% ============================================================
\pagestyle{fancy}
\fancyhf{}
\renewcommand{\headrulewidth}{0pt}
\renewcommand{\footrulewidth}{0pt}
\lhead{\raisebox{-0.15ex}{\poplogo\fontsize{13}{13}\selectfont\color{popink}OnBuch\color{poporange}+}}
\rhead{\poppill{\DOCMATIERE\ \textbullet\ \DOCNIVEAU\ \textbullet\ Leçon \DOCLECON}{white}{popink}}
\lfoot{\popxbold\fontsize{7.6}{7.6}\selectfont\color{popmuted}\MakeUppercase{Cours signé OnBuch+}\ \textbullet\ {\popsemi\DOCTITRE}}
\rfoot{\tikz[baseline=-0.6ex]\node[rounded corners=8.5pt,fill=popink,minimum height=17pt,minimum width=17pt,inner xsep=5pt,inner sep=0]{\popxbold\fontsize{8}{8}\selectfont\color{popcream}\thepage\,/\,\pageref*{LastPage}};}

% ============================================================
% Titres
% ============================================================
\newcommand{\chaptitre}[2]{%
  \begin{tcolorbox}[enhanced,colback=poporange,colframe=popink,boxrule=2.6pt,arc=11pt,
      drop shadow=popink,left=15pt,right=15pt,top=12pt,bottom=11pt,before skip=3pt,after skip=18pt]
    \poppill{#2}{popink}{popcream}\par\vspace{9pt}
    {\raggedright\hyphenpenalty=10000\exhyphenpenalty=10000\popdisplay\fontsize{22}{24}\selectfont\color{popcream}\MakeUppercase{#1}\par}
  \end{tcolorbox}
}
% Section numérotée : pastille orange avec le numéro + titre Archivo
\newcounter{popsec}
\newcommand{\coursec}[1]{%
  \stepcounter{popsec}\setcounter{popsub}{0}%
  \par\vspace{16pt}\noindent
  \begin{minipage}{\linewidth}
  \tikz[baseline=-0.7ex]\node[draw=popink,line width=1.6pt,fill=poporange,rounded corners=4pt,minimum size=22pt,inner sep=0]{\popdisplay\fontsize{12}{12}\selectfont\color{popcream}\thepopsec};%
  \hspace{8pt}{\popdisplay\fontsize{15}{17}\selectfont\color{popink}\MakeUppercase{#1}}\par
  \vspace{4pt}\noindent{\color{poporange}\rule{36pt}{3.6pt}}
  \end{minipage}\par\nopagebreak\vspace{8pt}
}
\newcounter{popsub}
\newcommand{\courssub}[1]{%
  \stepcounter{popsub}%
  \par\vspace{9pt}\noindent{\popxbold\fontsize{11.5}{13}\selectfont\color{popink}{\color{poporange}\thepopsec.\thepopsub}\hspace{6pt}#1}\par\nopagebreak\vspace{4pt}
}

% ============================================================
% Boîtes pédagogiques — même recette : fond blanc, bordure épaisse colorée,
% ombre dure encre, badge plein. Titre optionnel [#1].
% ============================================================
\tcbset{popbase/.style={breakable,enhanced,colback=white,boxrule=2.3pt,arc=9pt,
  shadow={3pt}{-3pt}{0pt}{popink},left=14pt,right=14pt,top=11pt,bottom=11pt,before skip=11pt,after skip=15pt,
  fontupper=\popsemi\fontsize{10.6}{14.5}\selectfont\color{popink},
  fontlower=\popsemi\fontsize{10.6}{14.5}\selectfont\color{popink}}}
% Coche dessinée (le glyphe ✓ n'existe pas dans les polices du document)
\newcommand{\popcheck}[1]{\tikz[baseline=-0.5ex]\draw[#1,line width=1.6pt,line cap=round,line join=round](-0.13,0.0)--(-0.04,-0.09)--(0.13,0.1);}
\providecommand{\coche}{\popcheck{popgreen}}
\providecommand{\resultat}[1]{\par\smallskip\noindent\fcolorbox{popgreen}{popgreenL}{\parbox{\dimexpr\linewidth-2\fboxsep-2\fboxrule\relax}{\textbf{Résultat.} #1}}\par\smallskip}
\newcommand{\popboxhead}[3]{\poppill{#1}{#2}{white}\ifx\relax#3\relax\else\hspace{6pt}{\popxbold\fontsize{10.6}{12}\selectfont\color{popink}#3}\fi\par\vspace{7pt}}

\newtcolorbox{aretenir}[1][]{popbase,colframe=popgreen,before upper={\popboxhead{À retenir}{popgreen}{#1}}}
% Texte en anglais (document de lecture, dialogue, extrait) : cadre
% d'étude. Un titre optionnel (source, auteur) s'affiche dans l'en-tête.
\newtcolorbox{textebox}[1][]{popbase,colframe=popblue,before upper={\popboxhead{Text}{popblue}{#1}\begin{otherlanguage}{english}},after upper={\end{otherlanguage}}}
% Marques ✓ / ✗ (forme correcte / fautive) souvent utilisées par les modèles
\providecommand{\cross}{\textcolor{poppink}{\ensuremath{\times}}}
\providecommand{\xmark}{\cross}\providecommand{\crossmark}{\cross}\providecommand{\cmark}{\ensuremath{\checkmark}}
% Ligne de réponse / justification à compléter (inventée par les modèles)
\providecommand{\justification}[1]{\par\noindent\textit{Justification~:} \dotfill\par}
% \textipa (package tipa non chargé) : la police ne couvre pas l'API -> simple passage
\providecommand{\textipa}[1]{#1}
% Soulignement (ulem non chargé) : \uline, \ul
\providecommand{\uline}[1]{\underline{#1}}\providecommand{\ul}[1]{\underline{#1}}
% \glossaire{\item ...} inventée par les modèles : liste de glossaire
\providecommand{\glossaire}[1]{\begin{itemize}#1\end{itemize}}
% \alert (beamer) inventée par les modèles : texte mis en évidence
\providecommand{\alert}[1]{\textbf{\textcolor{poppink}{#1}}}
% variantes ASCII d'ordinaux écrites en macro
\providecommand{\iere}{\textsuperscript{re}}\providecommand{\ieme}{\textsuperscript{e}}\providecommand{\ere}{\textsuperscript{re}}\providecommand{\eme}{\textsuperscript{e}}\providecommand{\er}{\textsuperscript{er}}\providecommand{\ie}{\textsuperscript{e}}
% Ordinaux français écrits en macro par les modèles : 1\ière, 3\ième
\newcommand{\ière}{\textsuperscript{re}}\newcommand{\ième}{\textsuperscript{e}}
% \exemple (inventée par les modèles) : étiquette « Exemple : »
\providecommand{\exemple}{\textbf{Exemple~:}\ }
% \square utilisable aussi hors mode math (cases à cocher)
\AtBeginDocument{\let\squareMath\square\renewcommand{\square}{\ensuremath{\squareMath}}}
% Mot ou phrase en anglais à l'intérieur d'un paragraphe français
\newcommand{\eng}[1]{\ifmmode\text{\textit{#1}}\else\begingroup\selectlanguage{english}\itshape #1\endgroup\fi}
\let\esp\eng % alias toléré
\newtcolorbox{exemplebox}[1][]{popbase,colframe=poporange,before upper={\popboxhead{Exemple}{poporange}{#1}}}
\newtcolorbox{definition}[1][]{popbase,colframe=popblue,before upper={\popboxhead{Définition}{popblue}{#1}}}
\newtcolorbox{propriete}[1][]{popbase,colframe=poppurple,before upper={\popboxhead{Propriété}{poppurple}{#1}}}
\newtcolorbox{methode}[1][]{popbase,colframe=popgold,before upper={\popboxhead{Méthode}{popgold}{#1}}}
\newtcolorbox{attention}[1][]{popbase,colframe=poppink,before upper={\popboxhead{Attention}{poppink}{#1}}}
\newtcolorbox{experience}[1][]{popbase,colframe=popblue,colback=popblueL,before upper={\popboxhead{Expérience}{popblue}{#1}}}
% « Le savais-tu ? » : carte violette pleine (contexte, culture, Cameroun)
\newtcolorbox{savaistu}[1][]{popbase,colback=poppurple,colframe=popink,
  fontupper=\popsemi\fontsize{10.4}{14.2}\selectfont\color{white},
  before upper={\poppill{Le savais-tu\,?}{white}{poppurple}\ifx\relax#1\relax\else\hspace{6pt}{\popxbold\color{white}#1}\fi\par\vspace{7pt}}}
% Exercice résolu : énoncé en haut, \tcblower puis la solution sur fond crème
\newcounter{popexo}\newcounter{popexr}
\newtcolorbox{exoresolu}[1][]{popbase,colframe=poporange,colbacklower=popcream,
  segmentation style={popink,line width=1.2pt,dashed},
  before upper={\stepcounter{popexr}\popboxhead{Exercice résolu \thepopexr}{poporange}{#1}},
  before lower={\poppill{Solution}{popink}{popcream}\par\vspace{6pt}}}
% Exercice (non résolu en place) — la correction est dans la partie Corrigés
\newtcolorbox{exercice}[1][]{popbase,colframe=popink,
  before upper={\stepcounter{popexo}\popboxhead{Exercice \thepopexo}{popink}{#1}}}
% Corrigé
\newtcolorbox{corrige}[1][]{popbase,colframe=popgreen,colback=popgreenL,
  before upper={\popboxhead{Corrigé}{popgreen}{#1}}}
% Objectifs / prérequis / bilan
\newtcolorbox{objectifs}{popbase,colframe=popink,colback=poporangeL,
  before upper={\popboxhead{Objectifs de la leçon}{popink}{}}}
\newtcolorbox{prerequis}{popbase,colframe=popink,colback=popblueL,
  before upper={\popboxhead{Prérequis}{popblue}{}}}
\newtcolorbox{bilan}{popbase,colframe=popink,colback=white,boxrule=2.8pt,
  before upper={\popboxhead{Fiche bilan}{poporange}{L'essentiel à connaître pour l'examen}}}
% Figure encadrée avec légende
\newtcolorbox{popfigure}[1][]{enhanced,breakable=false,colback=white,colframe=popline,boxrule=1.4pt,arc=8pt,
  left=10pt,right=10pt,top=10pt,bottom=8pt,before skip=10pt,after skip=12pt,halign=center,
  lower separated=false,halign lower=center,
  fontlower=\popsemi\fontsize{9}{11.5}\selectfont\color{popmuted},
  #1}
\newcommand{\legende}[1]{\par\vspace{6pt}{\popsemi\fontsize{9}{11.5}\selectfont\color{popmuted}#1}}

% ============================================================
% Signature OnBuch+
% ============================================================
\newcommand{\poplogoinline}[1]{{\poplogo\fontsize{#1}{#1}\selectfont\color{popink}OnBuch\color{poporange}+}}
\newcommand{\signatureonbuch}{%
  \begin{tcolorbox}[enhanced,colback=popink,colframe=popink,boxrule=2.2pt,arc=10pt,
      drop shadow=poporange,left=14pt,right=14pt,top=10pt,bottom=10pt,before skip=0pt,after skip=0pt]
    \tikz[baseline=-0.7ex]\node[circle,fill=popgreen,draw=popcream,line width=1.3pt,minimum size=22pt,inner sep=0]{\popcheck{white}};%
    \hspace{9pt}\begin{minipage}[c]{0.62\linewidth}
      {\popxbold\fontsize{12}{13}\selectfont\color{popcream}Signé\ \textbullet\ {\poplogo OnBuch\color{poporange}+}}\par\vspace{2pt}
      {\raggedright\popsemi\fontsize{8}{10.5}\selectfont\color{popline}\MakeUppercase{Équipe pédagogique OnBuch+}\\\MakeUppercase{Conforme au programme officiel MINESEC}\par}
    \end{minipage}\hfill
    {\poplogo\fontsize{22}{22}\selectfont\color{popcream}OnBuch\color{poporange}+}
  \end{tcolorbox}%
}

% ============================================================
% Page de garde
% #1 titre · #2 matière · #3 niveau · #4 module · #5 n° leçon · #6 durée ·
% #7 description / objectifs (texte ou itemize)
% ============================================================
\newcommand{\pagegarde}[7]{%
  \thispagestyle{empty}
  \newgeometry{a4paper,margin=1.7cm,top=1.6cm,bottom=1.4cm}%
  \begin{tikzpicture}[remember picture,overlay]
    \fill[poporange!18] ([xshift=-3.2cm,yshift=-3.6cm]current page.north east) circle (4.6cm);
    \fill[poppurple!14] ([xshift=2.4cm,yshift=7.6cm]current page.south west) circle (3.8cm);
  \end{tikzpicture}%
  {\poplogo\fontsize{30}{30}\selectfont\color{popink}OnBuch\color{poporange}+}\hfill
  \raisebox{0.6ex}{\poppill{Cours signé \textbullet\ Premium}{popink}{popcream}}\par
  \vspace{1pt}\popsquiggle{poppurple}\par
  \vspace{8pt}
  \begin{tcolorbox}[enhanced,colback=poporange,colframe=popink,boxrule=2.8pt,arc=13pt,
      drop shadow=popink,left=18pt,right=18pt,top=12pt,bottom=13pt,after skip=10pt]
    \poppill{#2 \textbullet\ #3}{popink}{popcream}\hspace{5pt}\poppill{#4}{white}{popink}\par\vspace{8pt}
    {\popdisplay\fontsize{14}{14}\selectfont\color{popink}LEÇON #5}\par\vspace{4pt}
    {\raggedright\hyphenpenalty=10000\exhyphenpenalty=10000\popdisplay\fontsize{25}{27}\selectfont\color{popcream}\MakeUppercase{#1}\par}
    \vspace{8pt}
    {\popxbold\fontsize{9.5}{9.5}\selectfont\color{popink}\MakeUppercase{Durée conseillée : #6}}
  \end{tcolorbox}
  \begin{tcolorbox}[enhanced,colback=white,colframe=popink,boxrule=2.2pt,arc=10pt,
      drop shadow=popink,left=16pt,right=16pt,top=10pt,bottom=10pt,after skip=8pt]
    \poppill{Dans cette leçon}{poppurple}{white}\par\vspace{5pt}
    \popsemi\fontsize{9.3}{12.6}\selectfont\color{popink}#7
  \end{tcolorbox}
  \noindent\begin{minipage}[t]{0.487\linewidth}
    \begin{tcolorbox}[enhanced,colback=popgreen,colframe=popink,boxrule=2.2pt,arc=10pt,
        drop shadow=popink,left=11pt,right=11pt,top=10pt,bottom=10pt,height=3.1cm,valign=center]
      \tcbox[colback=white,colframe=popink,boxrule=1.3pt,arc=5pt,boxsep=0pt,nobeforeafter,
        left=3pt,right=3pt,top=3pt,bottom=3pt,valign=center]{\qrcode[height=1.9cm]{https://play.google.com/store/apps/details?id=com.luvvix.onbuch&hl=fr}}%
      \hspace{8pt}\begin{minipage}[c]{0.5\linewidth}
        {\popdisplay\fontsize{11}{12.5}\selectfont\color{white}\MakeUppercase{Télécharge OnBuch}}\par\vspace{4pt}
        {\raggedright\popsemi\fontsize{8.4}{11}\selectfont\color{white}Gratuit \textbullet\ Google Play\\Tuteur IA, annales, concours, résultats\par}
      \end{minipage}
    \end{tcolorbox}
  \end{minipage}\hfill
  \begin{minipage}[t]{0.487\linewidth}
    \begin{tcolorbox}[enhanced,colback=poppurple,colframe=popink,boxrule=2.2pt,arc=10pt,
        drop shadow=popink,left=11pt,right=11pt,top=10pt,bottom=10pt,height=3.1cm,valign=center]
      \tikz[baseline=-0.7ex]\node[circle,fill=white,draw=popink,line width=1.3pt,minimum size=1.6cm,inner sep=0]{\popdisplay\fontsize{24}{24}\selectfont\color{poppurple}+};%
      \hspace{8pt}\begin{minipage}[c]{0.6\linewidth}
        {\popdisplay\fontsize{11}{12.5}\selectfont\color{white}\MakeUppercase{Passe à OnBuch+}}\par\vspace{4pt}
        {\raggedright\popsemi\fontsize{8.4}{11}\selectfont\color{white}Tous les cours signés, Tuteur IA illimité, fiches PDF — onglet OnBuch+ de l'app\par}
      \end{minipage}
    \end{tcolorbox}
  \end{minipage}
  \vspace{8pt}
  \popcontenu
  \vfill
  \signatureonbuch
  \restoregeometry
  \newpage
}

% Rangée « Ce cours contient » (page de garde)
\newcommand{\poptile}[3]{% #1 couleur #2 gros texte #3 légende
  \begin{tcolorbox}[enhanced,colback=#1,colframe=popink,boxrule=2pt,arc=9pt,shadow={2.5pt}{-2.5pt}{0pt}{popink},
      left=6pt,right=6pt,top=9pt,bottom=9pt,halign=center,valign=center,height=1.75cm,width=\linewidth,nobeforeafter]
    {\popdisplay\fontsize{12.5}{13}\selectfont\color{white}\MakeUppercase{#2}}\par\vspace{3pt}
    {\centering\popsemi\fontsize{7.8}{9.5}\selectfont\color{white}#3\par}
  \end{tcolorbox}}
\newcommand{\popcontenu}{%
  \par\noindent{\popxboldls\fontsize{7.5}{7.5}\selectfont\color{popmuted}CE COURS SIGNÉ CONTIENT}\par\vspace{7pt}
  \noindent\begin{minipage}[t]{0.235\linewidth}\poptile{popblue}{Cours}{complet, illustré, pas à pas}\end{minipage}\hfill
  \begin{minipage}[t]{0.235\linewidth}\poptile{poporange}{Méthodes}{et exercices résolus}\end{minipage}\hfill
  \begin{minipage}[t]{0.235\linewidth}\poptile{poppink}{Exercices}{type examen + corrigés}\end{minipage}\hfill
  \begin{minipage}[t]{0.235\linewidth}\poptile{popgreen}{Fiche bilan}{l'essentiel à retenir}\end{minipage}\par}

% Sommaire
\newtcolorbox{sommaire}{enhanced,colback=white,colframe=popink,boxrule=2.2pt,arc=10pt,
  drop shadow=popink,left=18pt,right=18pt,top=13pt,bottom=13pt,before skip=6pt,after skip=20pt,
  before upper={\poppill{Sommaire}{poppurple}{white}\par\vspace{9pt}\popsemi\fontsize{10.6}{14.5}\selectfont\color{popink}}}

% Signature de fin de cours — poussée tout en bas de la dernière page
\newcommand{\findecours}{%
  \par\vfill\nopagebreak
  \begin{center}\popsquiggle{poporange}\end{center}\vspace{4pt}
  \signatureonbuch
}
\AtBeginDocument{\def\llbracket{\mathopen{[\mkern-3.5mu[}}\def\rrbracket{\mathclose{]\mkern-3.5mu]}}} % intervalles d'entiers (glyphe ⟦ absent de la police)
% Glyphes absents de Fira Math : redéfinis à partir de glyphes disponibles
\AtBeginDocument{%
  \def\setminus{\mathbin{\backslash}}\let\smallsetminus\setminus
  \def\vdots{\mathord{\vbox{\baselineskip3.2pt\lineskiplimit0pt\kern4pt\hbox{.}\hbox{.}\hbox{.}}}}%
  \def\ddots{\mathinner{\mkern1mu\raise6.5pt\hbox{.}\mkern2mu\raise3.6pt\hbox{.}\mkern2mu\raise0.7pt\hbox{.}\mkern1mu}}%
  \def\triangle{\mathord{\scalebox{0.9}{$\symup{\Delta}$}}}%
  \def\nabla{\mathord{\raisebox{\depth}{\scalebox{1}[-1]{$\symup{\Delta}$}}}}%
  \def\checkmark{\popcheck{popdark}}%
  \def\star{\mathbin{\ast}}\def\bigstar{\mathord{\ast}}%
  \def\diamond{\mathbin{\scalebox{0.6}{$\Diamond$}}}%
  \def\bigcup{\mathop{\mathchoice{\raisebox{-0.3em}{\scalebox{1.7}{$\cup$}}}{\scalebox{1.25}{$\cup$}}{\cup}{\cup}}\displaylimits}%
  \def\bigcap{\mathop{\mathchoice{\raisebox{-0.3em}{\scalebox{1.7}{$\cap$}}}{\scalebox{1.25}{$\cap$}}{\cap}{\cap}}\displaylimits}%
  \def\lesseqqgtr{\mathrel{\lessgtr}}\def\bigoplus{\mathop{\oplus}\displaylimits}%
}
\providecommand{\ssi}{si et seulement si} % abréviation fréquente
\AtBeginDocument{\providecommand{\wideparen}{\overparen}} % arc de cercle

\begin{document}

\pagegarde{Lesson 29 — Grammar: Direct and indirect speech}{Anglais}{Tle A}{Chapter 3 : Environment, Well-being and Health}{EN29}{2 h}{Cette leçon de grammaire enseigne comment rapporter les paroles de quelqu'un en anglais : le discours direct (citation exacte) et le discours indirect (paroles rapportées). Les élèves apprennent les changements de temps, de pronoms et d'expressions de temps et de lieu, ainsi que les verbes introducteurs adaptés aux déclarations, questions et ordres.
\vspace{4pt}
\begin{multicols}{2}\small
\begin{itemize}
\item À la fin de cette leçon, je sais distinguer le discours direct du discours indirect
\item je sais identifier le verbe introducteur (say, tell, ask, order, advise...) et l'utiliser correctement
\item je sais reculer les temps (backshift) quand le verbe introducteur est au passé
\item je sais transformer les pronoms personnels et possessifs selon le contexte
\item je sais transformer les expressions de temps et de lieu (now → then, today → that day, here → there...)
\item je sais rapporter une question fermée avec 'if/whether' et une question ouverte avec le mot interrogatif
\end{itemize}\end{multicols}}

\begin{sommaire}
\begin{multicols}{2}
\begin{enumerate}
\item Observation : direct speech vs indirect speech
\item Le backshift : recul des temps
\item Pronoms et expressions de temps et de lieu
\item Rapporter questions, ordres et conseils
\item Verbes introducteurs et nuances de sens
\item Comparaison français/anglais et entraînement guidé
\item Activité d'intégration
\item Méthodes et exercices résolus
\item Exercices
\item Corrigés des exercices
\item Fiche bilan
\end{enumerate}
\end{multicols}
\end{sommaire}

\begin{objectifs}
\begin{itemize}
\item À la fin de cette leçon, je sais distinguer le discours direct du discours indirect
\item je sais identifier le verbe introducteur (say, tell, ask, order, advise...) et l'utiliser correctement
\item je sais reculer les temps (backshift) quand le verbe introducteur est au passé
\item je sais transformer les pronoms personnels et possessifs selon le contexte
\item je sais transformer les expressions de temps et de lieu (now → then, today → that day, here → there...)
\item je sais rapporter une question fermée avec 'if/whether' et une question ouverte avec le mot interrogatif
\item je sais rapporter un ordre ou un conseil avec la structure 'tell/ask/advise + objet + to-infinitive'
\item je sais éviter les erreurs typiques des francophones (say to me, que, oubli du backshift)
\item je sais utiliser le discours indirect dans un court texte narratif ou un compte rendu
\end{itemize}
\end{objectifs}

\begin{prerequis}
\begin{itemize}
\item Les temps principaux de l'anglais : present simple, present continuous, past simple, present perfect, future avec will (Première)
\item Les pronoms personnels sujets, objets et les adjectifs/pronoms possessifs
\item La structure de la phrase interrogative en anglais (inversion et auxiliaires do/does/did)
\item Les mots interrogatifs : what, where, when, why, who, how
\item L'infinitif avec 'to' après certains verbes (want, ask, tell...)
\end{itemize}
\end{prerequis}

\newpage

\coursec{Observation : direct speech vs indirect speech}

\courssub{Mise en situation : rapporter les paroles de l'oncle}

Pendant la réunion de famille à Yaoundé, ton oncle annonce : \eng{“I will send you to university in Buea next year.”} Le soir, ta petite sœur te demande ce qu'il a dit exactement. Tu dois rapporter ses paroles : \eng{“He said he would send me to university in Buea the following year.”} Mais attention : les temps changent (\eng{will} devient \eng{would}), les pronoms aussi (\eng{I} devient \eng{he}, \eng{you} devient \eng{me}), et même \eng{next year} devient \eng{the following year} !

\textbf{Problématique :} comment passer des paroles exactes prononcées par quelqu'un à un compte rendu fidèle de ces paroles, en respectant toutes les transformations grammaticales exigées par l'anglais ? Au Bac, rapporter correctement les paroles de quelqu'un est une compétence testée chaque année, dans les exercices de grammaire comme dans les rédactions (comptes rendus d'interview, articles de presse, lettres).

Dans cette première partie, nous allons d'abord \emph{observer} des exemples authentiques, puis dégager les deux notions fondamentales : le \cle{direct speech} et le \cle{indirect speech} (ou \cle{reported speech}).

\courssub{Texte d'observation : au lycée de Buea}

Lis attentivement ce court texte. Les passages en gras montrent les deux façons de rapporter des paroles.

\begin{textebox}[At Government High School Buea]
Last Monday, the English teacher was angry because many students had not done their homework. She entered the classroom and said, \textbf{“You must hand in your essays on Friday. I will not accept any late work.”} The students were worried. After the lesson, Divine told her friend Ngong: \textbf{“The teacher said we had to hand in our essays on Friday, and that she would not accept any late work.”}

In the afternoon, the principal met the students in the courtyard. He announced, \textbf{“I am going to close the library tomorrow for repairs.”} That evening, Ngong wrote in her diary: \textbf{“The principal announced that he was going to close the library the next day for repairs.”}

At home, Ngong's mother asked her, \textbf{“What did the principal say?”} Ngong replied, \textbf{“He said the library would be closed for repairs.”} Her mother smiled and said, \textbf{“Good. You can help me at the market then.”}
\end{textebox}

\textbf{Glossaire :} \emph{to hand in} (verbe à particule) : rendre (un devoir) ; \emph{late work} (nom) : travail rendu en retard ; \emph{courtyard} (nom) : cour ; \emph{repairs} (nom pluriel) : réparations ; \emph{diary} (nom) : journal intime ; \emph{to reply} (verbe) : répondre.

\textbf{Questions d'observation :}
\begin{enumerate}
\item Relevez toutes les phrases entre guillemets. Qui parle, et à qui ?
\item Comparez \eng{“I will not accept any late work”} avec \eng{“she would not accept any late work”}. Quels mots ont changé ?
\item Comparez \eng{“tomorrow”} et \eng{“the next day”}. Pourquoi ce changement ?
\end{enumerate}

\courssub{Les deux façons de rapporter des paroles}

\begin{definition}[Direct speech]
Le \cle{direct speech} (discours direct) consiste à citer les \textbf{paroles exactes} prononcées par quelqu'un. En anglais, ces paroles sont placées \textbf{entre guillemets} (“ … ”) et accompagnées d'un \textbf{verbe introducteur} (\eng{say}, \eng{tell}, \eng{ask}, \eng{announce}…).

\eng{The teacher said, “You must hand in your essays on Friday.”} « Le professeur a dit : “Vous devez rendre vos dissertations vendredi.” »
\end{definition}

\begin{definition}[Indirect / reported speech]
Le \cle{indirect speech} ou \cle{reported speech} (discours indirect) consiste à \textbf{rapporter le sens} des paroles de quelqu'un \textbf{sans les citer mot à mot}. Il n'y a \textbf{pas de guillemets} ; les paroles rapportées sont intégrées dans la phrase, souvent introduites par \cle{that} (que), avec des changements de temps, de pronoms et d'expressions de temps et de lieu.

\eng{The teacher said (that) we had to hand in our essays on Friday.} « Le professeur a dit que nous devions rendre nos dissertations vendredi. »
\end{definition}

\begin{popfigure}
\begin{tikzpicture}[
  grow=right,
  level distance=4.6cm,
  level 1/.style={sibling distance=2.4cm},
  every node/.style={draw, rounded corners, align=center, font=\small},
  edge from parent/.style={draw=popmuted, -{Stealth}, thick}
]
\node[fill=popblue, text=white, font=\bfseries, align=center]{Speech\\(paroles)}
  child { node[fill=poppinkL, align=center]{\textbf{Indirect speech}\\ reported, no quotes\\ changes of tense,\\ pronouns, time words} }
  child { node[fill=popblueL, align=center]{\textbf{Direct speech}\\ exact words\\ + quotation marks\\ “ … ”} };
\end{tikzpicture}
\legende{Les deux voies pour rapporter des paroles : citer exactement (direct) ou rapporter le sens (indirect).}
\end{popfigure}

\courssub{Observation guidée : quatre paires d'exemples en contexte camerounais}

Observons maintenant quatre paires de phrases. Dans chaque paire, la première phrase cite les paroles exactes (direct speech), la seconde les rapporte (indirect speech). Soulignons mentalement ce qui change.

\textbf{Paire 1 — Au marché de Mokolo (Yaoundé).}

\eng{The trader said, “I sell the best tomatoes in the market.”} « Le vendeur a dit : “Je vends les meilleures tomates du marché.” »

\eng{The trader said (that) he sold the best tomatoes in the market.} « Le vendeur a dit qu'il vendait les meilleures tomates du marché. »

Changements : \eng{I} $\rightarrow$ \eng{he} (pronom) ; \eng{sell} $\rightarrow$ \eng{sold} (temps : présent $\rightarrow$ prétérit).

\textbf{Paire 2 — À l'école de Bamenda.}

\eng{The headmaster announced, “The exams will start next week.”} « Le proviseur a annoncé : “Les examens commenceront la semaine prochaine.” »

\eng{The headmaster announced (that) the exams would start the following week.} « Le proviseur a annoncé que les examens commenceraient la semaine suivante. »

Changements : \eng{will start} $\rightarrow$ \eng{would start} (temps) ; \eng{next week} $\rightarrow$ \eng{the following week} (expression de temps).

\textbf{Paire 3 — En famille à Douala.}

\eng{My father said, “We are travelling to the village tomorrow.”} « Mon père a dit : “Nous partons au village demain.” »

\eng{My father said (that) they were travelling to the village the next day.} « Mon père a dit qu'ils partaient au village le lendemain. »

Changements : \eng{we} $\rightarrow$ \eng{they} (pronom) ; \eng{are travelling} $\rightarrow$ \eng{were travelling} (temps) ; \eng{tomorrow} $\rightarrow$ \eng{the next day} (expression de temps).

\textbf{Paire 4 — À l'hôpital de Buea.}

\eng{The nurse told me, “You have been here since this morning.”} « L'infirmière m'a dit : “Vous êtes ici depuis ce matin.” »

\eng{The nurse told me (that) I had been there since that morning.} « L'infirmière m'a dit que j'étais là depuis ce matin-là. »

Changements : \eng{you} $\rightarrow$ \eng{I} (pronom) ; \eng{have been} $\rightarrow$ \eng{had been} (temps) ; \eng{here} $\rightarrow$ \eng{there} (expression de lieu) ; \eng{this morning} $\rightarrow$ \eng{that morning} (expression de temps).

\begin{aretenir}[Les trois types de changements au discours indirect]
Quand on passe du discours direct au discours indirect (avec un verbe introducteur au passé), on observe \textbf{trois familles de changements} :
\begin{enumerate}
\item \textbf{Les temps verbaux} : ils « reculent » d'un cran vers le passé (\eng{sell} $\rightarrow$ \eng{sold}, \eng{will} $\rightarrow$ \eng{would}, \eng{have been} $\rightarrow$ \eng{had been}). C'est le \cle{backshift}.
\item \textbf{Les pronoms et possessifs} : ils s'adaptent au nouveau point de vue (\eng{I} $\rightarrow$ \eng{he/she}, \eng{you} $\rightarrow$ \eng{I/we}, \eng{we} $\rightarrow$ \eng{they}, \eng{my} $\rightarrow$ \eng{his/her}).
\item \textbf{Les expressions de temps et de lieu} : elles s'éloignent du moment et du lieu de la parole (\eng{tomorrow} $\rightarrow$ \eng{the next day}, \eng{here} $\rightarrow$ \eng{there}, \eng{now} $\rightarrow$ \eng{then}).
\end{enumerate}
\end{aretenir}

\begin{popfigure}
\begin{tikzpicture}[scale=1]
\draw[-{Stealth}, thick, popmuted] (0,0) -- (11,0) node[right, font=\small] {temps};
\foreach \x/\lab in {1.5/{past perfect}, 4/{past simple}, 6.5/{present}, 9/{future (will)}} {
  \fill[popblue] (\x,0) circle (2.2pt);
  \node[below, font=\small, align=center] at (\x,0) {\lab};
}
\draw[-{Stealth}, very thick, poporange] (6.5,0.35) to[bend right=35] node[above, font=\small\itshape] {backshift} (4,0.35);
\draw[-{Stealth}, very thick, poporange] (9,0.35) to[bend right=35] node[above, font=\small\itshape] {backshift} (6.5,0.35);
\draw[-{Stealth}, very thick, poporange] (4,0.35) to[bend right=35] node[above, font=\small\itshape] {backshift} (1.5,0.35);
\end{tikzpicture}
\legende{Le « recul » des temps (backshift) : chaque temps recule d'un cran vers le passé quand le verbe introducteur est au passé.}
\end{popfigure}

\courssub{Le backshift : tableau détaillé des reculs de temps}

\begin{propriete}[Règle du backshift (verbe introducteur au passé)]
Quand le verbe introducteur est au passé (\eng{said}, \eng{told}, \eng{asked}, \eng{replied}…), le temps de la phrase rapportée recule systématiquement d'un cran vers le passé.
\end{propriete}

\begin{center}
\begin{tabularx}{\linewidth}{|X|X|X|}
\hline
\rowcolor{popblueL} Direct Speech (temps) & Indirect Speech (temps) & Exemple \\
\hline
Present Simple (\eng{I work}) & Past Simple (\eng{he worked}) & \eng{“I work hard.”} $\rightarrow$ \eng{He said he worked hard.} \\
\hline
Present Continuous (\eng{I am working}) & Past Continuous (\eng{he was working}) & \eng{“I am working.”} $\rightarrow$ \eng{He said he was working.} \\
\hline
Present Perfect (\eng{I have worked}) & Past Perfect (\eng{he had worked}) & \eng{“I have finished.”} $\rightarrow$ \eng{He said he had finished.} \\
\hline
Past Simple (\eng{I worked}) & Past Perfect (\eng{he had worked}) & \eng{“I bought a car.”} $\rightarrow$ \eng{He said he had bought a car.} \\
\hline
Past Continuous (\eng{I was working}) & Past Perfect Continuous (\eng{he had been working}) & \eng{“I was reading.”} $\rightarrow$ \eng{He said he had been reading.} \\
\hline
Future \eng{will} (\eng{I will go}) & Conditional \eng{would} (\eng{he would go}) & \eng{“I will call you.”} $\rightarrow$ \eng{He said he would call me.} \\
\hline
\eng{can} & \eng{could} & \eng{“I can swim.”} $\rightarrow$ \eng{He said he could swim.} \\
\hline
\eng{may} & \eng{might} & \eng{“It may rain.”} $\rightarrow$ \eng{He said it might rain.} \\
\hline
\eng{must} (obligation) & \eng{had to} & \eng{“I must go.”} $\rightarrow$ \eng{He said he had to go.} \\
\hline
\eng{needn't} & \eng{didn't have to} / \eng{needn't} & \eng{“You needn't wait.”} $\rightarrow$ \eng{He said I didn't have to wait.} \\
\hline
\end{tabularx}
\end{center}

\begin{attention}[Exceptions au backshift]
\begin{itemize}
\item \textbf{Vérité générale / fait scientifique permanent} : le présent reste présent. \eng{The teacher said (that) the earth revolves around the sun.}
\item \textbf{Fait encore vrai au moment du rapport} : \eng{“I live in Douala,” he said.} $\rightarrow$ \eng{He said he lives in Douala.} (s'il y habite encore).
\item \textbf{Verbes modaux \eng{could}, \eng{might}, \eng{should}, \eng{ought to}, \eng{used to}} : ils ne changent pas. \eng{“I could help,” she said.} $\rightarrow$ \eng{She said she could help.}
\item \textbf{Conditionnelles 2\up{ème} et 3\up{ème} forme} : inchangées. \eng{“If I had money, I would travel.”} $\rightarrow$ \eng{He said if he had money he would travel.}
\end{itemize}
\end{attention}

\courssub{Pronoms, possessifs et adjectifs possessifs}

\begin{propriete}[Adaptation du point de vue]
Les pronoms et déterminants possessifs changent pour refléter la personne qui rapporte les paroles par rapport à l'énonciateur original.
\end{propriete}

\begin{center}
\begin{tabular}{|l|l|l|}
\hline
\rowcolor{popblueL} Direct & Indirect (exemples courants) & Note \\
\hline
I & he / she & selon le genre du locuteur \\
\hline
you & I / we / they & selon qui rapporte et à qui \\
\hline
we & they & \\
\hline
my & his / her & \\
\hline
your & my / our / their & \\
\hline
our & their & \\
\hline
this / these & that / those & démonstratifs \\
\hline
\end{tabular}
\end{center}

\courssub{Expressions de temps et de lieu : le décalage spatio-temporel}

\begin{propriete}[Expressions de temps et de lieu]
Les repères spatio-temporels s'éloignent du « ici et maintenant » de l'énonciation originale.
\end{propriete}

\begin{center}
\begin{tabularx}{\linewidth}{|X|X|}
\hline
\rowcolor{popblueL} Direct Speech & Indirect Speech \\
\hline
now & then / at that moment \\
\hline
today & that day \\
\hline
yesterday & the day before / the previous day \\
\hline
tomorrow & the next day / the following day \\
\hline
last week / month / year & the previous week / month / year / the week / month / year before \\
\hline
next week / month / year & the following week / month / year \\
\hline
ago & before / previously \\
\hline
here & there \\
\hline
this / these & that / those \\
\hline
\end{tabularx}
\end{center}

\courssub{Le rôle du verbe introducteur : présent ou passé}

\begin{propriete}[Verbe introducteur au présent / present perfect : pas de backshift]
Si le verbe introducteur est au \textbf{présent} (\eng{says}, \eng{tells}, \eng{asks}) ou au \textbf{present perfect} (\eng{has said}, \eng{has told}), les temps des paroles rapportées \textbf{ne changent pas}. Seuls les pronoms et possessifs s'adaptent si nécessaire.

\eng{She says, “I am tired.”} $\rightarrow$ \eng{She says (that) she is tired.} « Elle dit qu'elle est fatiguée. »

\eng{He has told me, “I live in Douala.”} $\rightarrow$ \eng{He has told me (that) he lives in Douala.} « Il m'a dit qu'il vit à Douala. »
\end{propriete}

\begin{propriete}[Verbe introducteur au passé : backshift obligatoire]
Si le verbe introducteur est au \textbf{passé} (\eng{said}, \eng{told}, \eng{asked}, \eng{announced}, \eng{replied}), les temps des paroles rapportées \textbf{reculent} d'un cran vers le passé : c'est le \cle{backshift}.

\eng{“I am tired,” she said.} $\rightarrow$ \eng{She said (that) she was tired.}

\textbf{Exception :} vérité générale ou fait toujours vrai $\rightarrow$ présent maintenu. \eng{The teacher said that water boils at 100 degrees.}
\end{propriete}

\courssub{Rapporter les questions (indirect questions)}

\begin{propriete}[Transformation des questions]
Au discours indirect, la question devient une proposition subordonnée : \textbf{pas d'inversion sujet-verbe}, \textbf{pas d'auxiliaire \eng{do/does/did}}, \textbf{point final} (pas de point d'interrogation), backshift si verbe introducteur au passé.
\end{propriete}

\textbf{1. Questions fermées (Yes/No questions) $\rightarrow$ \eng{if} ou \eng{whether}.}

\begin{exemplebox}[Questions fermées]
\eng{“Do you like cocoa?” he asked.} $\rightarrow$ \eng{He asked \textbf{if/whether} I \textbf{liked} cocoa.}

\eng{“Are you coming tomorrow?” she asked.} $\rightarrow$ \eng{She asked \textbf{if/whether} I \textbf{was coming} the next day.}

\eng{“Have you finished?” the teacher asked.} $\rightarrow$ \eng{The teacher asked \textbf{if/whether} I \textbf{had finished}.}
\end{exemplebox}

\textbf{2. Questions ouvertes (Wh-questions) $\rightarrow$ mot interrogatif conservé.}

\begin{exemplebox}[Questions ouvertes]
\eng{“Where do you live?” he asked.} $\rightarrow$ \eng{He asked \textbf{where} I \textbf{lived}.}

\eng{“What time does the train leave?” she asked.} $\rightarrow$ \eng{She asked \textbf{what time} the train \textbf{left}.}

\eng{“Why are you late?” the boss asked.} $\rightarrow$ \eng{The boss asked \textbf{why} I \textbf{was} late.}
\end{exemplebox}

\begin{attention}[Pièges fréquents : questions indirectes]
\begin{itemize}
\item \textbf{Ordre des mots :} Sujet + Verbe (affirmatif). \emph{Pas} \eng{He asked where did I live.} $\rightarrow$ \eng{He asked where I lived.}
\item \textbf{Auxiliaires \eng{do/does/did} :} disparaissent. \eng{“Where do you go?”} $\rightarrow$ \eng{He asked where I went.} (pas \eng{did go}).
\item \textbf{Point d'interrogation :} interdit en discours indirect. \eng{He asked where I lived.}
\item \textbf{\eng{whether} vs \eng{if} :} \eng{whether} est plus formel ; \eng{whether ... or not} possible. \eng{He asked whether I liked it or not.}
\end{itemize}
\end{attention}

\courssub{Rapporter les ordres, conseils, demandes et suggestions}

\begin{propriete}[Ordres, demandes, conseils $\rightarrow$ Infinitif]
Structure : \textbf{Verbe introducteur + Objet (COI) + \eng{to} + Base Verbale} (ou \textbf{not to} pour la négation).
Verbes courants : \eng{tell}, \eng{order}, \eng{ask}, \eng{advise}, \eng{warn}, \eng{beg}, \eng{forbid}, \eng{remind}, \eng{encourage}.
\end{propriete}

\begin{exemplebox}[Ordres et demandes]
\eng{The officer said, “Show your ID.”} $\rightarrow$ \eng{The officer \textbf{told} me \textbf{to show} my ID.}

\eng{“Don't touch that wire!” the electrician said.} $\rightarrow$ \eng{The electrician \textbf{warned} me \textbf{not to touch} that wire.}

\eng{“Please, help me,” the old man said.} $\rightarrow$ \eng{The old man \textbf{asked} me \textbf{to help} him.}

\eng{“You should revise your lessons,” my father said.} $\rightarrow$ \eng{My father \textbf{advised} me \textbf{to revise} my lessons.}
\end{exemplebox}

\begin{propriete}[Suggestions $\rightarrow$ \eng{gerund} ou \eng{that} + \eng{should}]
Verbe : \eng{suggest}, \eng{propose}, \eng{recommend}.
\end{propriete}

\begin{exemplebox}[Suggestions]
\eng{“Let's go to the beach,” she said.} $\rightarrow$ \eng{She \textbf{suggested going} to the beach.} (ou \eng{suggested that we should go}).

\eng{“Why don't we eat out?” he said.} $\rightarrow$ \eng{He \textbf{suggested eating} out.}
\end{exemplebox}

\begin{attention}[\eng{Say} vs \eng{Tell} : erreur classique]
\begin{itemize}
\item \eng{Say} \textbf{ne prend pas de COI personne} : \eng{He said (that) he was tired.} / \eng{He said \textbf{to me} that he was tired.}
\item \eng{Tell} \textbf{exige un COI personne} : \eng{He \textbf{told me} (that) he was tired.} (Pas \emph{He told that...} ni \emph{He told to me...}).
\item Pour les ordres : seulement \eng{tell} / \eng{order} + objet + \eng{to-inf}. \emph{Pas} \eng{say to me to go}.
\end{itemize}
\end{attention}

\courssub{Verbes introducteurs et nuances de sens (reporting verbs)}

Au-delà de \eng{say}, \eng{tell}, \eng{ask}, l'anglais utilise une grande variété de verbes qui résument l'intention du locuteur. Chaque verbe impose sa structure.

\begin{center}
\begin{tabularx}{\linewidth}{|X|X|X|}
\hline
\rowcolor{popblueL} Verbe introducteur & Structure & Exemple \\
\hline
\eng{agree}, \eng{refuse}, \eng{promise}, \eng{decide}, \eng{threaten}, \eng{offer} & \eng{to} + infinitif & \eng{“I'll help you.”} $\rightarrow$ \eng{He \textbf{promised to help} me.} \\
\hline
\eng{admit}, \eng{deny}, \eng{regret}, \eng{suggest}, \eng{propose} & \eng{V-ing} (gérondif) & \eng{“I stole the money.”} $\rightarrow$ \eng{He \textbf{admitted stealing} the money.} \\
\hline
\eng{accuse sb of}, \eng{blame sb for}, \eng{apologise for}, \eng{insist on} & Préposition + \eng{V-ing} & \eng{“You broke the vase.”} $\rightarrow$ \eng{He \textbf{accused me of breaking} the vase.} \\
\hline
\eng{claim}, \eng{explain}, \eng{complain}, \eng{state}, \eng{announce} & \eng{that} + clause (backshift) & \eng{“I am innocent.”} $\rightarrow$ \eng{He \textbf{claimed that} he \textbf{was} innocent.} \\
\hline
\eng{advise}, \eng{warn}, \eng{urge}, \eng{encourage}, \eng{persuade} & Objet + \eng{to} + infinitif & \eng{“Don't go.”} $\rightarrow$ \eng{She \textbf{warned me not to go}.} \\
\hline
\eng{remind} & Objet + \eng{to} + infinitif / \eng{that} & \eng{“Don't forget your ticket.”} $\rightarrow$ \eng{He \textbf{reminded me to take} my ticket.} \\
\hline
\end{tabularx}
\end{center}

\begin{methode}[Transformer une phrase au discours indirect : 5 étapes]
\begin{enumerate}
\item \textbf{Identifier} le verbe introducteur (temps ? personne ?) $\rightarrow$ détermine backshift ou non.
\item \textbf{Supprimer} les guillemets, mettre \eng{that} (optionnel pour énoncés) ou \eng{if/whether}/\eng{wh-} (questions) ou \eng{to} (ordres).
\item \textbf{Changer} les pronoms et possessifs selon le nouveau point de vue.
\item \textbf{Appliquer} le backshift aux temps (si verbe au passé) et adapter \eng{now/here/this} $\rightarrow$ \eng{then/there/that}.
\item \textbf{Vérifier} l'ordre des mots (sujet + verbe pour questions) et la ponctuation finale (point).
\end{enumerate}
\end{methode}

\courssub{La ponctuation du discours direct en anglais}

\begin{propriete}[Guillemets, virgule, majuscule]
\begin{enumerate}
\item Les paroles exactes sont placées entre \textbf{guillemets anglais} : “ … ” (en anglais britannique on trouve aussi ‘ … ’ pour les citations imbriquées).
\item Quand le verbe introducteur \textbf{précède} la citation, il est suivi d'une \textbf{virgule} : \eng{She said, “I am tired.”}
\item La citation commence toujours par une \textbf{majuscule}.
\item Quand la citation \textbf{précède} le verbe introducteur, on met une virgule (ou un point d'interrogation/exclamation) \textbf{à l'intérieur} des guillemets : \eng{“I am tired,” she said.} / \eng{“Are you tired?” she asked.}
\item Le point final se place \textbf{à l'intérieur} des guillemets : \eng{She said, “I am tired.”}
\end{enumerate}
\end{propriete}

\begin{attention}[Pièges de ponctuation pour les francophones]
\begin{itemize}
\item En français on écrit : Elle a dit : « Je suis fatiguée. » En anglais, \textbf{pas de deux-points} avant la citation, mais une \textbf{virgule} : \eng{She said, “I am tired.”}
\item Pas d'espace insécable à l'intérieur des guillemets anglais : on écrit \eng{“I am tired.”} et non « I am tired. »
\item Ne mets \textbf{jamais de guillemets} au discours indirect : \eng{She said she was tired.} (et non \emph{She said “she was tired”}).
\end{itemize}
\end{attention}

\courssub{Tableau récapitulatif et comparaison avec le français}

\begin{center}
\begin{tabularx}{\linewidth}{|X|X|X|}
\hline
\rowcolor{popblueL} Français & English & Remarque \\
\hline
Discours direct & Direct speech & Paroles exactes entre guillemets \\
\hline
Discours indirect & Indirect / reported speech & Sens rapporté, sans guillemets \\
\hline
Il a dit : « Je viens. » & He said, “I am coming.” & Virgule en anglais, deux-points en français \\
\hline
Il a dit qu'il venait. & He said (that) he was coming. & \eng{that} = « que », souvent omis en anglais \\
\hline
Elle dit qu'elle est fatiguée. & She says (that) she is tired. & Verbe au présent : pas de backshift \\
\hline
Elle a dit qu'elle était fatiguée. & She said (that) she was tired. & Verbe au passé : backshift \\
\hline
« Viens ! » & “Come!” (direct) & Impératif direct \\
\hline
Il m'a dit de venir. & He told me to come. & Infinitif au discours indirect \\
\hline
« Est-ce qu'il vient ? » & “Is he coming?” (direct) & Question directe \\
\hline
Il a demandé s'il venait. & He asked if he was coming. & \eng{if/whether} + ordre affirmatif \\
\hline
\end{tabularx}
\end{center}

\begin{attention}[Erreur fréquente : oublier le backshift]
Le français et l'anglais reculent tous les deux les temps (« Je suis fatiguée » $\rightarrow$ « Elle a dit qu'elle \emph{était} fatiguée »), donc le mécanisme est familier. Mais les francophones oublient souvent le backshift en anglais parce que la concordance des temps y est \textbf{plus stricte}. On ne peut pas dire \emph{She said she is tired} si le verbe introducteur est \eng{said} (sauf vérité encore valable au moment où l'on parle). De même, \eng{will} devient \eng{would} : \eng{He said he would come.} (et non \emph{He said he will come}).
\end{attention}

\begin{attention}[Autres erreurs fréquentes des francophones]
\begin{itemize}
\item \textbf{\eng{Say} + COI :} \emph{He said me…} $\rightarrow$ \textbf{He told me…} / \textbf{He said to me…}
\item \textbf{Ordre :} \emph{He said me to go.} $\rightarrow$ \textbf{He told me to go.} / \textbf{He ordered me to go.}
\item \textbf{Question indirecte :} \emph{He asked where did I go.} $\rightarrow$ \textbf{He asked where I went.}
\item \textbf{\eng{Suggest} + infinitif :} \emph{He suggested to go.} $\rightarrow$ \textbf{He suggested going.} / \textbf{He suggested that we (should) go.}
\item \textbf{\eng{Must} au passé :} \emph{He said he must go.} $\rightarrow$ \textbf{He said he had to go.} (sauf déduction logique : \eng{He said it must be true} $\rightarrow$ \eng{He said it must be true}).
\end{itemize}
\end{attention}

\courssub{Exercice résolu}

\begin{exoresolu}[Direct ou indirect ? Transformation]
\textbf{Énoncé.} Pour chaque phrase, indique s'il s'agit de direct speech (D) ou d'indirect speech (I), puis transforme les phrases au discours direct en discours indirect.

1. \eng{The farmer said, “I grow cocoa in the South West Region.”}

2. \eng{My sister told me that she had passed her GCE A Level.}

3. \eng{“We will visit our grandmother tomorrow,” the children said.}

4. \eng{The journalist announced that the match would start at four o'clock.}
\tcblower
\textbf{Solution détaillée.}

1. \textbf{D} (guillemets + paroles exactes). Transformation : verbe introducteur au passé (\eng{said}) donc backshift : \eng{grow} (présent) $\rightarrow$ \eng{grew} (prétérit) ; pronom \eng{I} $\rightarrow$ \eng{he}. Résultat : \eng{The farmer said (that) he grew cocoa in the South West Region.} « Le planteur a dit qu'il cultivait le cacao dans la région du Sud-Ouest. »

2. \textbf{I} (pas de guillemets, \eng{that}, backshift \eng{had passed}). Pas de transformation demandée.

3. \textbf{D}. Transformation : \eng{will visit} $\rightarrow$ \eng{would visit} ; \eng{our} $\rightarrow$ \eng{their} ; \eng{tomorrow} $\rightarrow$ \eng{the next day}. Résultat : \eng{The children said (that) they would visit their grandmother the next day.} « Les enfants ont dit qu'ils rendraient visite à leur grand-mère le lendemain. »

4. \textbf{I} (\eng{that} + backshift \eng{would start}). Pas de transformation demandée.
\end{exoresolu}

\courssub{Entraînement guidé : énoncés, questions, ordres}

\begin{exercice}[À toi de jouer — Énoncés]
Transforme les phrases suivantes au discours indirect (verbe introducteur au passé) :

1. \eng{The doctor said, “You must rest for three days.”}

2. \eng{“I am preparing my lessons now,” the student said.}

3. \eng{Our neighbours said, “We bought a new car yesterday.”}

4. \eng{The bus driver told the passengers, “We will arrive in Bafoussam this evening.”}
\end{exercice}

\begin{exercice}[À toi de jouer — Questions]
Transforme les questions directes en questions indirectes (verbe introducteur au passé) :

5. \eng{The customs officer asked, “Do you have anything to declare?”}

6. \eng{“Where did you put my passport?” my brother asked me.}

7. \eng{“Why are you crying?” the teacher asked the pupil.}

8. \eng{“Will you be at the meeting tomorrow?” she asked him.}
\end{exercice}

\begin{exercice}[À toi de jouer — Ordres, conseils, suggestions]
Transforme au discours indirect :

9. \eng{The coach said, “Run faster!”}

10. \eng{“Don't forget to lock the door,” my mother said.}

11. \eng{“You should see a dentist,” my friend said.}

12. \eng{“Let's organise a party for the end of term,” the class representative said.}
\end{exercice}

\begin{corrige}[Exercices « À toi de jouer »]
\textbf{Énoncés :}
1. \eng{The doctor said (that) I had to rest for three days.} (\eng{must} $\rightarrow$ \eng{had to} ; \eng{you} $\rightarrow$ \eng{I}.)

2. \eng{The student said (that) he/she was preparing his/her lessons then.} (présent continu $\rightarrow$ prétérit continu ; \eng{now} $\rightarrow$ \eng{then}.)

3. \eng{Our neighbours said (that) they had bought a new car the day before.} (prétérit simple $\rightarrow$ past perfect ; \eng{yesterday} $\rightarrow$ \eng{the day before}.)

4. \eng{The bus driver told the passengers (that) they would arrive in Bafoussam that evening.} (\eng{will} $\rightarrow$ \eng{would} ; \eng{this evening} $\rightarrow$ \eng{that evening}.)

\textbf{Questions :}
5. \eng{The customs officer asked if/whether I had anything to declare.} (\eng{do have} $\rightarrow$ \eng{had} ; ordre affirmatif.)

6. \eng{My brother asked me where I had put his passport.} (\eng{did put} $\rightarrow$ \eng{had put} ; pas d'auxiliaire \eng{did}.)

7. \eng{The teacher asked the pupil why he/she was crying.} (présent continu $\rightarrow$ prétérit continu ; ordre affirmatif.)

8. \eng{She asked him if/whether he would be at the meeting the next day.} (\eng{will be} $\rightarrow$ \eng{would be} ; \eng{tomorrow} $\rightarrow$ \eng{the next day}.)

\textbf{Ordres, conseils, suggestions :}
9. \eng{The coach told me to run faster.} (\eng{run} $\rightarrow$ \eng{to run} ; ordre $\rightarrow$ \eng{tell} + objet + \eng{to-inf}.)

10. \eng{My mother reminded me to lock the door.} (négatif \eng{don't forget} $\rightarrow$ \eng{remind} + \eng{to-inf} positif) ou \eng{My mother told me not to forget to lock the door.}

11. \eng{My friend advised me to see a dentist.} (\eng{should see} $\rightarrow$ \eng{advise} + \eng{to-inf}.)

12. \eng{The class representative suggested organising a party for the end of term.} (\eng{Let's} $\rightarrow$ \eng{suggest} + \eng{V-ing}.) Ou : \eng{... suggested that they (should) organise a party...}
\end{corrige}

\begin{experience}[Jeu de rôle : Interview « People »]
\textbf{Consigne :} En binôme. L'élève A est une star (footballeur, chanteur, ministre…). L'élève B est journaliste.
\begin{enumerate}
\item B interviewe A (5 questions variées : passé, futur, opinion, conseil, projet).
\item A répond naturellement.
\item B rédige un court article (5-6 lignes) pour \eng{Cameroon Tribune} en ne utilisant \textbf{que du discours indirect} (verbes variés : \eng{stated, revealed, admitted, promised, warned, suggested}…).
\item Échange des rôles.
\end{enumerate}
\textbf{Exemple de début d'article :} \eng{The Minister of Youth revealed that the government had invested heavily in sports infrastructure. He promised that new stadiums would open next year. He warned young people not to neglect their studies. He suggested creating more local leagues…}
\end{experience}

\begin{savaistu}[Le discours indirect est partout dans la presse anglophone]
Ouvre un journal anglophone — la BBC, \eng{Cameroon Tribune} ou \eng{The Guardian} — et tu trouveras du discours indirect dans presque chaque article : \eng{“The Minister announced that schools would reopen in September.”} « Le ministre a annoncé que les écoles rouvriraient en septembre. » Les journalistes citent rarement les paroles mot à mot ; ils les rapportent. C'est pourquoi les titres et les dépêches regorgent de \eng{said that}, \eng{announced that}, \eng{told reporters that}… Maîtriser le reported speech, c'est donc aussi mieux \emph{comprendre} la presse anglophone — une compétence précieuse pour l'épreuve de reading comprehension du Bac !
\end{savaistu}

\coursec{Le backshift : recul des temps}

Dans la section précédente, nous avons observé la différence fondamentale entre le discours direct (\eng{direct speech}), qui cite les paroles exactes entre guillemets, et le discours indirect (\eng{reported speech}), qui rapporte ces paroles de loin. Nous avons remarqué que quelque chose changeait dans les verbes : \eng{“I am tired,” she said} devient \eng{She said she was tired}. Ce phénomène a un nom : le \cle{backshift}, c'est-à-dire le « recul des temps ». C'est le cœur de cette leçon.

\courssub{Observation : un texte pour voir le backshift en action}

Lisons d'abord ce court texte original, puis observons ce qui arrive aux verbes quand les paroles sont rapportées.

\begin{textebox}[A health campaign in Yaoundé — original text]
Last week, a team of doctors visited a secondary school in Yaoundé to talk about well-being. Before the visit, the head teacher had announced: “The doctors will arrive on Monday.” On Monday morning, Dr. Ewusi told the students: “We are testing your eyes today.” One student, Amina, said: “I have never worn glasses.” Her friend Boris explained: “I lost my glasses last month, and I cannot read the board.” The nurse added: “You must drink more water, and you should sleep eight hours every night.” At the end of the day, the principal reported everything to the parents. She said that the doctors had arrived on time, that they had tested the students' eyes, and that Amina had never worn glasses. She also said that Boris had lost his glasses and could not read the board. Finally, she told the parents that their children had to drink more water.
\end{textebox}

\textbf{Glossaire}

\begin{center}
\begin{tabularx}{\linewidth}{|l|l|X|}
\hline
\rowcolor{popblueL} \textbf{Mot anglais} & \textbf{Nature} & \textbf{Traduction française} \\
\hline
well-being & nom & le bien-être \\
\hline
head teacher & nom & le proviseur / la proviseure \\
\hline
to test & verbe & examiner, dépister \\
\hline
glasses & nom pluriel & les lunettes \\
\hline
the board & nom & le tableau \\
\hline
nurse & nom & l'infirmier / l'infirmière \\
\hline
to report & verbe & rapporter, faire un compte rendu \\
\hline
\end{tabularx}
\end{center}

\textbf{Questions de compréhension}
\begin{enumerate}
\item What did the head teacher announce before the visit?
\item What did Dr. Ewusi tell the students on Monday morning?
\item How does the principal report Boris's words at the end of the text?
\end{enumerate}

\textbf{Observation grammaticale.} Comparons les paroles directes et les paroles rapportées (le backshift mécanique) :

\begin{center}
\begin{tabularx}{\linewidth}{|X|X|}
\hline
\rowcolor{popblueL} \textbf{Direct speech} & \textbf{Reported speech (backshift)} \\
\hline
\eng{“The doctors will arrive on Monday.”} & \eng{She said the doctors would arrive on Monday.} \\
\hline
\eng{“I have never worn glasses.”} & \eng{She said that Amina had never worn glasses.} \\
\hline
\eng{“I cannot read the board.”} & \eng{She said that Boris could not read the board.} \\
\hline
\eng{“You must drink more water.”} & \eng{She told the parents that their children had to drink more water.} \\
\hline
\end{tabularx}
\end{center}

On constate que chaque temps verbal a « reculé » d'un cran vers le passé : \eng{will} devient \eng{would}, \eng{have worn} devient \eng{had worn}, \eng{cannot} devient \eng{could not}, \eng{must} devient \eng{had to}. C'est exactement le mécanisme du backshift.

\courssub{La règle générale du backshift}

\begin{propriete}[Règle du backshift]
Quand le \cle{verbe introducteur} est au \textbf{passé} (\eng{said}, \eng{told}, \eng{asked}, \eng{explained}, \eng{announced}...), les temps du discours direct \textbf{reculent d'un cran vers le passé} dans le discours rapporté :
\begin{itemize}
\item présent $\rightarrow$ passé ;
\item passé ou present perfect $\rightarrow$ past perfect (pluperfect) ;
\item futur (\eng{will}) $\rightarrow$ conditionnel (\eng{would}).
\end{itemize}
Ce recul est logique : les paroles ont été prononcées \emph{avant} le moment où on les rapporte, donc on les « enfonce » davantage dans le passé.
\end{propriete}

\begin{aretenir}[L'idée-clé]
Le backshift n'est pas une complication artificielle : il exprime la \textbf{distance temporelle} entre le moment de la parole et le moment du rapport. Le français fait exactement la même chose : « Je \emph{suis} fatigué », a-t-elle dit $\rightarrow$ Elle a dit qu'elle \emph{était} fatiguée.
\end{aretenir}

\begin{popfigure}
\begin{center}
\begin{tikzpicture}[
  stepS/.style={draw=popblue, fill=popblueL, rounded corners=3pt, minimum width=4.6cm, minimum height=0.9cm, align=center, font=\small},
  lbl/.style={font=\small\itshape, text=popmuted}
]
\node[stepS, align=center] (s1) at (0,0){\textbf{Present simple}\\ \eng{“I work hard.”}};
\node[stepS, fill=poporangeL, draw=poporange, align=center] (s2) at (5.4,-1.4){\textbf{Past simple}\\ \eng{He said he worked hard.}};
\node[stepS, fill=poppinkL, draw=poppink, align=center] (s3) at (10.8,-2.8){\textbf{Past perfect}\\ \eng{He said he had worked hard.}};
\draw[-{Stealth[length=3mm]}, very thick, popblue] (s1.south east) -- (s2.north west) node[midway, above, sloped, lbl] {recul d'un cran};
\draw[-{Stealth[length=3mm]}, very thick, poporange] (s2.south east) -- (s3.north west) node[midway, above, sloped, lbl] {recul d'un cran};
\node[lbl, align=center] at (5.4,1.1) {Chaque marche = un cran vers le passé};
\end{tikzpicture}
\end{center}
\legende{L'escalier des temps : present simple $\rightarrow$ past simple $\rightarrow$ past perfect.}
\end{popfigure}

\courssub{Les transformations, temps par temps}

\textbf{1. Present simple $\rightarrow$ past simple}

\begin{exemplebox}[Present simple rapporté]
\eng{“I work hard,” he said.} $\rightarrow$ \eng{He said (that) he worked hard.}\\
« Je travaille dur », a-t-il dit. $\rightarrow$ Il a dit qu'il travaillait dur.

\eng{“We live in Douala,” they said.} $\rightarrow$ \eng{They said they lived in Douala.}\\
Ils ont dit qu'ils vivaient à Douala.
\end{exemplebox}

\textbf{2. Present continuous $\rightarrow$ past continuous}

\begin{exemplebox}[Present continuous rapporté]
\eng{“I am reading a book about health,” she said.} $\rightarrow$ \eng{She said she was reading a book about health.}\\
« Je lis un livre sur la santé », a-t-elle dit. $\rightarrow$ Elle a dit qu'elle lisait un livre sur la santé.

\eng{“The children are playing outside,” the teacher said.} $\rightarrow$ \eng{The teacher said the children were playing outside.}
\end{exemplebox}

\textbf{3. Past simple et present perfect $\rightarrow$ past perfect}

C'est un point crucial : \textbf{deux temps différents du discours direct} (past simple et present perfect) \textbf{convergent vers un seul temps} dans le discours rapporté : le past perfect (\eng{had + participe passé}).

\begin{exemplebox}[Past simple et present perfect rapportés]
\eng{“I lost my key,” he said.} $\rightarrow$ \eng{He said he had lost his key.} (past simple $\rightarrow$ past perfect)

\eng{“I have lost my key,” he said.} $\rightarrow$ \eng{He said he had lost his key.} (present perfect $\rightarrow$ past perfect)

\eng{“I have finished my homework,” Peter said.} $\rightarrow$ \eng{Peter said he had finished his homework.}\\
« J'ai fini mes devoirs », a dit Peter. $\rightarrow$ Peter a dit qu'il avait fini ses devoirs.
\end{exemplebox}

\begin{attention}[Deux temps, une seule forme rapportée]
En français, la distinction se retrouve parfois (« il a dit qu'il \emph{avait perdu} sa clé »), mais en anglais, impossible de distinguer à l'arrivée si la phrase originale était au past simple ou au present perfect : les deux donnent \eng{had lost}. Ne cherchez donc jamais à « traduire » le present perfect par autre chose que le past perfect au discours rapporté.
\end{attention}

\textbf{4. Will $\rightarrow$ would ; can $\rightarrow$ could ; may $\rightarrow$ might}

Les auxiliaires modaux reculent eux aussi :

\begin{exemplebox}[Modaux rapportés]
\eng{“I will help you,” she said.} $\rightarrow$ \eng{She said she would help me.}\\
« Je t'aiderai », a-t-elle dit. $\rightarrow$ Elle a dit qu'elle m'aiderait.

\eng{“I can swim very well,” the boy said.} $\rightarrow$ \eng{The boy said he could swim very well.}\\
Le garçon a dit qu'il savait nager très bien.

\eng{“It may rain tomorrow,” the farmer said.} $\rightarrow$ \eng{The farmer said it might rain the next day.}\\
Le fermier a dit qu'il pleuvrait peut-être le lendemain.
\end{exemplebox}

\textbf{5. Must $\rightarrow$ had to ; shall $\rightarrow$ should / would}

\eng{Must} n'a pas de passé propre : on le remplace par \eng{had to} (obligation passée). \eng{Shall} devient \eng{should} (conseil, question) ou \eng{would} (futur).

\begin{exemplebox}[Must et shall rapportés]
\eng{“You must wash your hands before eating,” the nurse said.} $\rightarrow$ \eng{The nurse said we had to wash our hands before eating.}\\
L'infirmière a dit que nous devions nous laver les mains avant de manger.

\eng{“Shall we meet at the hospital?” he asked.} $\rightarrow$ \eng{He asked if they should meet at the hospital.}
\end{exemplebox}

\begin{popfigure}
\begin{center}
\begin{tikzpicture}[
  box/.style={draw=popblue, fill=popblueL, rounded corners=3pt, align=center, font=\small, minimum height=0.85cm},
  rbox/.style={draw=popgreen, fill=popgreenL, rounded corners=3pt, align=center, font=\small, minimum height=0.85cm},
  arr/.style={-{Stealth[length=2.5mm]}, very thick, poporange}
]
\node[box] (d1) at (0,0) {\eng{work}};
\node[rbox] (r1) at (6.2,0) {\eng{worked}};
\draw[arr] (d1) -- (r1);
\node[box] (d2) at (0,-1.2) {\eng{am reading}};
\node[rbox] (r2) at (6.2,-1.2) {\eng{was reading}};
\draw[arr] (d2) -- (r2);
\node[box] (d3) at (0,-2.4) {\eng{lost / have lost}};
\node[rbox] (r3) at (6.2,-2.4) {\eng{had lost}};
\draw[arr] (d3) -- (r3);
\node[box] (d4) at (0,-3.6) {\eng{will / can / may}};
\node[rbox] (r4) at (6.2,-3.6) {\eng{would / could / might}};
\draw[arr] (d4) -- (r4);
\node[box] (d5) at (0,-4.8) {\eng{must}};
\node[rbox] (r5) at (6.2,-4.8) {\eng{had to}};
\draw[arr] (d5) -- (r5);
\node[font=\small\bfseries, text=popblue] at (0,0.8) {Direct speech};
\node[font=\small\bfseries, text=popgreen] at (6.2,0.8) {Reported speech};
\end{tikzpicture}
\end{center}
\legende{Tableau de transformation : chaque flèche représente le recul d'un cran vers le passé.}
\end{popfigure}

\courssub{Tableau récapitulatif complet des transformations}

\begin{center}
\begin{tabularx}{\linewidth}{|l|l|X|}
\hline
\rowcolor{popblueL} \textbf{Direct speech} & \textbf{Reported speech} & \textbf{Exemple} \\
\hline
present simple & past simple & \eng{“I eat well.” $\rightarrow$ He said he ate well.} \\
\hline
present continuous & past continuous & \eng{“I am resting.” $\rightarrow$ She said she was resting.} \\
\hline
past simple & past perfect & \eng{“I fell ill.” $\rightarrow$ He said he had fallen ill.} \\
\hline
present perfect & past perfect & \eng{“I have recovered.” $\rightarrow$ She said she had recovered.} \\
\hline
past perfect & past perfect (inchangé) & \eng{“I had never smoked.” $\rightarrow$ He said he had never smoked.} \\
\hline
will & would & \eng{“I will exercise.” $\rightarrow$ She said she would exercise.} \\
\hline
can & could & \eng{“I can run fast.” $\rightarrow$ He said he could run fast.} \\
\hline
may & might & \eng{“It may help.” $\rightarrow$ She said it might help.} \\
\hline
must & had to & \eng{“You must rest.” $\rightarrow$ The doctor said I had to rest.} \\
\hline
shall & should / would & \eng{“Shall I come?” $\rightarrow$ He asked if he should come.} \\
\hline
\end{tabularx}
\end{center}

\begin{aretenir}[Le past perfect ne recule plus]
Le past perfect est déjà le temps « le plus reculé » de l'anglais : il ne peut pas reculer davantage. \eng{“I had already left,” she said} $\rightarrow$ \eng{She said she had already left.} Aucun changement.
\end{aretenir}

\courssub{Quand le backshift est facultatif ou absent}

La règle du recul connaît trois exceptions importantes.

\textbf{1. Les vérités générales et les lois scientifiques.} Ce qui est vrai en tout temps ne recule pas.

\begin{exemplebox}[Vérité générale : pas de backshift]
\eng{“Water boils at 100°C,” the teacher said.} $\rightarrow$ \eng{The teacher said water boils at 100°C.}\\
« L'eau bout à 100 degrés », a dit le professeur. $\rightarrow$ Le professeur a dit que l'eau bout à 100 degrés. (Le backshift \eng{boiled} serait possible mais inutile : l'eau bout toujours à 100 degrés.)
\end{exemplebox}

\textbf{2. Les faits encore vrais au moment du rapport.} Si la situation décrite est toujours actuelle, on peut garder le temps d'origine (ou appliquer le backshift : les deux sont corrects).

\begin{exemplebox}[Fait encore vrai : backshift facultatif]
\eng{“I live in Bamenda,” she said.} $\rightarrow$ \eng{She said she lives in Bamenda.} (elle y habite toujours) ou \eng{She said she lived in Bamenda.}

\eng{“The hospital is near the market,” he said.} $\rightarrow$ \eng{He said the hospital is near the market.} (l'hôpital y est toujours)
\end{exemplebox}

\textbf{3. Le verbe introducteur au présent.} Si l'on rapporte avec \eng{says}, \eng{tells} (présent), il n'y a \textbf{aucun backshift}.

\begin{exemplebox}[Verbe introducteur au présent : pas de backshift]
\eng{“I am tired,” she says.} $\rightarrow$ \eng{She says she is tired.}\\
Elle dit qu'elle est fatiguée. (Rien ne change !)
\end{exemplebox}

\courssub{Comparaison avec le français}

\begin{center}
\begin{tabularx}{\linewidth}{|X|X|X|}
\hline
\rowcolor{popblueL} \textbf{Français} & \textbf{English} & \textbf{Remarque} \\
\hline
Il a dit qu'il \emph{était} malade. & \eng{He said he was ill.} & imparfait = past simple : parallèle parfait \\
\hline
Elle a dit qu'elle \emph{avait fini}. & \eng{She said she had finished.} & plus-que-parfait = past perfect : parallèle parfait \\
\hline
Il a dit qu'il \emph{viendrait}. & \eng{He said he would come.} & conditionnel français = \eng{would} : parallèle parfait \\
\hline
Elle dit qu'elle \emph{est} fatiguée. & \eng{She says she is tired.} & introducteur au présent : pas de recul dans les deux langues \\
\hline
\end{tabularx}
\end{center}

Bonne nouvelle : le français et l'anglais fonctionnent presque à l'identique. Le danger n'est donc pas la logique, mais l'\textbf{oubli} du recul quand on parle vite.

\begin{attention}[Erreur fréquente n°1 : oublier le backshift]
Quand la situation rapportée est passée, le temps DOIT reculer :\\
✗ \eng{He said he is tired.} (situation passée)\\
✓ \eng{He said he was tired.}\\
Il a dit qu'il était fatigué. L'erreur vient de ce qu'on traduit mot à mot depuis le français parlé relâché ; en anglais correct, le backshift est obligatoire ici.
\end{attention}

\begin{attention}[Erreur fréquente n°2 : confondre \eng{would} rapporté et conditionnel]
\eng{Would} dans le discours rapporté n'est PAS un conditionnel hypothétique : c'est le futur \eng{will} vu du passé.\\
\eng{She said she would come.} = Elle a dit qu'elle viendrait. (futur rapporté, la promesse était réelle)\\
Ne traduisez pas : « elle viendrait (si...) ». Il n'y a aucune condition ici.
\end{attention}

\courssub{Méthode et entraînement guidé}

\begin{methode}[Appliquer le backshift en quatre étapes]
\begin{enumerate}
\item \textbf{Repérer le temps du verbe introducteur.} S'il est au présent (\eng{says}), ne rien changer. S'il est au passé (\eng{said}, \eng{told}), continuer.
\item \textbf{Identifier le temps de la citation.} Present simple ? Present perfect ? \eng{will} ? \eng{must} ?
\item \textbf{Reculer d'un cran} selon le tableau récapitulatif : present $\rightarrow$ past ; past / present perfect $\rightarrow$ past perfect ; \eng{will} $\rightarrow$ \eng{would} ; \eng{can} $\rightarrow$ \eng{could} ; \eng{must} $\rightarrow$ \eng{had to}.
\item \textbf{Vérifier les exceptions.} Vérité générale ou fait encore vrai ? Le backshift est alors facultatif. Past perfect déjà présent ? Il reste inchangé.
\end{enumerate}
\end{methode}

\begin{exoresolu}[Transformation guidée]
\textbf{Énoncé.} Mettez au discours rapporté, en commençant par \eng{The doctor told me...} :\\
\eng{“You smoke too much, you have gained weight, and you will feel better if you exercise. You must sleep more.”}
\tcblower
\textbf{Solution détaillée.}
\begin{enumerate}
\item Verbe introducteur : \eng{told}, au passé $\rightarrow$ backshift obligatoire.
\item \eng{smoke} (present simple) $\rightarrow$ \eng{smoked} ; \eng{have gained} (present perfect) $\rightarrow$ \eng{had gained} ; \eng{will feel} $\rightarrow$ \eng{would feel} ; \eng{must sleep} $\rightarrow$ \eng{had to sleep}.
\item Phrase complète : \eng{The doctor told me (that) I smoked too much, (that) I had gained weight, and (that) I would feel better if I exercised. He added that I had to sleep more.}
\end{enumerate}
Traduction : Le médecin m'a dit que je fumais trop, que j'avais pris du poids et que je me sentirais mieux si je faisais de l'exercice. Il a ajouté que je devais dormir davantage.
\end{exoresolu}

\begin{exercice}[À vous de jouer]
Mettez les phrases suivantes au discours rapporté avec le verbe introducteur indiqué :
\begin{enumerate}
\item \eng{“I am training for the football match,” Eric said.}
\item \eng{“We have cleaned the water tank,” the students said.}
\item \eng{“I will visit my grandmother in Buea,” she said.}
\item \eng{“You must eat more fruit,” the nurse told the children.}
\item \eng{“Malaria is dangerous,” the doctor said.} (attention : vérité générale)
\end{enumerate}
\end{exercice}

\begin{corrige}[Exercice « À vous de jouer »]
\begin{enumerate}
\item \eng{Eric said he was training for the football match.} (present continuous $\rightarrow$ past continuous)
\item \eng{The students said they had cleaned the water tank.} (present perfect $\rightarrow$ past perfect)
\item \eng{She said she would visit her grandmother in Buea.} (\eng{will} $\rightarrow$ \eng{would})
\item \eng{The nurse told the children they had to eat more fruit.} (\eng{must} $\rightarrow$ \eng{had to})
\item \eng{The doctor said malaria is dangerous.} (vérité générale : pas de backshift ; \eng{was} serait toléré mais moins naturel)
\end{enumerate}
\end{corrige}

\begin{aretenir}[L'essentiel de la section]
\begin{itemize}
\item Verbe introducteur au passé $\rightarrow$ les temps reculent d'un cran : present $\rightarrow$ past, past / present perfect $\rightarrow$ past perfect, \eng{will} $\rightarrow$ \eng{would}, \eng{can} $\rightarrow$ \eng{could}, \eng{may} $\rightarrow$ \eng{might}, \eng{must} $\rightarrow$ \eng{had to}.
\item Le past perfect ne recule pas (il est déjà au maximum).
\item Pas de backshift si le verbe introducteur est au présent, et backshift facultatif pour les vérités générales et les faits encore vrais.
\item Le français suit la même logique : fiez-vous à votre intuition de francophone... mais n'oubliez pas d'appliquer le recul !
\end{itemize}
\end{aretenir}

\coursec{Pronoms et expressions de temps et de lieu}

Dans la section précédente, nous avons vu que lorsqu'on rapporte des paroles, les \cle{temps verbaux} reculent d'un cran vers le passé : c'est le \emph{backshift}. Mais ce n'est pas tout ce qui change. Imagine la scène suivante : ton ami Peter te dit \eng{“I will bring my books here tomorrow.”} Le lendemain, tu racontes cette conversation à un camarade. Tu ne peux plus dire \eng{I}, \eng{my}, \eng{here} ou \eng{tomorrow} : ces mots avaient un sens \emph{au moment où Peter parlait}, mais ils n'ont plus le même sens \emph{au moment où toi tu rapportes}. Cette section t'apprend à ajuster les \cle{pronoms}, les \cle{possessifs} et les \cle{expressions de temps et de lieu}.

\courssub{Observation : pourquoi les mots changent-ils ?}

\begin{textebox}[At the school canteen — a dialogue]
\textbf{Monday, 12.30 p.m.} — Sarah and Peter are chatting at the school canteen in Yaoundé.

\textbf{Sarah:} “I forgot my money at home today. Can you lend me some?”

\textbf{Peter:} “Of course! I will pay for your lunch now, and you can give me my money back tomorrow.”

\textbf{Sarah:} “Thank you so much. We bought these sandwiches here yesterday and they were delicious.”

\medskip
\textbf{Tuesday} — Sarah's friend Ann asks her about the conversation. Sarah reports:

\emph{“Peter said he would pay for my lunch then, and that I could give him his money back the next day. I told him we had bought those sandwiches there the day before.”}
\end{textebox}

\begin{center}
\begin{tabular}{lll}
\toprule
\textbf{Mot} & \textbf{Nature} & \textbf{Traduction} \\
\midrule
to lend (lent, lent) & verbe irrégulier & prêter \\
to pay back & verbe à particule & rembourser \\
delicious & adjectif & délicieux \\
to report & verbe & rapporter (des paroles) \\
\bottomrule
\end{tabular}
\end{center}

\textbf{Questions d'observation :}
\begin{enumerate}
\item Dans le discours direct, Peter dit \eng{“I will pay”}. Dans le discours rapporté, que devient \eng{I} ? Pourquoi ?
\item Que devient \eng{tomorrow} quand Sarah rapporte la conversation le mardi ?
\item Que deviennent \eng{these sandwiches} et \eng{here} ?
\end{enumerate}

Tu observes que \eng{I} devient \eng{he} (parce que Sarah parle \emph{de} Peter), que \eng{tomorrow} devient \eng{the next day}, et que \eng{these} et \eng{here} deviennent \eng{those} et \eng{there}. Rien de mécanique : chaque changement suit une \textbf{logique de point de vue}.

\courssub{Le triangle de la communication : qui parle, qui rapporte ?}

La clé de toute cette section tient en une phrase : \textbf{le pronom dépend de QUI parle et de QUI rapporte}. Il y a trois acteurs dans toute situation de discours rapporté :

\begin{itemize}
\item le \cle{speaker} (le locuteur d'origine) : celui qui a prononcé les paroles ;
\item le \cle{reporter} (le rapporteur) : celui qui raconte ensuite ces paroles ;
\item le \cle{listener} (le destinataire) : celui qui écoute le rapport.
\end{itemize}

\begin{popfigure}
\begin{tikzpicture}[
  actor/.style={rectangle, rounded corners=6pt, draw=popblue, fill=popblueL, thick, align=center, minimum width=3.2cm, minimum height=1cm, font=\small\bfseries},
  lbl/.style={font=\footnotesize\itshape, text=popdark}
]
\node[actor, align=center] (S) at (0,3){Speaker\\ “\textbf{I} love \textbf{my} school”};
\node[actor, fill=poporangeL, draw=poporange, align=center] (R) at (5,3){Reporter\\ “she said...”};
\node[actor, fill=popgreenL, draw=popgreen, align=center] (L) at (2.5,0){Listener\\ hears the report};
\draw[-{Stealth}, thick, popblue] (S) -- node[lbl, above, sloped]{original words} (R);
\draw[-{Stealth}, thick, poporange] (R) -- node[lbl, above, sloped]{reported words} (L);
\draw[-{Stealth}, thick, popgreen, dashed] (L) -- node[lbl, left]{understands} (S);
\node[font=\footnotesize, text=popmuted, align=center] at (2.5,-1.1) {“I” of the speaker becomes “he/she” in the mouth of the reporter.};
\end{tikzpicture}
\legende{Le triangle de la communication : les pronoms pivotent selon le point de vue du rapporteur.}
\end{popfigure}

\begin{propriete}[Règle des pronoms : adapte-les à la situation de celui qui rapporte]
Quand on rapporte des paroles, les pronoms personnels et possessifs changent pour refléter le point de vue du \textbf{rapporteur}, et non plus celui du locuteur d'origine :
\begin{itemize}
\item \eng{I} (le locuteur) $\rightarrow$ \eng{he / she} : \eng{“I love my school,” said Mary.} $\rightarrow$ \eng{Mary said (that) she loved her school.} « Mary a dit qu'elle aimait son école. »
\item \eng{you} $\rightarrow$ \eng{I / me / he / she / we / us}, selon à qui s'adressait le locuteur ;
\item \eng{we} $\rightarrow$ \eng{they} (si le rapporteur n'est pas inclus dans le groupe) ;
\item \eng{my} $\rightarrow$ \eng{his / her} ; \eng{our} $\rightarrow$ \eng{their}.
\end{itemize}
\textbf{Exception logique :} si le rapporteur est lui-même concerné, le pronom ne change pas : \eng{“I will help you,” Peter told me.} $\rightarrow$ \eng{Peter told me he would help me.} (c'est bien \emph{moi} qui suis aidé : \eng{you} devient \eng{me}).
\end{propriete}

\begin{center}
\begin{tabular}{|l|l|l|}
\hline
\rowcolor{popblueL} \textbf{Discours direct} & \textbf{Discours rapporté} & \textbf{Exemple} \\
\hline
I / me & he / she / him / her & “I am tired,” said John. $\rightarrow$ He said \textbf{he} was tired. \\
\hline
my / mine & his / her & “This is my pen,” said Ann. $\rightarrow$ She said it was \textbf{her} pen. \\
\hline
we / us & they / them & “We are ready,” they said. $\rightarrow$ They said \textbf{they} were ready. \\
\hline
our / ours & their / theirs & “Our team won,” they said. $\rightarrow$ They said \textbf{their} team had won. \\
\hline
you & I / me / he / she / we / us & “I will visit you,” he told me. $\rightarrow$ He told me he would visit \textbf{me}. \\
\hline
your & my / his / her / our / their & “I like your bag,” she told him. $\rightarrow$ She told him she liked \textbf{his} bag. \\
\hline
\end{tabular}
\end{center}

\begin{exemplebox}[Des exemples traduits]
\begin{itemize}
\item \eng{“I will visit you tomorrow,” John told me.} $\rightarrow$ \eng{John told me he would visit me the next day.} « John m'a dit qu'il me rendrait visite le lendemain. »
\item \eng{“We bought these books here yesterday,” they said.} $\rightarrow$ \eng{They said they had bought those books there the day before.} « Ils ont dit qu'ils avaient acheté ces livres-là là-bas la veille. »
\item \eng{“I have lost my keys,” said my brother.} $\rightarrow$ \eng{My brother said he had lost his keys.} « Mon frère a dit qu'il avait perdu ses clés. »
\item \eng{“You can use my phone,” she told us.} $\rightarrow$ \eng{She told us we could use her phone.} « Elle nous a dit que nous pouvions utiliser son téléphone. »
\end{itemize}
\end{exemplebox}

\courssub{Les expressions de temps : le « recul » temporel}

Quand les paroles sont rapportées \emph{plus tard}, les repères de temps doivent reculer, exactement comme les temps verbaux. En français, tu fais déjà la même chose : « je viendrai \emph{demain} » rapporté devient « il viendrait \emph{le lendemain} ». L'anglais suit la même logique — mais il est plus systématique que le français.

\begin{propriete}[Changements des expressions de temps]
Quand le moment du rapport est différent du moment de l'énonciation :
\begin{itemize}
\item \eng{now} $\rightarrow$ \eng{then} ;
\item \eng{today} $\rightarrow$ \eng{that day} ;
\item \eng{yesterday} $\rightarrow$ \eng{the day before} ou \eng{the previous day} ;
\item \eng{tomorrow} $\rightarrow$ \eng{the next day} ou \eng{the following day} ;
\item \eng{ago} $\rightarrow$ \eng{before} (\eng{two days ago} $\rightarrow$ \eng{two days before}) ;
\item \eng{last week / month / year} $\rightarrow$ \eng{the week before / the previous week}, etc. ;
\item \eng{next week / year} $\rightarrow$ \eng{the following week / year} ;
\item \eng{tonight} $\rightarrow$ \eng{that night}.
\end{itemize}
\textbf{Cas où rien ne change :} si tu rapportes les paroles \emph{le même jour} et \emph{au même endroit}, les expressions peuvent rester identiques : \eng{“I am busy today,” she said.} (rapporté le jour même) $\rightarrow$ \eng{She said she is busy today.}
\end{propriete}

\courssub{Les expressions de lieu et les démonstratifs}

Le même raisonnement s'applique à l'espace : ce qui était « ici » pour le locuteur devient « là-bas » pour le rapporteur, et ce qui était « proche » (\eng{this/these}) devient « éloigné » (\eng{that/those}).

\begin{propriete}[Lieu et démonstratifs]
\begin{itemize}
\item \eng{here} $\rightarrow$ \eng{there} : \eng{“I live here,” he said.} $\rightarrow$ \eng{He said he lived there.} « Il a dit qu'il vivait là-bas. »
\item \eng{this} $\rightarrow$ \eng{that} : \eng{“This exercise is easy,” she said.} $\rightarrow$ \eng{She said that exercise was easy.}
\item \eng{these} $\rightarrow$ \eng{those} : \eng{“These mangoes are ripe,” he said.} $\rightarrow$ \eng{He said those mangoes were ripe.}
\end{itemize}
En français, la distinction \emph{ceci/cela} est moins marquée ; en anglais, le changement \eng{this} $\rightarrow$ \eng{that} est quasi obligatoire dans le discours rapporté.
\end{propriete}

\begin{popfigure}
\begin{tikzpicture}[
  dir/.style={rectangle, rounded corners=4pt, draw=popblue, fill=popblueL, thick, align=center, font=\small, minimum width=3.6cm, minimum height=0.7cm},
  ind/.style={rectangle, rounded corners=4pt, draw=poporange, fill=poporangeL, thick, align=center, font=\small, minimum width=4.6cm, minimum height=0.7cm},
  head/.style={font=\small\bfseries, text=popdark}
]
\node[head] at (-2.2,0.3) {Direct speech};
\node[head] at (3.4,0.3) {Indirect speech};
\foreach \y/\a/\b in {
  -0.6/now/then,
  -1.4/today/that day,
  -2.2/yesterday/the day before,
  -3.0/tomorrow/the next day,
  -3.8/two days ago/two days before,
  -4.6/last week/the week before,
  -5.4/next year/the following year,
  -6.2/tonight/that night,
  -7.0/here/there,
  -7.8/this and these/that and those}
{
  \node[dir] at (-2.2,\y) {\a};
  \node[ind] at (3.4,\y) {\b};
  \draw[-{Stealth}, thick, poppurple] (0.1,\y) -- (1.0,\y);
}
\end{tikzpicture}
\legende{Time and place changes : du discours direct au discours rapporté.}
\end{popfigure}

\begin{center}
\begin{tabularx}{\linewidth}{|X|X|X|}
\hline
\rowcolor{popblueL} \textbf{Français} & \textbf{English (direct)} & \textbf{English (reported)} \\
\hline
maintenant & now & then \\
\hline
aujourd'hui / ce jour-là & today & that day \\
\hline
hier / la veille & yesterday & the day before \\
\hline
demain / le lendemain & tomorrow & the next day \\
\hline
il y a deux jours & two days ago & two days before \\
\hline
la semaine dernière / précédente & last week & the previous week \\
\hline
l'année prochaine / suivante & next year & the following year \\
\hline
ce soir / ce soir-là & tonight & that night \\
\hline
ici / là-bas & here & there \\
\hline
ceci, ces / cela, ces...-là & this, these & that, those \\
\hline
\end{tabularx}
\end{center}

\courssub{Méthode, entraînement et pièges}

\begin{methode}[Transformer une phrase du discours direct au discours rapporté]
\begin{enumerate}
\item Mets le \textbf{verbe introducteur} au passé : \eng{say} $\rightarrow$ \eng{said}, \eng{tell} $\rightarrow$ \eng{told} (+ complément de personne).
\item \textbf{Supprime les guillemets} et les deux-points.
\item Ajoute \eng{that} (\textbf{facultatif}) après le verbe introducteur.
\item Applique le \textbf{backshift} des temps verbaux (section précédente).
\item Adapte les \textbf{pronoms et possessifs} au point de vue du rapporteur.
\item Adapte les expressions de \textbf{temps et de lieu} (\eng{tomorrow} $\rightarrow$ \eng{the next day}, \eng{here} $\rightarrow$ \eng{there}...).
\end{enumerate}
\end{methode}

\begin{exoresolu}[Transformation guidée]
Énoncé — Mets au discours rapporté : \eng{“I will show you these photos here tomorrow,” Mary told me.}
\tcblower
Solution détaillée :
\begin{enumerate}
\item Verbe introducteur au passé : \eng{Mary told me} (déjà au passé).
\item On supprime les guillemets.
\item On ajoute \eng{that} (facultatif) : \eng{Mary told me (that)...}
\item Backshift : \eng{will show} $\rightarrow$ \eng{would show}.
\item Pronoms : \eng{I} $\rightarrow$ \eng{she} (Mary) ; \eng{you} $\rightarrow$ \eng{me} (c'est à moi qu'elle parlait).
\item Temps et lieu : \eng{these photos} $\rightarrow$ \eng{those photos} ; \eng{here} $\rightarrow$ \eng{there} ; \eng{tomorrow} $\rightarrow$ \eng{the next day}.
\end{enumerate}
\textbf{Résultat :} \eng{Mary told me (that) she would show me those photos there the next day.} « Mary m'a dit qu'elle me montrerait ces photos-là là-bas le lendemain. »
\end{exoresolu}

\begin{attention}[Erreur fréquente des francophones : garder les mots de temps inchangés]
En français parlé, on entend parfois « il a dit qu'il viendrait demain ». En anglais, c'est une \textbf{erreur} : les expressions de temps doivent reculer comme les temps verbaux.
\begin{itemize}
\item \emph{Faux :} \eng{He said he would come tomorrow.}
\item \emph{Correct :} \eng{He said he would come the next day.} « Il a dit qu'il viendrait le lendemain. »
\item \emph{Faux :} \eng{She said she had seen him yesterday.}
\item \emph{Correct :} \eng{She said she had seen him the day before.}
\end{itemize}
Même vigilance pour \eng{here} : ne dis pas \eng{He said he lived here} si tu n'es pas sur les lieux ; dis \eng{He said he lived there}.
\end{attention}

\begin{exercice}[À toi de jouer]
Mets les phrases suivantes au discours rapporté :
\begin{enumerate}
\item \eng{“I am writing my essay now,” said Peter.}
\item \eng{“We visited our grandmother last week,” the twins said.}
\item \eng{“I will finish this exercise tonight,” she told me.}
\item \eng{“You left your bag here two days ago,” the teacher told him.}
\end{enumerate}
\end{exercice}

\begin{corrige}[Exercice — À toi de jouer]
\begin{enumerate}
\item \eng{Peter said (that) he was writing his essay then.} « Peter a dit qu'il écrivait sa rédaction à ce moment-là. »
\item \eng{The twins said (that) they had visited their grandmother the week before.} « Les jumeaux ont dit qu'ils avaient rendu visite à leur grand-mère la semaine précédente. »
\item \eng{She told me (that) she would finish that exercise that night.} « Elle m'a dit qu'elle finirait cet exercice-là ce soir-là. »
\item \eng{The teacher told him (that) he had left his bag there two days before.} « Le professeur lui a dit qu'il avait laissé son sac là-bas deux jours auparavant. »
\end{enumerate}
\end{corrige}

\begin{aretenir}[L'essentiel de la section]
\begin{itemize}
\item Les pronoms et possessifs s'adaptent au \textbf{point de vue du rapporteur} : \eng{I} $\rightarrow$ \eng{he/she}, \eng{my} $\rightarrow$ \eng{his/her}, \eng{we} $\rightarrow$ \eng{they}, \eng{our} $\rightarrow$ \eng{their}.
\item Les expressions de temps reculent : \eng{now} $\rightarrow$ \eng{then}, \eng{today} $\rightarrow$ \eng{that day}, \eng{yesterday} $\rightarrow$ \eng{the day before}, \eng{tomorrow} $\rightarrow$ \eng{the next day}, \eng{ago} $\rightarrow$ \eng{before}.
\item Lieu et démonstratifs : \eng{here} $\rightarrow$ \eng{there}, \eng{this/these} $\rightarrow$ \eng{that/those}.
\item Pas de changement si le rapport est fait \emph{au même moment et au même endroit}.
\end{itemize}
\end{aretenir}

\coursec{Rapporter questions, ordres et conseils}

Dans la section précédente, nous avons vu comment les pronoms et les expressions de temps et de lieu changent lorsqu'on rapporte des paroles (\eng{I} devient \eng{he}, \eng{today} devient \eng{that day}, etc.). Mais jusqu'ici, nous n'avons rapporté que des \cle{phrases déclaratives} avec \eng{say} ou \eng{tell} + \eng{that}. Que se passe-t-il lorsque la personne pose une \cle{question}, donne un \cle{ordre} ou un \cle{conseil} ? C'est l'objet de cette section : chaque type de phrase a sa propre structure de discours indirect.

\courssub{Observation : un dialogue au dispensaire de Buea}

\begin{textebox}[At the health centre in Buea — a dialogue]
Nurse: “Are you feeling better today, Mr. Ekane?”

Mr. Ekane: “Not really. My head still hurts.”

Nurse: “When did the pain start?”

Mr. Ekane: “Three days ago, after the football match.”

Nurse: “Take these tablets twice a day, and don't carry heavy loads.”

Mr. Ekane: “Should I stop working?”

Nurse: “Yes. Rest for at least two days. And drink a lot of water.”

\medskip
Later, Mr. Ekane reported the conversation to his wife:

“The nurse asked me \textbf{if I was feeling} better. She asked me \textbf{when the pain had started}. Then she \textbf{told me to take} the tablets twice a day and \textbf{not to carry} heavy loads. I asked her \textbf{if I should stop} working, and she \textbf{advised me to rest} for at least two days.”
\end{textebox}

\textbf{Glossaire :} \emph{nurse} (nom, infirmière) ; \emph{to hurt} (verbe, faire mal) ; \emph{tablets} (nom, comprimés) ; \emph{load} (nom, charge) ; \emph{to rest} (verbe, se reposer).

Observez les phrases en gras : la question \eng{Are you feeling better?} devient \eng{if I was feeling better} ; l'ordre \eng{Take these tablets} devient \eng{told me to take}. Remarquez deux choses essentielles : dans les questions rapportées, \textbf{l'inversion a disparu} et il n'y a \textbf{plus de point d'interrogation}. Décortiquons chaque cas.

\courssub{Les questions fermées (oui/non) : ask + if/whether}

\begin{propriete}[Rapporter une question fermée]
Pour rapporter une question à laquelle on répond par \eng{yes} ou \eng{no}, on utilise :

\begin{center}
\textbf{ask (+ complément de personne) + if / whether + sujet + verbe}
\end{center}

\begin{itemize}
\item \eng{if} ou \eng{whether} (= « si ») remplace le point d'interrogation ;
\item on retrouve l'\textbf{ordre de la phrase affirmative} : sujet + verbe, \textbf{sans inversion} ;
\item les auxiliaires \eng{do / does / did} \textbf{disparaissent} (le verbe principal reprend la marque du temps) ;
\item le \cle{backshift} s'applique normalement (present simple $\rightarrow$ past simple, etc.) ;
\item \textbf{pas de point d'interrogation} : c'est une phrase déclarative.
\end{itemize}
\end{propriete}

\begin{exemplebox}[Questions fermées rapportées]
\eng{“Do you like football?” he asked me.} $\rightarrow$ \eng{He asked me if I liked football.} « Il m'a demandé si j'aimais le football. »

\eng{“Are you ready?” the coach asked the players.} $\rightarrow$ \eng{The coach asked the players if they were ready.} « L'entraîneur a demandé aux joueurs s'ils étaient prêts. »

\eng{“Have you finished your homework?” mother asked.} $\rightarrow$ \eng{Mother asked whether I had finished my homework.} « Maman a demandé si j'avais fini mes devoirs. »

\eng{“Will the bus leave on time?” a passenger asked.} $\rightarrow$ \eng{A passenger asked if the bus would leave on time.} « Un passager a demandé si le car partirait à l'heure. »
\end{exemplebox}

\begin{attention}[Ne gardez pas l'inversion !]
L'erreur n°1 des francophones est de conserver l'inversion sujet–verbe de la question directe :

✗ \eng{He asked me if did I like football.} / ✗ \eng{He asked me if was I ready.}

✓ \eng{He asked me if I liked football.} / ✓ \eng{He asked me if I was ready.}

En français, on dit naturellement « il m'a demandé \emph{si j'aimais} le football » — sans inversion. Faites de même en anglais : après \eng{if/whether}, ordre sujet + verbe.
\end{attention}

\courssub{Les questions ouvertes (wh-questions) : ask + mot interrogatif}

\begin{propriete}[Rapporter une question ouverte]
Pour rapporter une question introduite par un mot interrogatif (\eng{where, when, why, who, what, how, how much, how old}...), on garde ce mot interrogatif comme connecteur :

\begin{center}
\textbf{ask (+ personne) + wh-word + sujet + verbe}
\end{center}

Là encore : \textbf{pas d'inversion}, \textbf{pas d'auxiliaire do/does/did}, \textbf{pas de point d'interrogation}, et backshift des temps.
\end{propriete}

\begin{exemplebox}[Questions ouvertes rapportées]
\eng{“Where do you live?” she asked me.} $\rightarrow$ \eng{She asked me where I lived.} « Elle m'a demandé où j'habitais. »

\eng{“Why are you late?” the principal asked the student.} $\rightarrow$ \eng{The principal asked the student why he was late.} « Le proviseur a demandé à l'élève pourquoi il était en retard. »

\eng{“How much does this bag of rice cost?” the customer asked.} $\rightarrow$ \eng{The customer asked how much the bag of rice cost.} « Le client a demandé combien coûtait le sac de riz. »

\eng{“What have you done with my keys?” he asked his sister.} $\rightarrow$ \eng{He asked his sister what she had done with his keys.} « Il a demandé à sa sœur ce qu'elle avait fait de ses clés. »
\end{exemplebox}

\begin{attention}[✗ where did I live → ✓ where I lived]
Même piège que pour les questions fermées, aggravé par l'auxiliaire :

✗ \eng{He asked me where did I live.}

✓ \eng{He asked me where I lived.}

L'auxiliaire \eng{do/does/did} n'existe que pour construire une question \emph{directe}. Dans la question rapportée, il disparaît et le verbe porte seul la marque du passé : \eng{do you live} $\rightarrow$ \eng{I lived}.
\end{attention}

\courssub{Ordres et demandes : tell / order / ask + objet + to-infinitive}

\begin{propriete}[Rapporter un ordre ou une demande]
Les impératifs se rapportent avec l'\textbf{infinitif en to} :

\begin{center}
\textbf{tell / order / ask + complément de personne + to + base verbale}
\end{center}

\begin{itemize}
\item \eng{tell} : verbe neutre, le plus fréquent (« dire de ») ;
\item \eng{order} : ordre autoritaire (militaire, juge, professeur strict) ;
\item \eng{ask} : demande polie (« demander de »).
\end{itemize}

Pour un \textbf{ordre négatif}, on place \eng{not} \textbf{avant} l'infinitif :

\begin{center}
\textbf{tell + personne + not to + base verbale}
\end{center}
\end{propriete}

\begin{exemplebox}[Ordres rapportés]
\eng{“Open your books!” the teacher said.} $\rightarrow$ \eng{The teacher told us to open our books.} « Le professeur nous a dit d'ouvrir nos livres. »

\eng{“Don't make noise!” she said to them.} $\rightarrow$ \eng{She told them not to make noise.} « Elle leur a dit de ne pas faire de bruit. »

\eng{“Stand up and give me your pass!” the policeman said.} $\rightarrow$ \eng{The policeman ordered him to stand up and give him his pass.} « Le policier lui a ordonné de se lever et de lui donner son laissez-passer. »

\eng{“Please help me carry this box,” she said to me.} $\rightarrow$ \eng{She asked me to help her carry the box.} « Elle m'a demandé de l'aider à porter le carton. »

\eng{“Don't touch the hot pot!” mother said.} $\rightarrow$ \eng{Mother told the children not to touch the hot pot.} « Maman a dit aux enfants de ne pas toucher la marmite chaude. »
\end{exemplebox}

\begin{attention}[say ne prend jamais de complément de personne]
✗ \eng{He said me to sit down.} / ✗ \eng{He said me that he was tired.}

✓ \eng{He told me to sit down.} / ✓ \eng{He told me that he was tired.} ou \eng{He said that he was tired.}

Retenez : \eng{say something (to somebody)} mais \eng{tell somebody something}. Dès qu'il y a un complément de personne (\eng{me, us, the students}...), utilisez \eng{tell} (ou \eng{ask, order, advise}), jamais \eng{say}.
\end{attention}

\courssub{Conseils et suggestions : advise et suggest}

\begin{propriete}[Rapporter un conseil : advise + objet + to-infinitive]
\eng{“You should rest.”} $\rightarrow$ \eng{The doctor advised me to rest.} « Le médecin m'a conseillé de me reposer. »

Forme : \textbf{advise + personne + (not) to + base verbale}. Le modal \eng{should} disparaît, absorbé par \eng{advise}.
\end{propriete}

\begin{propriete}[Rapporter une suggestion : suggest]
\eng{suggest} se construit différemment — \textbf{jamais} avec la structure « suggest + personne + to » :

\begin{itemize}
\item \textbf{suggest + V-ing} : \eng{“Let's go to the market.”} $\rightarrow$ \eng{He suggested going to the market.} « Il a suggéré d'aller au marché. »
\item \textbf{suggest + that + sujet + (should) + base verbale} : \eng{He suggested that we (should) go to the market.}
\end{itemize}

✗ \eng{He suggested me to go.} — construction impossible en anglais !
\end{propriete}

\begin{exemplebox}[Conseils et suggestions]
\eng{“You should drink more water,” the nurse said.} $\rightarrow$ \eng{The nurse advised him to drink more water.} « L'infirmière lui a conseillé de boire plus d'eau. »

\eng{“Don't stay in the sun,” she said.} $\rightarrow$ \eng{She advised me not to stay in the sun.} « Elle m'a conseillé de ne pas rester au soleil. »

\eng{“Let's plant trees around the school,” the prefect said.} $\rightarrow$ \eng{The prefect suggested planting trees around the school.} / \eng{The prefect suggested that we (should) plant trees around the school.} « Le préfet a suggéré de planter des arbres autour de l'école. »
\end{exemplebox}

\begin{savaistu}[If ou whether ?]
\eng{whether} est légèrement plus formel que \eng{if} et très apprécié dans les rédactions du Bac : \eng{She asked whether the exam would be difficult.} « Elle a demandé si l'examen serait difficile. » Attention : \eng{whether} se prononce « oué-zeur », exactement comme \eng{weather} (le temps qu'il fait) — deux homophones à ne pas confondre à l'écrit !
\end{savaistu}

\courssub{Tableau récapitulatif et organigramme de décision}

\begin{center}
\begin{tabularx}{\linewidth}{|X|X|X|}
\hline
\rowcolor{popblueL} Type de phrase directe & Structure rapportée & Exemple \\
\hline
Déclaration & say / tell + (that) + sujet + verbe & \eng{He told me (that) he was tired.} \\
\hline
Question oui/non & ask + if/whether + sujet + verbe & \eng{She asked if I was ready.} \\
\hline
Question wh- & ask + wh- + sujet + verbe & \eng{He asked where I lived.} \\
\hline
Ordre / demande & tell / order / ask + pers. + to-inf. & \eng{She told us to sit down.} \\
\hline
Ordre négatif & tell + pers. + not to-inf. & \eng{He told them not to run.} \\
\hline
Conseil & advise + pers. + to-inf. & \eng{The doctor advised me to rest.} \\
\hline
Suggestion & suggest + V-ing / that + (should) & \eng{He suggested going home.} \\
\hline
\end{tabularx}
\end{center}

\begin{popfigure}
\begin{tikzpicture}[
  node distance=0.9cm and 1.4cm,
  decision/.style={diamond, draw=popdark, fill=poppurpleL, text width=3.2cm, align=center, inner sep=1pt, font=\small},
  box/.style={rectangle, rounded corners, draw=popdark, fill=popblueL, text width=4.6cm, align=center, font=\small, inner sep=4pt},
  arr/.style={-{Stealth}, thick, popdark}]
\node[decision] (start) {Type de phrase à rapporter ?};
\node[box, above right=0.6cm and 1.6cm of start, align=center] (decl){Déclaration \\ \eng{say / tell + (that) + sujet + verbe}};
\node[box, right=1.6cm of start, yshift=0.9cm, align=center] (yesno){Question oui/non \\ \eng{ask + if/whether + sujet + verbe}};
\node[box, right=1.6cm of start, yshift=-0.9cm, align=center] (wh){Question wh- \\ \eng{ask + wh- + sujet + verbe}};
\node[box, below right=0.6cm and 1.6cm of start, align=center] (order){Ordre / conseil \\ \eng{tell / advise + pers. + to-infinitif}};
\draw[arr] (start) -- (decl);
\draw[arr] (start) -- (yesno);
\draw[arr] (start) -- (wh);
\draw[arr] (start) -- (order);
\end{tikzpicture}
\legende{Organigramme de décision : choisir la bonne structure de discours indirect.}
\end{popfigure}

\courssub{Entraînement guidé}

\begin{exoresolu}[Transformation au discours indirect]
Mettez au discours indirect : \eng{“Why didn't you come to school yesterday?” the teacher asked Ngono.}

\tcblower
\textbf{Étape 1 — Type de phrase :} question ouverte introduite par \eng{why} $\rightarrow$ structure \eng{ask + why + sujet + verbe}.

\textbf{Étape 2 — Suppression de l'inversion et de l'auxiliaire :} \eng{why didn't you come} $\rightarrow$ \eng{why ... had not come} (past simple $\rightarrow$ past perfect, backshift).

\textbf{Étape 3 — Pronoms et temps :} \eng{you} $\rightarrow$ \eng{he} (Ngono) ; \eng{yesterday} $\rightarrow$ \eng{the day before}.

\textbf{Solution :} \eng{The teacher asked Ngono why he had not come to school the day before.} « Le professeur a demandé à Ngono pourquoi il n'était pas venu à l'école la veille. »
\end{exoresolu}

\begin{exercice}[À vous de jouer]
Rapportez les paroles suivantes :

\begin{enumerate}
\item \eng{“Do you play handball?” she asked me.}
\item \eng{“Where is the nearest pharmacy?” the tourist asked.}
\item \eng{“Don't cross the road here!” the policeman said to the children.}
\item \eng{“You should see a doctor,” my friend said.}
\item \eng{“Let's clean the classroom,” the class captain said.}
\end{enumerate}
\end{exercice}

\begin{corrige}[Exercice — À vous de jouer]
\begin{enumerate}
\item \eng{She asked me if/whether I played handball.} « Elle m'a demandé si je jouais au handball. »
\item \eng{The tourist asked where the nearest pharmacy was.} « Le touriste a demandé où était la pharmacie la plus proche. » (Pas d'inversion : \eng{where ... was}, pas \eng{where was...}.)
\item \eng{The policeman told the children not to cross the road there.} « Le policier a dit aux enfants de ne pas traverser la route à cet endroit. » (\eng{not to} + \eng{here} $\rightarrow$ \eng{there}.)
\item \eng{My friend advised me to see a doctor.} « Mon ami m'a conseillé de voir un médecin. »
\item \eng{The class captain suggested cleaning the classroom.} ou \eng{The class captain suggested that we (should) clean the classroom.} « Le chef de classe a suggéré de nettoyer la salle. »
\end{enumerate}
\end{corrige}

\begin{aretenir}[L'essentiel de la section]
\begin{itemize}
\item Question oui/non $\rightarrow$ \eng{ask + if/whether + sujet + verbe} ; question wh- $\rightarrow$ \eng{ask + wh- + sujet + verbe}. Dans les deux cas : \textbf{ni inversion, ni do/does/did, ni point d'interrogation}.
\item Ordre $\rightarrow$ \eng{tell/order/ask + personne + to-infinitive} ; ordre négatif $\rightarrow$ \eng{not to-infinitive}.
\item Conseil $\rightarrow$ \eng{advise + personne + to-infinitive} ; suggestion $\rightarrow$ \eng{suggest + V-ing} ou \eng{suggest that + (should)}.
\item \eng{say} ne prend jamais de complément de personne : \eng{tell me}, jamais \eng{say me}.
\end{itemize}
\end{aretenir}

\coursec{Verbes introducteurs et nuances de sens}

Dans la section précédente, nous avons appris à rapporter les questions (\eng{ask}, \eng{wonder}), les ordres et les conseils (\eng{tell}, \eng{order}, \eng{advise} + objet + to-infinitive). Mais dire simplement \eng{he said...} à chaque phrase est pauvre et répétitif. L'anglais possède une famille très riche de \cle{verbes introducteurs} (\eng{reporting verbs}) qui précisent l'\textbf{intention} du locuteur : promettre, accuser, refuser, conseiller, nier... Choisir le bon verbe, c'est rendre le discours rapporté plus vivant, plus précis — et c'est une compétence directement valorisée au Bac.

\courssub{Observation : un mini-dialogue rapporté}

\begin{textebox}[At the market in Mokolo — a reported conversation]
Yesterday at Mokolo market, I met my friend Ali. “I'll pay you back next week,” \textbf{promised} Ali. He \textbf{explained} that his salary had been delayed. Then a shopkeeper shouted at a boy: “You stole my pen!” She \textbf{accused} him of stealing her pen, but the boy \textbf{denied} taking anything. “You should ask before you touch things,” his mother said. She \textbf{advised} him to be more careful. “Don't touch the wire!” the electrician \textbf{warned} us not to touch the wire near the stall. Finally, the shopkeeper \textbf{apologised} and \textbf{admitted} that she had found her pen in her own bag. “I suggest we all calm down,” \textbf{suggested} a customer. Everyone \textbf{agreed} that it had been a misunderstanding.
\end{textebox}

\begin{definition}[Glossaire du texte]
\begin{itemize}
\item \eng{to pay back} (v.) : rembourser
\item \eng{salary} (n.) : salaire
\item \eng{delayed} (adj.) : retardé
\item \eng{to steal, stole, stolen} (v.) : voler, dérober
\item \eng{to deny} (v.) : nier
\item \eng{wire} (n.) : fil (électrique)
\item \eng{stall} (n.) : étal, stand
\item \eng{to apologise} (v.) : s'excuser
\item \eng{misunderstanding} (n.) : malentendu
\end{itemize}
\end{definition}

Observe bien : chaque verbe en gras dans le texte (\eng{promised}, \eng{explained}, \eng{accused}, \eng{denied}, \eng{advised}, \eng{warned}, \eng{admitted}, \eng{suggested}) remplace un simple \eng{said} et ajoute une \textbf{nuance de sens}. C'est exactement l'objet de cette leçon.

\courssub{Say et tell : la différence de construction}

C'est la distinction la plus fondamentale — et la plus fautive chez les francophones.

\begin{propriete}[Règle SAY / TELL]
\begin{itemize}
\item \textbf{SAY + (that) + phrase} : on dit quelque chose, sans nommer obligatoirement la personne.\\
\eng{He said (that) he was ill.} « Il a dit qu'il était malade. »
\item \textbf{TELL + personne + (that) + phrase} : on dit quelque chose \textbf{à quelqu'un} ; le complément de personne est \textbf{obligatoire}.\\
\eng{He told me (that) he was ill.} « Il m'a dit qu'il était malade. »
\end{itemize}
Exception utile : \eng{say} peut introduire une personne seulement avec \eng{to} : \eng{He said to me that he was tired.} (correct mais lourd ; préfère \eng{told me}).
\end{propriete}

\begin{attention}[Erreurs fréquentes des francophones]
\begin{itemize}
\item ✗ \eng{He said me that he was ill.} → ✓ \eng{He told me that he was ill.} (jamais de personne directement après \eng{say}).
\item ✗ \eng{He told that he was tired.} → ✓ \eng{He told \textbf{us} that he was tired.} (\eng{tell} exige un complément de personne).
\item Expressions figées avec \eng{tell} : \eng{tell the truth} (dire la vérité), \eng{tell a lie} (mentir), \eng{tell the time} (dire l'heure), \eng{tell a story} (raconter une histoire). Mais : \eng{say hello}, \eng{say goodbye}, \eng{say sorry}.
\end{itemize}
\end{attention}

\courssub{Carte mentale : la famille des verbes introducteurs}

\begin{popfigure}
\begin{tikzpicture}[mindmap, grow cyclic, every node/.style={concept, text=white, font=\small\bfseries}, root concept/.append style={concept color=popblue, font=\large\bfseries}, level 1/.append style={level distance=3.6cm, sibling angle=90}, level 2/.append style={level distance=2.4cm, sibling angle=45, font=\scriptsize\bfseries}]
\node[root concept]{Reporting verbs}
  child[concept color=popgreen]{ node {declaring} 
    child { node {announce} }
    child { node {explain} }
    child { node {promise} }
    child { node {warn} }
  }
  child[concept color=poporange]{ node {asking}
    child { node {ask} }
    child { node {wonder} }
    child { node {inquire} }
  }
  child[concept color=poppurple]{ node {ordering / advising}
    child { node {order} }
    child { node {advise} }
    child { node {beg} }
    child { node {remind} }
  }
  child[concept color=poppink]{ node {+ gerund}
    child { node {suggest} }
    child { node {admit} }
    child { node {deny} }
  };
\end{tikzpicture}
\legende{Carte mentale des verbes introducteurs : quatre grandes familles selon la construction.}
\end{popfigure}

\courssub{Les verbes pour déclarer (declaring)}

Ces verbes rapportent des déclarations, avec la construction \textbf{verbe + (that) + phrase} (avec recul des temps, vu à la section 2).

\begin{center}
\begin{tabularx}{\linewidth}{|l|X|X|}
\hline
\rowcolor{popblueL} Verbe & Exemple anglais & Traduction \\
\hline
announce & The headmaster announced that the exams would start on Monday. & Le proviseur a annoncé que les examens commenceraient lundi. \\
\hline
explain & She explained that the water was polluted. & Elle a expliqué que l'eau était polluée. \\
\hline
state & The minister stated that the hospital would be renovated. & Le ministre a déclaré que l'hôpital serait rénové. \\
\hline
reply & He replied that he was too busy. & Il a répondu qu'il était trop occupé. \\
\hline
add & “And bring your own cup,” she added. & « Et apporte ton propre gobelet », a-t-elle ajouté. \\
\hline
complain & The patients complained that the queue was too long. & Les patients se sont plaints que la file était trop longue. \\
\hline
promise & Ali promised that he would pay me back. & Ali a promis qu'il me rembourserait. \\
\hline
warn & The doctor warned that the disease was spreading. & Le médecin a averti que la maladie se propageait. \\
\hline
\end{tabularx}
\end{center}

\begin{attention}[Construction de explain]
✗ \eng{She explained me the problem.} → ✓ \eng{She explained the problem \textbf{to} me.} Le verbe \eng{explain} ne prend jamais de complément de personne direct : il faut \eng{explain something to somebody}.
\end{attention}

\courssub{Les verbes pour demander (asking)}

Rappel de la section précédente : les questions rapportées perdent l'inversion et le point d'interrogation.

\begin{itemize}
\item \eng{ask} : \eng{“Where do you live?” → She asked me where I lived.} « Elle m'a demandé où j'habitais. »
\item \eng{wonder} (se demander — le locuteur se pose la question à lui-même) : \eng{“Why is the river so dirty?” → He wondered why the river was so dirty.} « Il s'est demandé pourquoi la rivière était si sale. »
\item \eng{want to know} : \eng{They wanted to know if the vaccine was free.} « Ils voulaient savoir si le vaccin était gratuit. »
\item \eng{inquire} (soutenu, administratif) : \eng{She inquired whether the clinic was open on Sundays.} « Elle s'est renseignée pour savoir si le dispensaire était ouvert le dimanche. »
\end{itemize}

\courssub{Les verbes pour ordonner et conseiller : verbe + objet + to-infinitive}

\begin{propriete}[Construction verbe + objet + to-infinitive]
Pour rapporter un ordre, un conseil, une invitation ou une prière :\\
\textbf{verbe + complément de personne + (not) to + base verbale}\\
\eng{“Don't touch the wire!” → The electrician warned us \textbf{not to touch} the wire.} « L'électricien nous a avertis de ne pas toucher le fil. »
\end{propriete}

\begin{center}
\begin{tabularx}{\linewidth}{|l|X|X|}
\hline
\rowcolor{popblueL} Verbe & Exemple anglais & Traduction \\
\hline
tell & The nurse told me to sit down. & L'infirmière m'a dit de m'asseoir. \\
\hline
order & The officer ordered the driver to stop. & L'officier a ordonné au chauffeur de s'arrêter. \\
\hline
command & The captain commanded his men to advance. & Le capitaine a commandé à ses hommes d'avancer. \\
\hline
advise & The doctor advised her to drink more water. & Le médecin lui a conseillé de boire plus d'eau. \\
\hline
beg & The child begged his father to buy him a ball. & L'enfant a supplié son père de lui acheter un ballon. \\
\hline
encourage & The teacher encouraged us to plant trees. & Le professeur nous a encouragés à planter des arbres. \\
\hline
invite & They invited us to attend the ceremony. & Ils nous ont invités à assister à la cérémonie. \\
\hline
remind & She reminded me to take my medicine. & Elle m'a rappelé de prendre mes médicaments. \\
\hline
warn & He warned them not to swim in the river. & Il les a avertis de ne pas nager dans la rivière. \\
\hline
\end{tabularx}
\end{center}

\courssub{Les verbes suivis du gérondif}

Certains verbes introducteurs se construisent avec un \textbf{gérondif} (-ing) et non avec une subordonnée. C'est une construction très élégante pour le Bac.

\begin{propriete}[Verbe + gérondif]
\begin{itemize}
\item \eng{suggest + -ing} : \eng{“Let's plant trees.” → He suggested \textbf{planting} trees.} « Il a suggéré de planter des arbres. »
\item \eng{admit + -ing} : \eng{“I broke the window.” → He admitted \textbf{breaking} the window.} « Il a admis avoir cassé la vitre. »
\item \eng{deny + -ing} : \eng{“I didn't take the money.” → She denied \textbf{taking} the money.} « Elle a nié avoir pris l'argent. »
\end{itemize}
Autre construction utile : \eng{accuse somebody of + -ing} : \eng{“You stole my pen!” → She accused him \textbf{of stealing} her pen.} « Elle l'a accusé d'avoir volé son stylo. »
\end{propriete}

\begin{attention}[Le piège de suggest]
✗ \eng{He suggested me to go.} — construction impossible en anglais !\\
✓ \eng{He suggested \textbf{that I (should) go}.} « Il a suggéré que j'y aille. »\\
✓ \eng{He suggested \textbf{going}.} « Il a suggéré d'y aller. »\\
Le français « il m'a suggéré de partir » se traduit donc par \eng{he suggested that I should leave}, jamais avec \eng{me to}.
\end{attention}

\courssub{Choisir le verbe selon l'intention du locuteur}

La même phrase de départ peut être rapportée par plusieurs verbes, selon ce que tu veux souligner. C'est un choix d'interprétation.

\begin{exemplebox}[Une phrase, plusieurs rapports]
Phrase directe : \eng{“I'll help you clean the schoolyard tomorrow.”}\\
\begin{itemize}
\item \eng{He \textbf{promised} to help us clean the schoolyard the next day.} (on insiste sur l'engagement)
\item \eng{He \textbf{offered} to help us clean the schoolyard the next day.} (on insiste sur la proposition volontaire)
\item \eng{He \textbf{said} he would help us clean the schoolyard the next day.} (neutre)
\end{itemize}
Autre exemple : \eng{“You never listen to me!”} → \eng{She \textbf{complained} that he never listened to her.} (reproche) ou \eng{She \textbf{accused} him of never listening to her.} (accusation).
\end{exemplebox}

\begin{methode}[Comment choisir le bon verbe introducteur]
\begin{enumerate}
\item Lis la phrase au discours direct et demande-toi : que \emph{fait} le locuteur ? Il promet ? Il refuse ? Il ordonne ? Il s'excuse ?
\item Identifie les indices : \eng{I'll...} (promesse), \eng{Don't...!} (ordre ou avertissement), \eng{You should...} (conseil), \eng{I'm sorry...} (excuse), \eng{I didn't do it!} (dénégation).
\item Choisis le verbe de la bonne famille : déclaration → \eng{+ (that)} ; question → \eng{ask/wonder + if ou mot interrogatif} ; ordre/conseil → \eng{+ objet + to-infinitive} ; suggestion, aveu, déni → \eng{+ gérondif}.
\item Applique ensuite le recul des temps et les changements de pronoms et d'expressions de temps (sections 2 et 3).
\end{enumerate}
\end{methode}

\begin{exoresolu}[Transformer avec le verbe introducteur adapté]
Énoncé — Rapporte chaque phrase au discours indirect en utilisant le verbe entre parenthèses.\\
1. \eng{“I will visit the health centre tomorrow,” said Mary.} (promise)\\
2. \eng{“Don't swim in this river,” the guide said to us.} (warn)\\
3. \eng{“I didn't break the thermometer,” said the boy.} (deny)\\
4. \eng{“You should eat more fruit,” the doctor said to me.} (advise)\\
5. \eng{“Why is the clinic closed?” she asked herself.} (wonder)
\tcblower
Solution détaillée —\\
1. \eng{Mary promised (that) she would visit the health centre the next day.} — \eng{will} recule en \eng{would} ; \eng{tomorrow} devient \eng{the next day}.\\
2. \eng{The guide warned us not to swim in that river.} — impératif négatif → \eng{warn + objet + not to + base verbale}.\\
3. \eng{The boy denied breaking the thermometer.} — dénégation → \eng{deny + gérondif}.\\
4. \eng{The doctor advised me to eat more fruit.} — conseil → \eng{advise + objet + to-infinitive}.\\
5. \eng{She wondered why the clinic was closed.} — question rapportée : pas d'inversion, \eng{is} recule en \eng{was}.
\end{exoresolu}

\begin{exercice}[À toi de jouer]
Rapporte les phrases suivantes avec un verbe introducteur varié et précis (évite \eng{said}).\\
1. \eng{“I'll bring the books back on Friday,” said Peter.}\\
2. \eng{“You took my umbrella!” she shouted at her brother.}\\
3. \eng{“Please, please, lend me your bicycle,” said John to his friend.}\\
4. \eng{“Let's clean the classroom before the inspector arrives,” said the class prefect.}\\
5. \eng{“Don't forget to boil the water,” the nurse said to the villagers.}
\end{exercice}

\begin{corrige}[Exercice : verbes introducteurs]
1. \eng{Peter promised to bring the books back on Friday.} (ou \eng{promised that he would bring...}).\\
2. \eng{She accused her brother of taking her umbrella.}\\
3. \eng{John begged his friend to lend him his bicycle.}\\
4. \eng{The class prefect suggested cleaning the classroom before the inspector arrived.} (ou \eng{suggested that they (should) clean...}).\\
5. \eng{The nurse reminded the villagers to boil the water.} (ou \eng{warned them not to forget...}).
\end{corrige}

\begin{aretenir}[L'essentiel de la section]
\begin{itemize}
\item \eng{say + (that)} ; \eng{tell + personne + (that)} — jamais \eng{say me}, jamais \eng{tell} sans personne.
\item Déclarer : \eng{announce, explain, state, reply, add, complain, promise, warn} + (that).
\item Demander : \eng{ask, wonder, want to know, inquire} + if/whether ou mot interrogatif.
\item Ordonner/conseiller : \eng{tell, order, command, advise, beg, encourage, invite, remind, warn} + objet + to-infinitive.
\item Gérondif : \eng{suggest, admit, deny + -ing} ; \eng{accuse somebody of + -ing}.
\item Choisis le verbe selon l'\textbf{intention} du locuteur, pas seulement ses mots.
\end{itemize}
\end{aretenir}

\begin{savaistu}[Un atout pour le Bac]
Dans les épreuves d'expression écrite du Bac (rédaction narrative, compte rendu, dialogue rapporté), les correcteurs valorisent la \textbf{richesse lexicale}. Une copie qui répète \eng{said, said, said} paraît pauvre ; une copie qui utilise \eng{promised, warned, advised, denied, admitted, suggested} montre une maîtrise fine de la langue. Entraîne-toi à glisser deux ou trois verbes introducteurs variés dans chaque production écrite : c'est un investissement à haut rendement !
\end{savaistu}

\coursec{Comparaison français/anglais et entraînement guidé}

Après avoir étudié les verbes introducteurs (\eng{say}, \eng{tell}, \eng{ask}, \eng{advise}, \eng{warn}...) et leurs nuances de sens, il est temps de confronter les deux systèmes : le français et l'anglais rapportent les paroles de manière très proche, mais avec des différences précises qui sont la source de la plupart des erreurs des élèves francophones. Cette dernière section te donne les clés pour transformer n'importe quel énoncé en discours indirect, étape par étape, comme à l'examen.

\courssub{Similitude fondamentale : les deux langues reculent les temps}

Bonne nouvelle : le français et l'anglais fonctionnent de la même façon pour les temps. Quand le verbe introducteur est au passé, les deux langues reculent les temps de la parole rapportée.

\begin{exemplebox}[Le même recul dans les deux langues]
\eng{He said, “I am tired.”} → \eng{He said he was tired.}\\
« Il a dit : “Je suis fatigué.” » → « Il a dit qu'il \textbf{était} fatigué. »\\[0.3em]
\eng{She said, “I will come.”} → \eng{She said she would come.}\\
« Elle a dit : “Je viendrai.” » → « Elle a dit qu'elle \textbf{viendrait}. »
\end{exemplebox}

\begin{savaistu}[La concordance des temps, ton alliée secrète]
Tu as appris au collège la concordance des temps en français : « il a dit qu'il \emph{viendrait} », « elle a expliqué qu'elle \emph{avait fini} ». Ce mécanisme est exactement le même en anglais : \eng{will} → \eng{would}, \eng{have finished} → \eng{had finished}. Quand tu hésites sur une réponse en anglais, traduis-la mentalement en français correct : si le français recule le temps, l'anglais le recule aussi. Utilise ce réflexe pour vérifier tes productions à l'examen.
\end{savaistu}

\begin{exemplebox}[Présent : pas de backshift, dans les deux langues]
« Je peux t'aider », dit-elle. → \eng{“I can help you,” she says.} → \eng{She says she can help me.}\\
Quand le verbe introducteur est au présent (\eng{says}), rien ne change dans la parole rapportée — exactement comme en français : « elle dit qu'elle \emph{peut} m'aider ».
\end{exemplebox}

\courssub{Différence 1 : « que » obligatoire, « that » facultatif}

\begin{propriete}[That : à dire ou à ne pas dire ?]
En français, la conjonction « que » est \textbf{obligatoire} : « il a dit qu'il était fatigué » (on ne peut pas dire « il a dit il était fatigué »).\\
En anglais, \cle{that} est \textbf{facultatif} après \eng{say}, \eng{tell}, \eng{think}, \eng{know}, etc. :\\
\eng{He said (that) he was tired.} — « Il a dit qu'il était fatigué. »\\
À l'écrit formel (examen, dissertation), on préfère souvent garder \eng{that} ; à l'oral, on l'omet fréquemment.
\end{propriete}

\begin{attention}[Ne mets jamais « that » dans une question rapportée]
Erreur classique du francophone : traduire « que » partout. Mais dans les \textbf{questions rapportées}, il n'y a pas de « que » en français... et donc pas de \eng{that} en anglais !\\
✗ \eng{He asked that if I was ready.}\\
✓ \eng{He asked if I was ready.} — « Il a demandé si j'étais prêt. »\\
✗ \eng{She wanted to know that where I lived.}\\
✓ \eng{She wanted to know where I lived.} — « Elle voulait savoir où j'habitais. »
\end{attention}

\courssub{Différence 2 : l'infinitif après « demander à quelqu'un de »}

\begin{propriete}[Ask + objet + to-infinitive]
Le français dit : « le professeur a demandé aux élèves \textbf{d'ouvrir} leurs cahiers » (infinitif direct). L'anglais exige la structure \cle{ask + complément d'objet + to-infinitive} :\\
\eng{The teacher asked the students to open their notebooks.}\\
Même logique avec \eng{tell}, \eng{advise}, \eng{warn}, \eng{order}, \eng{invite} :\\
\eng{She told me to be careful.} — « Elle m'a dit de faire attention. »\\
\eng{The doctor advised him to stop smoking.} — « Le médecin lui a conseillé d'arrêter de fumer. »\\
Négation : \eng{He warned us \textbf{not to touch} the wires.} — « Il nous a avertis de ne pas toucher les fils. »
\end{propriete}

\courssub{Différence 3 : « il m'a dit que oui / que non »}

En français, on dit « il m'a répondu que oui », « elle a dit que non ». En anglais, on ne dit jamais \eng{he said that yes}. Les formes correctes sont :

\begin{center}
\begin{tabularx}{\linewidth}{|X|X|X|}
\hline
\rowcolor{popblueL} Français & English & Remarque \\
\hline
Il a dit que oui. & \eng{He said yes.} & pas de \eng{that} \\
\hline
Elle a répondu que non. & \eng{She said no.} / \eng{She answered no.} & \\
\hline
Il m'a dit que oui. & \eng{He told me so.} & \eng{so} = « cela, oui » \\
\hline
Je ne pense pas. & \eng{I don't think so.} & jamais \eng{I don't think} seul \\
\hline
J'espère que oui. & \eng{I hope so.} & \\
\hline
\end{tabularx}
\end{center}

\courssub{Piège : « demander si » et l'ordre des mots}

Traduire « demander si » par \eng{ask if} est correct — mais le piège se cache dans l'\textbf{ordre des mots} de la question rapportée : on revient à l'ordre de la phrase affirmative (sujet + verbe), sans inversion et sans auxiliaire \eng{do/does/did}.

\begin{attention}[Question rapportée = ordre affirmatif]
✗ \eng{He asked if was I ready.} (inversion du français « si j'étais prêt » mal reconstituée)\\
✓ \eng{He asked if I was ready.}\\[0.3em]
✗ \eng{She asked where did I live.}\\
✓ \eng{She asked where I lived.} — « Elle a demandé où j'habitais. »\\[0.3em]
Retiens : après \eng{if}, \eng{whether} ou un mot interrogatif (\eng{where}, \eng{when}, \eng{why}, \eng{how}), l'anglais suit l'ordre \textbf{sujet + verbe}, comme en français.
\end{attention}

\begin{attention}[Tell ou say ? Le français « dire » sert aux deux]
En français, « dire » convient partout ; en anglais, il faut choisir :\\
\cle{say} (+ \eng{to} + personne, souvent omise) : \eng{He said (to me) that he was tired.}\\
\cle{tell} + \textbf{personne obligatoirement exprimée} : \eng{He told me that he was tired.}\\
✗ \eng{He told that he was tired.} / ✗ \eng{He said me that he was tired.}\\
✓ \eng{He told me...} / \eng{He said (that)...} / \eng{He said to me...}
\end{attention}

\begin{popfigure}
\begin{center}
\begin{tikzpicture}[font=\small]
\fill[popblueL, rounded corners=4pt] (0,0) rectangle (6.4,4.6);
\fill[poporangeL, rounded corners=4pt] (6.8,0) rectangle (13.2,4.6);
\node[align=center, font=\bfseries] at (3.2,4.15) {Français};
\node[align=center, font=\bfseries] at (10,4.15) {English};
\node[align=left, text width=5.6cm] at (3.2,2.0) {Il dit qu'il est fatigué.\\[0.4em] Il a dit qu'il \textbf{était} fatigué.\\[0.4em] Elle m'a dit de venir.\\[0.4em] Il a demandé si j'étais prêt.};
\node[align=left, text width=5.6cm] at (10,2.0) {\eng{He says (that) he is tired.}\\[0.4em] \eng{He said he \textbf{was} tired.}\\[0.4em] \eng{She told me to come.}\\[0.4em] \eng{He asked if I was ready.}};
\draw[-{Stealth}, thick, poppurple] (6.5,3.35) -- (6.75,3.35);
\draw[-{Stealth}, thick, poppurple] (6.5,2.55) -- (6.75,2.55);
\draw[-{Stealth}, thick, poppurple] (6.5,1.75) -- (6.75,1.75);
\draw[-{Stealth}, thick, poppurple] (6.5,0.95) -- (6.75,0.95);
\end{tikzpicture}
\end{center}
\legende{Correspondances français/anglais : mêmes reculs de temps, mais \eng{that} facultatif, \eng{tell} + objet + \eng{to}-infinitive, et ordre affirmatif après \eng{if}.}
\end{popfigure}

\courssub{Exercice guidé pas à pas : transformer un mini-dialogue}

\begin{methode}[La fiche de révision en 5 étapes]
\begin{enumerate}
\item \textbf{Souligne le verbe introducteur} et note son temps : présent (\eng{says}) → pas de changement ; passé (\eng{said}) → backshift.
\item \textbf{Identifie le type de phrase} : déclaration (\eng{say/tell + that}), question fermée (\eng{ask + if/whether}), question ouverte (\eng{ask +} mot interrogatif), ordre ou conseil (\eng{tell/ask/advise + objet + to}-infinitive).
\item \textbf{Recule les temps} (present → past, \eng{will} → \eng{would}, \eng{have} → \eng{had}) et adapte les pronoms et possessifs.
\item \textbf{Adapte les marqueurs de temps et de lieu} : \eng{now} → \eng{then}, \eng{today} → \eng{that day}, \eng{here} → \eng{there}, \eng{tomorrow} → \eng{the next day}.
\item \textbf{Relis} en vérifiant : ordre des mots affirmatif dans les questions, présence de l'objet après \eng{tell}, et concordance avec le français.
\end{enumerate}
\end{methode}

\begin{exoresolu}[Transformation complète d'un mini-dialogue]
Énoncé — Report this dialogue in indirect speech, starting with the verbs given.\\
\eng{At the health centre in Buea, the nurse said to Peter: “Are you feeling better today?” Peter answered: “I still have a headache, but I will take my medicine now.” The nurse said: “Drink a lot of water and don't stay in the sun.”}\\[0.3em]
\tcblower
Solution détaillée, étape par étape :\\[0.3em]
\textbf{Étape 1} : tous les verbes introducteurs sont au passé (\eng{said}, \eng{answered}) → backshift partout.\\
\textbf{Étape 2} : phrase 1 = question fermée → \eng{ask + if} ; phrase 2 = déclaration → \eng{answer + that} ; phrase 3 = ordres → \eng{tell + objet + to}-infinitive.\\
\textbf{Étapes 3 et 4} : \eng{are you feeling} → \eng{if he was feeling} ; \eng{today} → \eng{that day} ; \eng{have} → \eng{had} ; \eng{will take} → \eng{would take} ; \eng{now} → \eng{then}.\\[0.3em]
\textbf{Version rapportée} :\\
\eng{At the health centre in Buea, the nurse asked Peter if he was feeling better that day. Peter answered that he still had a headache, but that he would take his medicine then. The nurse told him to drink a lot of water and not to stay in the sun.}\\[0.3em]
\textbf{Étape 5 — relecture} : ordre affirmatif après \eng{if} ✓ ; objet \eng{him} après \eng{told} ✓ ; négation \eng{not to stay} ✓ ; concordance avec le français (« elle lui demanda s'il se sentait mieux ce jour-là ») ✓.
\end{exoresolu}

\courssub{Repérage d'erreurs : phrases d'élèves francophones}

\begin{exercice}[Trouve et corrige l'erreur]
Chaque phrase ci-dessous contient une erreur typique d'élève francophone. Corrige-la.
\begin{enumerate}
\item \eng{He told that he was sick.}
\item \eng{She asked me where did I come from.}
\item \eng{The teacher said us to open our books.}
\item \eng{He asked that if I was ready.}
\item \eng{My mother said me that she was proud.}
\item \eng{He said that yes, he would come.}
\end{enumerate}
\end{exercice}

\begin{corrige}[Exercice — repérage d'erreurs]
\begin{enumerate}
\item \eng{He told \textbf{me} that he was sick.} — \eng{tell} exige un complément de personne (ou utilise \eng{said}).
\item \eng{She asked me where I \textbf{came} from.} — ordre affirmatif, pas d'auxiliaire \eng{did}.
\item \eng{The teacher \textbf{told} us to open our books.} — \eng{say} ne prend jamais d'objet personne direct ; ici il faut \eng{tell} + objet + \eng{to}-infinitive.
\item \eng{He asked if I was ready.} — pas de \eng{that} dans une question rapportée.
\item \eng{My mother \textbf{told} me that she was proud.} (ou \eng{said to me}) — jamais \eng{said me}.
\item \eng{He said \textbf{yes}, he would come.} (ou \eng{He told me so.}) — jamais \eng{said that yes}.
\end{enumerate}
\end{corrige}

\courssub{Synthèse finale : ta fiche de révision}

\begin{aretenir}[Direct → indirect : l'essentiel en cinq points]
\begin{enumerate}
\item Verbe au \textbf{présent} (\eng{says}) → aucun changement de temps. Verbe au \textbf{passé} (\eng{said}) → backshift : present → past, past/present perfect → past perfect, \eng{will} → \eng{would}, \eng{can} → \eng{could}.
\item Déclaration : \eng{say (that)} / \eng{tell + objet (that)}. \eng{That} est facultatif, « que » est obligatoire.
\item Question fermée : \eng{ask + if/whether} + ordre affirmatif. Question ouverte : \eng{ask +} mot interrogatif + ordre affirmatif. Jamais de \eng{that}, jamais d'inversion.
\item Ordre/conseil : \eng{tell/ask/advise/warn + objet + to}-infinitive ; négation : \eng{not to} + verbe.
\item Adapte pronoms, possessifs et marqueurs : \eng{now} → \eng{then}, \eng{today} → \eng{that day}, \eng{here} → \eng{there}, \eng{this} → \eng{that}, \eng{ago} → \eng{before}.
\end{enumerate}
En cas de doute, traduis en français correct et vérifie la concordance des temps : les deux langues reculent ensemble.
\end{aretenir}

\begin{experience}[Entraînement guidé en binômes]
Par groupes de deux : l'élève A joue un journaliste qui interviewe un élève de ta classe sur ses habitudes de santé (sommeil, sport, alimentation) et pose trois questions directes. L'élève B écoute, puis rapporte l'interview à la classe au discours indirect en appliquant la fiche en 5 étapes : \eng{The journalist asked him if he slept enough. He answered that he usually went to bed late because he watched football...} La classe vérifie chaque phrase avec la fiche : verbe introducteur, type de phrase, backshift, pronoms, marqueurs de temps.
\end{experience}

\newpage

\coursec{Activité d'intégration}

\coursec{Lesson 29 — Grammar: Direct and Indirect Speech}

\courssub{Objectifs de la leçon}
\begin{itemize}
    \item Maîtriser les règles du \cle{discours indirect} (backshift, pronoms, marqueurs spatio-temporels, verbes introducteurs) pour transformer un énoncé du style direct au style indirect.
    \item Passer de la compréhension orale/écrite (conférence de presse, questions-réponses) à la production écrite (article de presse) et orale (restitution).
    \item Identifier et corriger les erreurs fréquentes des francophones (oubli du backshift, mauvais choix de pronom, verbe introducteur inadapté, ordre des mots dans la clause indirecte).
    \item Varier les verbes introducteurs (\eng{announce, state, reveal, promise, advise, warn, urge, explain, confirm, add}) et adapter la structure qui suit (infinitif, \eng{that}-clause, gérondif).
\end{itemize}

\courssub{Situation : Conférence de presse au Lycée Bilingue de Buea}
\begin{textebox}[Script de la conférence de presse — Salle des Actes]
The Principal, Mr. John Ngwa, stands at the podium. Microphones from CRTV and Cameroon Tribune are directed at him. He looks serious but confident.

\textbf{Principal Ngwa:} “Good morning, distinguished journalists, dear students, and members of the teaching staff. I have three major announcements to make today regarding the welfare of our students and the upcoming examination session.

First, I am happy to announce that the school canteen will reopen next Monday after six months of closure for renovation. The new kitchen meets all hygiene standards. A nutritionist has designed the menus. We will serve balanced meals at a subsidised price of 300 FCFA per plate. I want every student to benefit from this.

Second, the new computer laboratory, funded by the Ministry of Secondary Education and our Partner NGO ‘TechForAll’, will be operational from the 15th of next month. It contains fifty new desktop computers with high-speed internet access. Computer science teachers will organise mandatory initiation sessions for all Form Five and Upper Sixth students. This is a great opportunity for you to acquire digital skills before university.

Third, concerning the GCE Advanced Level examinations: the written papers will start on the 3rd of June. Practical examinations for science subjects begin two weeks earlier. I strongly advise all candidates to collect their timetables from the Exams Officer by this Friday. Do not wait until the last minute. I warned the previous batch about late registration; some missed their papers. Learn from their mistakes.

Finally, I urge parents to pay the PTA levies before the end of this term. The funds are needed to complete the perimeter fence. I promised the PTA chairman last week that the administration would enforce this deadline.

Thank you. I will now take three questions.”
\end{textebox}

\begin{textebox}[Extrait de la séance de questions-réponses — Journalistes élèves]
\textbf{Journalist 1 (Cameroon Tribune):} “Sir, you said the canteen reopens next Monday. Does that mean the menu is fixed for the whole term?”

\textbf{Principal Ngwa:} “The nutritionist will review the menu monthly according to seasonal produce. We announced the price today; it will not change this term.”

\textbf{Journalist 2 (CRTV):} “Concerning the computer lab, will Form Three students have access?”

\textbf{Principal Ngwa:} “Priority is given to examination classes. If slots remain available, Form Three students may apply through their form masters next term.”

\textbf{Journalist 3 (School Magazine):} “You advised candidates to collect timetables by Friday. What happens if they don't?”

\textbf{Principal Ngwa:} “They risk missing crucial updates on venue changes. The Exams Officer told me yesterday that the timetable is provisional.”
\end{textebox}

\courssub{Consignes de l'activité intégrée (2 h)}
\begin{enumerate}
    \item \textbf{Préparation (10 min) :} En groupes de 4 à 5 élèves. Désignez un \emph{secrétaire} (noteur), un \emph{porte-parole} (lecteur final), deux \emph{journalistes} (spécialisés : cantine/informatique, examens/discipline) et un \emph{correcteur} (vérificateur grammatical). Lisez les deux documents ci-dessus. Soulignez tous les verbes introducteurs potentiels (\eng{announce, advise, warn, promise, urge, explain, tell, say}) et les marqueurs de temps/lieu à reculer (\eng{next Monday, next month, this Friday, yesterday, today, this term, last week, two weeks earlier}).
    
    \item \textbf{Prise de notes / Écoute active (15 min) :} L'enseignant (ou un élève volontaire) lit le script de la conférence \textbf{une seule fois} à vitesse normale. Les « journalistes » prennent des notes en \textbf{discours direct} (citations clés) sous forme de mots-clés. Exemple : \eng{Canteen reopens Monday — 300 FCFA — subsidised.}
    
    \item \textbf{Transformation grammaticale (25 min) :} Chaque groupe rédige un \textbf{article de presse de 8 à 10 lignes} (environ 100-120 mots) pour le \eng{Cameroon Tribune} ou la \eng{CRTV Online}. L'article doit :
    \begin{itemize}
        \item Avoir un titre accrocheur (ex. \eng{“Principal Unveils Major Reforms Ahead of GCE Exams”}).
        \item Rapporter les trois annonces principales \textbf{au discours indirect}.
        \item Utiliser \textbf{au moins 5 verbes introducteurs différents} parmi : \eng{announced, stated, revealed, promised, advised, warned, urged, explained, added, confirmed}.
        \item Appliquer correctement le \cle{backshift} (présent $\rightarrow$ prétérit, prétérit/present perfect $\rightarrow$ plus-que-parfait, \eng{will} $\rightarrow$ \eng{would}, \eng{can} $\rightarrow$ \eng{could}, \eng{must} $\rightarrow$ \eng{had to}).
        \item Adapter pronoms (\eng{I} $\rightarrow$ \eng{he}, \eng{we} $\rightarrow$ \eng{they}, \eng{our} $\rightarrow$ \eng{their}, \eng{you} $\rightarrow$ \eng{students/candidates}) et marqueurs spatio-temporels (\eng{next Monday} $\rightarrow$ \eng{the following Monday}, \eng{this Friday} $\rightarrow$ \eng{that Friday}, \eng{yesterday} $\rightarrow$ \eng{the previous day}, \eng{last week} $\rightarrow$ \eng{the previous week}).
    \end{itemize}
    
    \item \textbf{Révision par les pairs / Auto-correction (15 min) :} Échangez votre brouillon avec un groupe voisin. Le \emph{correcteur} de chaque groupe utilise la grille ci-dessous (Ressources) pour cocher les points validés et entourer les erreurs. Retournez le papier avec un commentaire constructif.
    
    \item \textbf{Production finale et Restitution orale (15 min) :} Finalisez l'article. Le \emph{porte-parole} le lit à haute voix. La classe note les éventuelles erreurs de \cle{backshift} ou de concordance des temps entendues. L'enseignant fait une correction collective au tableau sur 2-3 phrases types.
\end{enumerate}

\courssub{Ressources grammaticales à mobiliser}

\begin{propriete}[Règle de base : Le Backshift (Recul des temps)]
Si le verbe introducteur est au \textbf{passé} (\eng{said, told, announced, advised, warned, promised}), les temps du discours rapporté reculent d'un cran :
\begin{center}
\begin{tabularx}{\linewidth}{|X|X|X|}
\hline
\rowcolor{popblueL} \textbf{Discours Direct} & \textbf{Discours Indirect} & \textbf{Exemple (Contexte)} \\
\hline
Present Simple (\eng{reopens}) & Past Simple (\eng{reopened}) & He \eng{announced} the canteen \eng{reopened} the following Monday. \\
\hline
Present Continuous (\eng{is designing}) & Past Continuous (\eng{was designing}) & He \eng{said} a nutritionist \eng{was designing} the menus. \\
\hline
Present Perfect (\eng{has designed}) / Past Simple (\eng{designed}) & Past Perfect (\eng{had designed}) & He \eng{explained} a nutritionist \eng{had designed} the menus. \\
\hline
\eng{will} (\eng{will serve}) & \eng{would} (\eng{would serve}) & He \eng{promised} they \eng{would} serve balanced meals. \\
\hline
\eng{can} (\eng{can acquire}) & \eng{could} (\eng{could acquire}) & He \eng{stated} students \eng{could} acquire digital skills. \\
\hline
\eng{must} / \eng{have to} (\eng{must collect}) & \eng{had to} (\eng{had to collect}) & He \eng{advised} candidates they \eng{had to} collect timetables. \\
\hline
\eng{may} (\eng{may apply}) & \eng{might} (\eng{might apply}) & He \eng{added} Form Three students \eng{might} apply later. \\
\hline
\end{tabularx}
\end{center}
\textbf{Exceptions :} Pas de backshift si le fait est \textbf{toujours vrai} au moment du rapport (\eng{“The earth \textbf{is} round”} $\rightarrow$ He said the earth \textbf{is} round) ou pour un \textbf{futur très proche/imminent} (style journalistique parfois). Ici, tout est au passé $\rightarrow$ backshift obligatoire.
\end{propriete}

\begin{propriete}[Changements des pronoms et marqueurs spatio-temporels]
\begin{center}
\begin{tabularx}{\linewidth}{|X|X|}
\hline
\rowcolor{popblueL} \textbf{Direct} & \textbf{Indirect (contexte passé)} \\
\hline
\eng{I, we, my, our, us} & \eng{he, they, his, their, them} (selon le sujet introducteur) \\
\hline
\eng{you, your} & \eng{the students, the candidates, their} (selon le destinataire) \\
\hline
\eng{this, these} & \eng{that, those} \\
\hline
\eng{here} & \eng{there} \\
\hline
\eng{now} & \eng{then} \\
\hline
\eng{today, tonight} & \eng{that day, that night} \\
\hline
\eng{tomorrow} & \eng{the next day / the following day} \\
\hline
\eng{next Monday / week / month} & \eng{the following Monday / week / month} \\
\hline
\eng{yesterday} & \eng{the previous day / the day before} \\
\hline
\eng{last week / month / year} & \eng{the previous week / month / year} \\
\hline
\eng{ago} & \eng{before} \\
\hline
\end{tabularx}
\end{center}
\end{propriete}

\begin{propriete}[Verbes introducteurs et structures associées]
Le choix du verbe nuance le sens. Attention à la structure qui suit (infinitif, \eng{that} + indicatif, gérondif, préposition).
\begin{center}
\begin{tabularx}{\linewidth}{|l|X|X|}
\hline
\rowcolor{popblueL} \textbf{Verbe} & \textbf{Structure} & \textbf{Exemple adapté au contexte} \\
\hline
\eng{announce / state / reveal} & \eng{that} + clause & He \eng{announced that} the canteen \eng{would reopen}... \\
\hline
\eng{promise / threaten} & \eng{to} + infinitif / \eng{that} + \eng{would} & He \eng{promised to enforce} the deadline. / He \eng{promised that} the price \eng{would not change}. \\
\hline
\eng{advise / warn / urge} & \eng{Object + to} + infinitif / \eng{that} + \eng{should} / \eng{subjunctive} & He \eng{advised candidates to collect}... / He \eng{warned them that} they \eng{risked missing}... \\
\hline
\eng{explain / confirm / add} & \eng{that} + clause & He \eng{explained that} the menu \eng{would be reviewed} monthly. \\
\hline
\eng{tell} & \eng{Object + that} / \eng{Object + to} + infinitif & He \eng{told journalists that}... / He \eng{told them to wait}. \\
\hline
\end{tabularx}
\end{center}
\end{propriete}

\begin{attention}[Pièges fréquents pour francophones]
\begin{itemize}
    \item \textbf{Oubli du backshift :} \eng{*He announced the canteen reopens Monday.} $\rightarrow$ \eng{He announced the canteen \textbf{would reopen} the following Monday.}
    \item \textbf{Mauvaise conversion de \eng{must} :} \eng{*He said they must collect.} $\rightarrow$ \eng{He said they \textbf{had to} collect.} (\eng{must} n'a pas de passé ; \eng{had to} remplace \eng{must} au passé).
    \item \textbf{Pronom \eng{you} non adapté :} \eng{*He advised you to collect.} $\rightarrow$ \eng{He advised \textbf{the candidates / them} to collect.}
    \item \textbf{Marqueurs de temps inchangés :} \eng{*He said next Monday.} $\rightarrow$ \eng{He said \textbf{the following Monday}.}
    \item \textbf{Verbe \eng{say} vs \eng{tell} :} \eng{*He said them...} $\rightarrow$ \eng{He \textbf{told} them...} / \eng{He \textbf{said that}...}
    \item \textbf{Ordre des mots dans la clause indirecte :} \eng{*He asked when will the exams start.} $\rightarrow$ \eng{He asked when the exams \textbf{would start}.} (Sujet + Verbe, pas d'inversion, pas d'auxiliaire \eng{do/does/did}).
\end{itemize}
\end{attention}

\courssub{Proposition de production et corrigé commenté}
\begin{exoresolu}[Article modèle — Cameroon Tribune / 9 lignes]
\textbf{Principal Unveils Major Reforms Ahead of GCE Exams}

\textbf{Buea, May 15.} The Principal of the Government Bilingual High School Buea, Mr. John Ngwa, \textbf{announced} yesterday \textbf{that} the school canteen \textbf{would reopen} \textbf{the following Monday} after six months of closure. He \textbf{stated that} a nutritionist \textbf{had designed} balanced menus and \textbf{promised that} the price \textbf{would not change} during the term, remaining at 300 FCFA per plate.

Concerning digital infrastructure, he \textbf{revealed that} the new computer laboratory, equipped with fifty computers, \textbf{would become} operational \textbf{the following month}. He \textbf{explained that} priority \textbf{would be given} to examination classes, adding \textbf{that} Form Three students \textbf{might apply} later if places \textbf{were} available.

Regarding the GCE Advanced Level session, Mr. Ngwa \textbf{advised candidates to collect} their provisional timetables \textbf{by that Friday} and \textbf{warned them that} latecomers \textbf{risked missing} crucial updates. He \textbf{urged parents to settle} PTA levies before the end of term, \textbf{confirming that} the funds \textbf{were needed} for the perimeter fence.
\tcblower
\textbf{Commentaire détaillé des choix linguistiques :}
\begin{enumerate}
    \item \textbf{Backshift systématique :} \eng{reopens $\rightarrow$ would reopen} (futur rapporté) ; \eng{has designed $\rightarrow$ had designed} (antériorité) ; \eng{will become $\rightarrow$ would become} ; \eng{will be given $\rightarrow$ would be given} (passif) ; \eng{are needed $\rightarrow$ were needed} (présent $\rightarrow$ prétérit). Le \eng{past perfect} (\eng{had designed}) marque l'antériorité par rapport au moment de l'annonce.
    \item \textbf{Marqueurs spatio-temporels :} \eng{next Monday $\rightarrow$ the following Monday} ; \eng{next month $\rightarrow$ the following month} ; \eng{this Friday $\rightarrow$ that Friday} ; \eng{yesterday} (dans le lead) ancre l'article dans le passé du journaliste.
    \item \textbf{Pronoms :} \eng{I $\rightarrow$ He (Mr. Ngwa)} ; \eng{We $\rightarrow$ The Principal / He} ; \eng{Our $\rightarrow$ The school / Their} ; \eng{You (candidates) $\rightarrow$ Candidates / Them}.
    \item \textbf{Variété des verbes introducteurs (6 utilisés) :} \eng{announced, stated, promised, revealed, explained, advised, warned, urged, confirmed}. Structures variées : \eng{that}-clauses (\eng{announced that}), infinitif avec objet (\eng{advised candidates to collect}, \eng{urged parents to settle}), clause \eng{that} + verbe (\eng{warned them that... risked missing}).
    \item \textbf{Registre journalistique :} Lead (\textbf{Buea, May 15.}), phrases courtes, vocabulaire précis (\eng{unveils, provisional, perimeter fence, subsidised}), ton neutre et objectif.
\end{enumerate}
\end{exoresolu}

\courssub{Bilan de la leçon}
\begin{aretenir}[Ce qu'il faut retenir de cette intégration]
\begin{itemize}
    \item Le passage du direct à l'indirect n'est pas un exercice mécanique : c'est un \textbf{changement de point de vue} (déictique) qui impose une révision globale de la phrase (temps, pronoms, lieu/temps).
    \item Le \textbf{choix du verbe introducteur} dicte la structure grammaticale qui suit (\eng{advise + to}, \eng{warn + that}, \eng{promise + to/that}). Un bon journaliste varie ces verbes pour éviter la répétition de \eng{said}.
    \item La \textbf{vérification par les pairs} (étape 4) est cruciale : l'erreur la plus fréquente en Terminale est l'oubli partiel du backshift (ex. reculer le verbe principal mais oublier le modal ou le marqueur de temps).
    \item En situation d'examen (Bac), on vous demandera souvent de \textbf{transformer un paragraphe} ou de \textbf{compléter un texte à trous} (cloze test) en discours indirect. La méthode reste identique : 1. Identifier le verbe introducteur (temps). 2. Lister les éléments à reculer. 3. Réécrire phrase par phrase.
\end{itemize}
\end{aretenir}

\begin{methode}[Pour réussir l'épreuve de transformation au Bac]
\begin{enumerate}
    \item \textbf{Surlignez} le verbe introducteur (souvent \eng{said, told, asked, advised, promised} au prétérit).
    \item \textbf{Entourez} dans le texte à transformer : tous les verbes, tous les pronoms (\eng{I, you, we, my, your}), tous les marqueurs (\eng{now, here, this, next, yesterday, tomorrow, ago}).
    \item \textbf{Appliquez} le tableau de correspondance (Backshift + Pronoms + Marqueurs) mot à mot.
    \item \textbf{Relisez} la phrase entière pour vérifier l'ordre des mots (sujet + verbe en clause indirecte) et la cohérence du sens.
\end{enumerate}
\end{methode}

\newpage

\coursec{Méthodes et exercices résolus}

\coursec{Lesson 29 — Grammar: Direct and Indirect Speech}

\courssub{Observation : discours direct et discours indirect}

\begin{textebox}[Observation — Pollution à Douala]
\textbf{Direct speech:} “I am worried about air pollution in Douala,” the student said.\\
\textbf{Indirect speech:} The student said (that) he was worried about air pollution in Douala.
\end{textebox}

\begin{definition}[Discours direct vs discours indirect]
Le \cle{discours direct} cite les mots exacts du locuteur : guillemets anglais “ … ”, deux-points, majuscule au début de la citation, ponctuation finale à l'intérieur des guillemets.\\
Le \cle{discours indirect} (ou rapporté) rapporte le contenu sans le citer mot pour mot : pas de guillemets, subordonnée introduite par \eng{that} (souvent omis à l'oral), adaptation des pronoms, recul des temps (\cle{backshift}) si le verbe introducteur est au passé, adaptation des marqueurs spatio-temporels.
\end{definition}

\courssub{Le backshift : recul des temps verbaux}

\begin{propriete}[Tableau du backshift (verbe introducteur au passé)]
\begin{center}
\begin{tabularx}{\linewidth}{|X|X|X|}\hline
\rowcolor{popblueL} \textbf{Direct speech} & \textbf{Indirect speech} & \textbf{Exemple} \\ \hline
Present Simple & Past Simple & \eng{“I live in Yaoundé.”} $\rightarrow$ \eng{He said he \textbf{lived} in Yaoundé.} \\ \hline
Present Continuous & Past Continuous & \eng{“I am reading.”} $\rightarrow$ \eng{She said she \textbf{was reading}.} \\ \hline
Present Perfect / Past Simple & Past Perfect & \eng{“I have finished.” / “I finished.”} $\rightarrow$ \eng{He said he \textbf{had finished}.} \\ \hline
Past Continuous & Past Perfect Continuous & \eng{“I was sleeping.”} $\rightarrow$ \eng{She said she \textbf{had been sleeping}.} \\ \hline
\eng{will} & \eng{would} & \eng{“I will come.”} $\rightarrow$ \eng{He said he \textbf{would} come.} \\ \hline
\eng{can} & \eng{could} & \eng{“I can swim.”} $\rightarrow$ \eng{She said she \textbf{could} swim.} \\ \hline
\eng{may} & \eng{might} & \eng{“It may rain.”} $\rightarrow$ \eng{He said it \textbf{might} rain.} \\ \hline
\eng{must} & \eng{had to} & \eng{“I must go.”} $\rightarrow$ \eng{She said she \textbf{had to} go.} \\ \hline
\eng{shall} (futur / suggestion) & \eng{should} / \eng{would} & \eng{“Shall I help?”} $\rightarrow$ \eng{He asked if he \textbf{should} help.} \\ \hline
\end{tabularx}
\end{center}
\end{propriete}

\begin{attention}[Exceptions au backshift — pas de recul nécessaire]
\begin{enumerate}
\item \cle{Vérités générales / faits permanents} : \eng{“The earth \textbf{revolves} around the sun.”} $\rightarrow$ \eng{The teacher said the earth \textbf{revolves} around the sun.}
\item \cle{Verbe introducteur au présent / futur / present perfect} : \eng{He \textbf{says} he \textbf{is} tired.} / \eng{He \textbf{will say} he \textbf{is} tired.} / \eng{He \textbf{has said} he \textbf{is} tired.}
\item \cle{Modaux sans équivalent passé ou déjà « passés »} : \eng{could, would, should, might, ought to, used to} ne changent pas. \eng{“I \textbf{could} help.”} $\rightarrow$ \eng{He said he \textbf{could} help.}
\item \cle{Report immédiat} (quelques secondes après) : le recul est optionnel.
\end{enumerate}
\end{attention}

\courssub{Identifier et construire le discours rapporté (déclarations)}

\begin{methode}[Passer du discours direct au discours indirect — Déclarations]
Pour rapporter une déclaration (\emph{statement}), suis toujours la même démarche en cinq étapes.
\begin{enumerate}
\item \cle{Identifier le verbe introducteur} : \eng{say} (sans complément de personne) ou \eng{tell} + complément de personne. \eng{He said (that)...} / \eng{He told me (that)...}
\item \cle{Supprimer les guillemets} et introduire la subordonnée par \eng{that} (facultatif à l'oral, recommandé à l'écrit).
\item \cle{Adapter les pronoms} selon la logique de la situation : \eng{I} devient \eng{he/she}, \eng{my} devient \eng{his/her}, \eng{we} devient \eng{they}, \eng{our} devient \eng{their}.
\item \cle{Appliquer le backshift} (recul des temps) si le verbe introducteur est au passé : voir tableau ci-dessus.
\item \cle{Adapter les marqueurs de temps et de lieu} : \eng{now} $\rightarrow$ \eng{then} ; \eng{today} $\rightarrow$ \eng{that day} ; \eng{yesterday} $\rightarrow$ \eng{the day before} / \eng{the previous day} ; \eng{tomorrow} $\rightarrow$ \eng{the next day} / \eng{the following day} ; \eng{here} $\rightarrow$ \eng{there} ; \eng{this/these} $\rightarrow$ \eng{that/those} ; \eng{ago} $\rightarrow$ \eng{before}.
\end{enumerate}
Réflexe d'auto-contrôle : relis la phrase rapportée et vérifie qu'il n'y a ni guillemets, ni deux-points, ni intonation interrogative.
\end{methode}

\begin{exoresolu}[Rapporter des déclarations — type Bac]
Report the following statements in indirect speech. Start with the words given.

1. “I am worried about air pollution in Douala,” the student said. $\rightarrow$ The student said (that) \ldots

2. “We have planted two hundred trees near the school,” the teacher told the journalist. $\rightarrow$ The teacher told the journalist (that) \ldots

3. “I will visit the health centre tomorrow,” Mary said. $\rightarrow$ Mary said (that) \ldots
\tcblower
\textbf{Étape 1 — Verbe introducteur.} Les trois verbes (\eng{said}, \eng{told}) sont au passé : le backshift est obligatoire. Attention : \eng{tell} exige un complément de personne (\eng{told the journalist}), \eng{say} n'en prend pas.

\textbf{Étape 2 — Phrase 1.} \eng{I am worried} : present continuous $\rightarrow$ past continuous ; \eng{I} devient \eng{he} (ou \eng{she}).

Réponse : \eng{The student said (that) he was worried about air pollution in Douala.}
« L'élève a dit qu'il était inquiet à propos de la pollution de l'air à Douala. »

\textbf{Étape 3 — Phrase 2.} \eng{have planted} : present perfect $\rightarrow$ past perfect (\eng{had planted}) ; \eng{we} devient \eng{they}.

Réponse : \eng{The teacher told the journalist (that) they had planted two hundred trees near the school.}
« Le professeur a dit au journaliste qu'ils avaient planté deux cents arbres près de l'école. »

\textbf{Étape 4 — Phrase 3.} \eng{will visit} $\rightarrow$ \eng{would visit} ; \eng{tomorrow} $\rightarrow$ \eng{the next day} (ou \eng{the following day}).

Réponse : \eng{Mary said (that) she would visit the health centre the next day.}
« Mary a dit qu'elle visiterait le centre de santé le lendemain. »

\textbf{Conclusion.} Trois transformations systématiques : pronoms, temps (backshift), marqueurs spatio-temporels. Aucun guillemet ne subsiste.
\end{exoresolu}

\courssub{Rapporter des questions}

\begin{methode}[Questions rapportées : yes/no et wh-questions]
\begin{enumerate}
\item \cle{Question fermée} (réponse yes/no) : introduis-la par \eng{if} ou \eng{whether} (plus formel) : \eng{He asked me if I was tired.}
\item \cle{Question ouverte} (wh-word) : garde le mot interrogatif comme connecteur : \eng{She asked where I lived.}
\item \cle{Ordre des mots} : la question rapportée reprend l'ordre de la déclaration (sujet + verbe), \textbf{sans inversion} et \textbf{sans auxiliaire do/does/did} : \eng{He asked where the hospital was} (jamais \eng{where was the hospital}).
\item \cle{Pas de point d'interrogation} à la fin : la phrase principale est déclarative.
\item Applique ensuite le backshift et l'adaptation des pronoms comme pour les déclarations.
\end{enumerate}
Verbes introducteurs typiques : \eng{ask}, \eng{wonder}, \eng{want to know}, \eng{enquire}.
\end{methode}

\begin{exoresolu}[Rapporter des questions — type Bac]
Rewrite these questions in indirect speech.

1. “Do you drink enough water every day?” the doctor asked Peter. $\rightarrow$ The doctor asked Peter \ldots

2. “Where did you buy these vegetables?” she asked me. $\rightarrow$ She asked me \ldots

3. “Why are so many children absent from school?” the headmaster wondered. $\rightarrow$ The headmaster wondered \ldots
\tcblower
\textbf{Phrase 1.} Question fermée (\eng{Do you...?}) $\rightarrow$ connecteur \eng{if} ; suppression de l'auxiliaire \eng{do} ; present simple $\rightarrow$ past simple ; \eng{you} $\rightarrow$ \eng{he} ; \eng{every day} peut rester.

Réponse : \eng{The doctor asked Peter if he drank enough water every day.}
« Le médecin a demandé à Peter s'il buvait assez d'eau chaque jour. »

\textbf{Phrase 2.} Question ouverte : on garde \eng{where} ; suppression de \eng{did} ; past simple $\rightarrow$ past perfect ; \eng{you} $\rightarrow$ \eng{I} ; \eng{these} $\rightarrow$ \eng{those}.

Réponse : \eng{She asked me where I had bought those vegetables.}
« Elle m'a demandé où j'avais acheté ces légumes. »

\textbf{Phrase 3.} On garde \eng{why} ; present continuous $\rightarrow$ past continuous ; ordre sujet + verbe : \eng{so many children were}.

Réponse : \eng{The headmaster wondered why so many children were absent from school.}
« Le directeur s'est demandé pourquoi tant d'enfants étaient absents de l'école. »

\textbf{Conclusion.} Dans une question rapportée, il n'y a jamais d'inversion ni de point d'interrogation : c'est une subordonnée déclarative.
\end{exoresolu}

\courssub{Rapporter des ordres, conseils et demandes}

\begin{methode}[Ordres et conseils rapportés : la structure à infinitif]
\begin{enumerate}
\item \cle{Structure affirmative} : verbe introducteur + complément de personne + \eng{to} + base verbale. \eng{“Open the window,” she said to me} $\rightarrow$ \eng{She told me to open the window.}
\item \cle{Structure négative} : \eng{not to} + base verbale. \eng{“Don't smoke here,” he said} $\rightarrow$ \eng{He told me not to smoke there.}
\item \cle{Choisir le verbe introducteur} selon le sens : \eng{tell} (ordre neutre), \eng{order} (ordre strict), \eng{ask} (demande polie), \eng{advise} (conseil), \eng{warn} (avertissement), \eng{encourage} (encouragement), \eng{invite} (invitation), \eng{beg} (supplication).
\item Pas de backshift sur l'infinitif : la base verbale après \eng{to} reste inchangée. Seuls les pronoms et marqueurs de lieu/temps s'adaptent.
\end{enumerate}
\end{methode}

\begin{exoresolu}[Ordres et conseils rapportés — type Bac]
Report the following using the verbs in brackets.

1. “Don't throw plastic bags on the ground,” the mayor said to the villagers. (warn)

2. “You should eat more fruit and vegetables,” the nurse said to me. (advise)

3. “Please help me carry this water,” the old woman said to the boy. (ask)
\tcblower
\textbf{Phrase 1.} Impératif négatif $\rightarrow$ \eng{not to} + base verbale ; \eng{warn} + complément + \eng{not to throw}.

Réponse : \eng{The mayor warned the villagers not to throw plastic bags on the ground.}
« Le maire a averti les villageois de ne pas jeter de sachets plastiques par terre. »

\textbf{Phrase 2.} Conseil exprimé avec \eng{should} : au discours indirect avec \eng{advise}, on utilise la structure à infinitif, plus naturelle que \eng{should} rapporté.

Réponse : \eng{The nurse advised me to eat more fruit and vegetables.}
« L'infirmière m'a conseillé de manger plus de fruits et légumes. »

\textbf{Phrase 3.} Demande polie (\eng{please}) $\rightarrow$ \eng{ask} + complément + \eng{to help} ; \eng{me} devient \eng{her} ; \eng{this} devient \eng{that}.

Réponse : \eng{The old woman asked the boy to help her carry that water.}
« La vieille dame a demandé au garçon de l'aider à porter cette eau. »

\textbf{Conclusion.} La structure \eng{verbe + complément + (not) to + base verbale} remplace tous les impératifs rapportés ; le choix du verbe introducteur traduit la nuance (ordre, conseil, prière).
\end{exoresolu}

\courssub{Choisir le verbe introducteur selon la nuance de sens}

\begin{methode}[Les verbes de parole et leurs constructions]
Chaque verbe introducteur a une construction propre ; apprends le verbe avec sa structure.
\begin{center}
\begin{tabularx}{\linewidth}{|l|X|}\hline
\rowcolor{popblueL} \textbf{Verbe} & \textbf{Construction} \\ \hline
\eng{say} & \eng{say (that)...} — jamais de complément de personne direct \\ \hline
\eng{tell} & \eng{tell somebody (that)...} — complément obligatoire \\ \hline
\eng{ask} & \eng{ask somebody if / wh-...} ; \eng{ask somebody to do} \\ \hline
\eng{advise / warn / order / encourage} & \eng{+ somebody + (not) to do} \\ \hline
\eng{suggest} & \eng{suggest (that) + sujet + should + BV} ou \eng{suggest + V-ing} \\ \hline
\eng{apologise for / accuse of / insist on} & \eng{+ V-ing} \\ \hline
\eng{promise / refuse / offer / agree} & \eng{+ to + BV} \\ \hline
\eng{exclaim / remark / add / explain} & \eng{+ (that)...} \\ \hline
\eng{deny / admit / regret} & \eng{+ V-ing} (ou \eng{that + sujet + verb}) \\ \hline
\end{tabularx}
\end{center}
Démarche : (1) lis la phrase directe et identifie l'intention (promesse ? reproche ? suggestion ?) ; (2) choisis le verbe correspondant ; (3) applique la construction exacte du verbe ; (4) adapte pronoms et temps.
\end{methode}

\begin{exoresolu}[Verbes introducteurs — type Bac]
Report each sentence using the verb in brackets.

1. “I'll bring the medicine this evening, I promise,” John said. (promise)

2. “You broke the water pump!” the villagers said to the mechanic. (accuse)

3. “Let's clean the market square on Saturday,” the youth leader said. (suggest)
\tcblower
\textbf{Phrase 1.} \eng{promise} se construit avec \eng{to} + base verbale ; \eng{will bring} $\rightarrow$ \eng{to bring} ; \eng{this evening} $\rightarrow$ \eng{that evening}.

Réponse : \eng{John promised to bring the medicine that evening.}
« John a promis d'apporter le médicament ce soir-là. »

\textbf{Phrase 2.} \eng{accuse} se construit : \eng{accuse somebody of + V-ing} ; \eng{broke} $\rightarrow$ \eng{breaking} (le past perfect est implicite dans le gérondif).

Réponse : \eng{The villagers accused the mechanic of breaking the water pump.}
« Les villageois ont accusé le mécanicien d'avoir cassé la pompe à eau. »

\textbf{Phrase 3.} \eng{suggest} ne se construit jamais avec un complément de personne direct : on utilise \eng{suggest + V-ing} ou \eng{suggest (that) + should}.

Réponse : \eng{The youth leader suggested cleaning the market square on Saturday.} (ou : \eng{The youth leader suggested that they should clean the market square on Saturday.})
« Le chef des jeunes a suggéré de nettoyer la place du marché samedi. »

\textbf{Conclusion.} Un verbe introducteur bien choisi rend le discours rapporté plus précis et plus naturel — compétence très valorisée à l'écrit du Bac.
\end{exoresolu}

\courssub{Comparaison avec le français et entraînement guidé}

\begin{methode}[Traduire un discours rapporté français $\rightarrow$ anglais]
\begin{enumerate}
\item Repère le verbe introducteur français : « dire que » $\rightarrow$ \eng{say/tell that} ; « demander si » $\rightarrow$ \eng{ask if} ; « demander de + infinitif » $\rightarrow$ \eng{ask/tell + to + BV}.
\item Applique le backshift anglais, qui est \textbf{plus strict} qu'en français : « il a dit qu'il viendrait » $\rightarrow$ \eng{He said he would come} (jamais \eng{will come} après un verbe au passé).
\item Attention à « demander » : français « demander à qqn de faire » $\rightarrow$ anglais \eng{ask somebody to do} (pas de préposition \eng{to} devant la personne).
\item Vérifie les marqueurs temporels : « hier » $\rightarrow$ \eng{the day before}, « demain » $\rightarrow$ \eng{the next day} dans le récit au passé.
\item Relis en vérifiant l'absence d'inversion dans les questions rapportées.
\end{enumerate}
\end{methode}

\begin{exoresolu}[Traduction guidée — type Bac]
Translate into English.

1. Le médecin a dit qu'il était fatigué parce qu'il avait travaillé toute la nuit.

2. Elle m'a demandé si j'avais déjà visité l'hôpital de Buea.

3. Le professeur nous a dit de ne pas jeter de déchets dans la rivière.
\tcblower
\textbf{Phrase 1.} « a dit » $\rightarrow$ \eng{said} ; « était » : l'anglais recule aussi — mais ici le français a déjà un imparfait, donc past simple \eng{was} ; « avait travaillé » $\rightarrow$ past perfect \eng{had worked} ; « toute la nuit » $\rightarrow$ \eng{all night}.

Réponse : \eng{The doctor said (that) he was tired because he had worked all night.}

\textbf{Phrase 2.} « demandé si » $\rightarrow$ \eng{asked if} ; « j'avais déjà visité » : le plus-que-parfait français correspond au past perfect \eng{had already visited} ; ordre déclaratif, pas d'inversion.

Réponse : \eng{She asked me if I had already visited the Buea hospital.}

\textbf{Phrase 3.} « nous a dit de ne pas jeter » $\rightarrow$ \eng{told us not to throw} ; « déchets » $\rightarrow$ \eng{rubbish} (ou \eng{litter}) ; « dans la rivière » $\rightarrow$ \eng{into the river}.

Réponse : \eng{The teacher told us not to throw rubbish into the river.}

\textbf{Conclusion.} Le français et l'anglais reculent tous deux les temps, mais l'anglais exige en plus : \eng{tell} + complément, \eng{if} pour les questions fermées, et la structure à infinitif pour les ordres.
\end{exoresolu}

\begin{attention}[Les 5 erreurs qui coûtent des points au Bac]
\begin{enumerate}
\item \cle{Oublier le backshift} : \eng{He said he \textbf{will} come} est faux ; écris \eng{He said he \textbf{would} come}. Après un verbe introducteur au passé, le recul des temps est obligatoire.
\item \cle{Confondre say et tell} : on dit \eng{He said that...} mais \eng{He told \textbf{me} that...}. Jamais \eng{He said me} ni \eng{He told that...} sans complément.
\item \cle{Garder l'inversion dans les questions rapportées} : \eng{She asked where \textbf{was} the hospital} est faux ; la bonne forme est \eng{She asked where the hospital \textbf{was}} — ordre sujet + verbe, sans point d'interrogation.
\item \cle{Conserver les marqueurs du discours direct} : \eng{He said he would come \textbf{tomorrow}} est incohérent dans un récit au passé ; écris \eng{the next day}. De même \eng{here} $\rightarrow$ \eng{there}, \eng{this} $\rightarrow$ \eng{that}.
\item \cle{Calquer les constructions françaises} : « demander à quelqu'un de faire » ne se traduit pas par \eng{ask to somebody to do} mais par \eng{ask somebody to do} ; et \eng{suggest} ne prend jamais de complément de personne : \eng{He suggested \textbf{going}} ou \eng{suggested that we should go}, jamais \eng{He suggested me to go}.
\end{enumerate}
\end{attention}

\begin{exercice}[Entraînement complet — Discours rapporté]
Report the following sentences into indirect speech. Choose an appropriate reporting verb where indicated.

1. “I have never seen such a clean market,” the visitor said. (remark)

2. “Don't forget to take your malaria treatment,” the doctor told me. (remind)

3. “Where is the nearest health centre?” the tourist asked. (want to know)

4. “I'll help you carry the water cans,” my brother said. (promise)

5. “You should boil drinking water,” the nurse said. (advise)

6. “Let's organise a cleaning campaign,” the class prefect said. (propose / suggest)

7. “Did you attend the meeting yesterday?” the headmaster asked me.

8. “I am writing a report on water pollution,” the journalist said.
\end{exercice}

\begin{corrige}[Exercice — Discours rapporté]
1. The visitor \textbf{remarked} (that) he had never seen such a clean market.\\
2. The doctor \textbf{reminded} me to take my malaria treatment.\\
3. The tourist \textbf{wanted to know} where the nearest health centre was.\\
4. My brother \textbf{promised} to help me carry the water cans.\\
5. The nurse \textbf{advised} me to boil drinking water.\\
6. The class prefect \textbf{proposed} organising a cleaning campaign. (ou \textbf{suggested} organising / \textbf{suggested that they should organise})\\
7. The headmaster \textbf{asked} me if I had attended the meeting the day before.\\
8. The journalist \textbf{said} (that) he was writing a report on water pollution.
\end{corrige}

\begin{savaistu}[Le discours rapporté dans les médias camerounais]
Dans les journaux comme \emph{The Cameroon Tribune} (édition anglaise) ou \emph{The Post}, les discours des officiels sont souvent rapportés : \eng{The Minister said that the government would increase health funding.} À la radio (CRTV English), on entend : \eng{The correspondent reported that the bridge had collapsed.} Maîtriser le discours indirect est indispensable pour comprendre l'actualité et réussir l'épreuve de \emph{Summary} ou de \emph{Composition} au Bac.
\end{savaistu}

\newpage

\coursec{Exercices}

\courssub{Je vérifie mes connaissances}

\begin{exercice}[QCM : la bonne forme rapportée]
Choisis la bonne réponse parmi les quatre options proposées.
\begin{enumerate}
\item “I am tired,” John said. $\rightarrow$ John said that \ldots
\begin{itemize}
\item[a)] he is tired \quad b)] he was tired \quad c)] he has been tired \quad d)] he be tired
\end{itemize}
\item “Where do you live?” she asked me. $\rightarrow$ She asked me \ldots
\begin{itemize}
\item[a)] where did I live \quad b)] where do I live \quad c)] where I lived \quad d)] where I am living
\end{itemize}
\item “Don't touch the wire!” the teacher told us. $\rightarrow$ The teacher told us \ldots
\begin{itemize}
\item[a)] to not touch the wire \quad b)] not to touch the wire \quad c)] don't touch the wire \quad d)] that we not touch the wire
\end{itemize}
\item “I will come tomorrow,” Mary promised. $\rightarrow$ Mary promised that she would come \ldots
\begin{itemize}
\item[a)] tomorrow \quad b)] the next day \quad c)] today \quad d)] now
\end{itemize}
\item “Have you finished your homework?” my father asked. $\rightarrow$ My father asked \ldots
\begin{itemize}
\item[a)] if had I finished my homework \quad b)] if I have finished my homework \quad c)] whether I had finished my homework \quad d)] that I finished my homework
\end{itemize}
\end{enumerate}
\end{exercice}

\begin{exercice}[Vrai ou faux ? Justifie ta réponse]
Dis si les affirmations suivantes sont vraies ou fausses, puis justifie en citant la règle.
\begin{enumerate}
\item In reported speech, “will” usually becomes “would”.
\item We always keep the question word order in reported questions: \eng{He asked me where did I go.}
\item “Say” is followed by an indirect object: \eng{He said me the truth.}
\item “Here” becomes “there” in indirect speech.
\item Orders are reported with “that” + clause: \eng{She told me that sit down.}
\item If the reporting verb is in the present (\eng{she says}), the tense of the original statement does not change.
\end{enumerate}
\end{exercice}

\begin{exercice}[Vocabulaire : les verbes introducteurs]
Complète chaque phrase avec le verbe introducteur qui convient, choisi dans la liste : \eng{advised — apologised — promised — suggested — warned — refused}.
\begin{enumerate}
\item “I'm sorry I broke your phone,” Peter said. $\rightarrow$ Peter \ldots\ for breaking my phone.
\item “You should see a doctor,” she told me. $\rightarrow$ She \ldots\ me to see a doctor.
\item “I will help you with your essay,” he said. $\rightarrow$ He \ldots\ to help me with my essay.
\item “No, I won't lend you my bicycle,” my brother said. $\rightarrow$ My brother \ldots\ to lend me his bicycle.
\item “Be careful! The road is slippery,” the driver said. $\rightarrow$ The driver \ldots\ us that the road was slippery.
\item “Why don't we plant trees around the school?” the student said. $\rightarrow$ The student \ldots\ planting trees around the school.
\end{enumerate}
\end{exercice}

\begin{exercice}[Conjugaison : le backshift]
Mets le verbe entre parenthèses au temps qui convient dans ce discours rapporté.
\begin{enumerate}
\item She said that she \ldots\ (live) in Bamenda since 2020.
\item He told me that he \ldots\ (already / see) the doctor.
\item They announced that the exams \ldots\ (begin) the following Monday.
\item The nurse explained that the patient \ldots\ (suffer) from malaria for three days.
\item My mother said she \ldots\ (can) not come to the meeting because she was ill.
\item The journalist reported that the government \ldots\ (build) a new hospital the previous year.
\end{enumerate}
\end{exercice}

\courssub{Je m'entraîne}

\begin{exercice}[Traduction]
Traduis les phrases suivantes en anglais, en utilisant le discours indirect.
\begin{enumerate}
\item Le médecin m'a dit que je devais boire beaucoup d'eau.
\item Elle m'a demandé si j'avais fait mes devoirs la veille.
\item Le directeur a annoncé que l'école fermerait le lendemain.
\item Il m'a conseillé de ne pas fumer.
\item Ils ont dit qu'ils étaient en train de nettoyer le marché.
\item Mon père m'a demandé où j'avais passé la nuit.
\end{enumerate}
\end{exercice}

\begin{exercice}[Transformation : du discours direct au discours indirect]
Rapporte les phrases suivantes en commençant par l'expression donnée entre parenthèses.
\begin{enumerate}
\item “I don't eat enough fruit,” the boy said. (The boy admitted that \ldots)
\item “Are you feeling better today?” the nurse asked me. (The nurse asked me \ldots)
\item “Why did you miss the health talk yesterday?” the teacher asked us. (The teacher wanted to know \ldots)
\item “Wash your hands before eating,” my grandmother always tells me. (My grandmother always tells me \ldots)
\item “Don't swim in this river; it is polluted,” the guide warned the tourists. (The guide warned the tourists \ldots)
\item “We have planted five hundred trees this year,” the mayor declared. (The mayor declared that \ldots)
\item “I will visit the new health centre next week,” my aunt said. (My aunt said that \ldots)
\item “Can you help me carry these boxes?” she asked me. (She asked me \ldots)
\end{enumerate}
\end{exercice}

\begin{exercice}[Texte à trous : le compte rendu d'une réunion]
Complète le compte rendu suivant en mettant les verbes entre parenthèses au temps qui convient (backshift) et en adaptant si nécessaire les expressions de temps.

\begin{textebox}[Report of the parents' meeting]
Last Friday, our class teacher called a meeting about students' health. She said that many students \ldots\ (come) to school without eating breakfast. She explained that this \ldots\ (affect) their concentration in class. The school nurse added that she \ldots\ (treat) several cases of malaria the previous month and that parents \ldots\ (must) buy mosquito nets. One parent asked whether the school \ldots\ (can) organise a vaccination campaign. The headmaster replied that he \ldots\ (already / write) to the district health officer and that a team \ldots\ (visit) the school the following term. Finally, the teacher told the parents \ldots\ (not / give) their children too much pocket money for snacks, and she promised that she \ldots\ (send) them a written report the next day.
\end{textebox}
\end{exercice}

\begin{exercice}[Repérage et correction d'erreurs]
Chacune des phrases suivantes contient une erreur typique de francophone. Souligne l'erreur et réécris la phrase correctement.
\begin{enumerate}
\item He said me that he was tired.
\item She asked me where did I live.
\item The doctor told me to not smoke.
\item My friend said he will come tomorrow.
\item They asked me if I have seen the film.
\item The teacher said us to open our books.
\item She told that she was happy.
\item He asked me what I am doing there now.
\end{enumerate}
\end{exercice}

\courssub{Je me prépare au Bac}

\begin{exercice}[Type Bac : compréhension et langue]
Lis le texte suivant, puis réponds aux questions.

\begin{textebox}[A health campaign in Buea]
Last month, the mayor of Buea launched a clean-up campaign in the town. Speaking to journalists, he declared: “Our town is becoming too dirty, and diseases are spreading. We must act now.” He promised that the council would provide free dustbins to every neighbourhood and added that volunteers had already cleared the gutters along the main market road.

A nurse from the regional hospital told the crowd: “Cholera cases increased last year because people drank untreated water.” She advised families to boil their drinking water and warned them not to throw rubbish into streams.

One young trader asked the mayor whether the market would get public toilets. The mayor replied that construction would start the following month. “If everyone plays their part,” he concluded, “Buea will be the cleanest town in the region.”
\end{textebox}

\textbf{Comprehension questions}
\begin{enumerate}
\item Why did the mayor launch a clean-up campaign? (Answer in your own words.)
\item What did the nurse advise families to do?
\item What promise did the mayor make about dustbins?
\item What did the young trader want to know?
\end{enumerate}

\textbf{Language work}
\begin{enumerate}
\item Report the following sentences from the text, starting as indicated:
\begin{itemize}
\item “Our town is becoming too dirty.” (The mayor declared that \ldots)
\item “We must act now.” (He said that they \ldots)
\item “Cholera cases increased last year.” (The nurse explained that \ldots)
\item “Don't throw rubbish into streams.” (She warned them \ldots)
\item “Will the market get public toilets?” (The trader asked \ldots)
\end{itemize}
\item Find in the text two expressions of time or place that would change in reported speech, and give their reported form.
\end{enumerate}
\end{exercice}

\begin{exercice}[Type Bac : traduction]
Traduis le passage suivant en anglais.

\begin{textebox}[Texte à traduire]
Hier, le directeur de l'hôpital a annoncé aux journalistes que l'hôpital recevrait de nouveaux équipements le mois suivant. Il a expliqué que le gouvernement avait déjà envoyé les fonds. Une infirmière a demandé si l'hôpital recruterait plus de personnel. Le directeur a répondu qu'il ne pouvait pas répondre à cette question ce jour-là, mais il a promis qu'il donnerait plus d'informations la semaine suivante. Enfin, il a conseillé à la population de se faire vacciner et de ne pas avoir peur des vaccins.
\end{textebox}
\end{exercice}

\begin{exercice}[Type Bac : rédaction guidée]
Tu viens d'avoir un entretien avec le conseiller d'orientation de ton lycée au sujet de tes études supérieures. Rédige, en \textbf{10 à 12 lignes} (environ 100 à 120 mots), un paragraphe au discours indirect dans lequel tu rapportes cet entretien à un camarade. Ton texte doit contenir :
\begin{itemize}
\item au moins trois déclarations rapportées (avec backshift) ;
\item au moins une question fermée rapportée (\eng{if / whether}) ;
\item au moins une question ouverte rapportée (\eng{wh-} question) ;
\item au moins un conseil ou un ordre rapporté (\eng{advise / tell + to}-infinitive) ;
\item au moins deux expressions de temps transformées (\eng{tomorrow} $\rightarrow$ \eng{the next day}, etc.).
\end{itemize}
Commence ton paragraphe par : \eng{Yesterday, I met the guidance counsellor and I asked him about my future studies\ldots}
\end{exercice}

\newpage

\coursec{Corrigés des exercices}

\begin{corrige}[Exercice 1 — QCM : la bonne forme rapportée]
\begin{enumerate}
\item \textbf{b) he was tired.} Le verbe introducteur \eng{said} est au passé : le présent simple \eng{am} recule au prétérit \eng{was} (backshift).
\item \textbf{c) where I lived.} Dans une question rapportée, on retrouve l'ordre sujet + verbe (pas d'inversion) et le présent \eng{live} recule au prétérit \eng{lived}.
\item \textbf{b) not to touch the wire.} Un ordre négatif se rapporte avec \eng{tell + complément + not to} + base verbale. La négation se place \emph{avant} \eng{to}.
\item \textbf{b) the next day.} \eng{Tomorrow} devient \eng{the next day} (ou \eng{the following day}) au discours indirect ; \eng{will} est déjà devenu \eng{would}.
\item \textbf{c) whether I had finished my homework.} Question fermée rapportée avec \eng{if/whether} + ordre sujet-verbe ; le present perfect \eng{have finished} recule au past perfect \eng{had finished}.
\end{enumerate}
\end{corrige}

\begin{corrige}[Exercice 2 — Vrai ou faux ? Justifie ta réponse]
\begin{enumerate}
\item \textbf{Vrai.} Au discours indirect, \eng{will} recule normalement en \eng{would} : \eng{“I will come.”} $\rightarrow$ \eng{He said he would come.}
\item \textbf{Faux.} Dans une question rapportée, on utilise l'ordre de la déclarative (sujet + verbe), sans inversion ni auxiliaire \eng{do/did} : \eng{He asked me where I went.} La phrase donnée est incorrecte.
\item \textbf{Faux.} \eng{Say} ne prend jamais de complément d'objet indirect de personne sans la préposition \eng{to} : on dit \eng{He told me the truth} ou \eng{He said to me that\ldots} La structure \eng{said me} n'existe pas.
\item \textbf{Vrai.} Les expressions de lieu proche s'éloignent au discours rapporté : \eng{here} $\rightarrow$ \eng{there}, \eng{this} $\rightarrow$ \eng{that}.
\item \textbf{Faux.} Les ordres se rapportent avec \eng{tell/order + complément + to}-infinitive : \eng{She told me to sit down.} La structure \eng{that} + proposition ne convient pas pour un ordre.
\item \textbf{Vrai.} Si le verbe introducteur est au présent (\eng{she says}, \eng{he tells me}), il n'y a pas de backshift : \eng{“I am tired.”} $\rightarrow$ \eng{She says she is tired.}
\end{enumerate}
\end{corrige}

\begin{corrige}[Exercice 3 — Vocabulaire : les verbes introducteurs]
\begin{enumerate}
\item Peter \textbf{apologised} for breaking my phone. (\eng{apologise for} + \emph{-ing} : s'excuser de)
\item She \textbf{advised} me to see a doctor. (\eng{advise + complément + to}-infinitive)
\item He \textbf{promised} to help me with my essay. (\eng{promise to} + base verbale)
\item My brother \textbf{refused} to lend me his bicycle. (\eng{refuse to} + base verbale)
\item The driver \textbf{warned} us that the road was slippery. (\eng{warn + complément + that} ; on aurait aussi : \eng{warned us to be careful})
\item The student \textbf{suggested} planting trees around the school. (\eng{suggest} + \emph{-ing} ; jamais \eng{suggest to do})
\end{enumerate}
\end{corrige}

\begin{corrige}[Exercice 4 — Conjugaison : le backshift]
\begin{enumerate}
\item She said that she \textbf{had lived} in Bamenda since 2020. (present perfect $\rightarrow$ past perfect)
\item He told me that he \textbf{had already seen} the doctor. (present perfect $\rightarrow$ past perfect ; \eng{already} peut être conservé)
\item They announced that the exams \textbf{would begin} the following Monday. (\eng{will} $\rightarrow$ \eng{would} ; \eng{next} $\rightarrow$ \eng{the following})
\item The nurse explained that the patient \textbf{had been suffering} (ou \textbf{had suffered}) from malaria for three days. (present perfect continuous $\rightarrow$ past perfect continuous)
\item My mother said she \textbf{could} not come to the meeting because she was ill. (\eng{can} $\rightarrow$ \eng{could})
\item The journalist reported that the government \textbf{had built} a new hospital the previous year. (prétérit $\rightarrow$ past perfect ; \eng{last year} $\rightarrow$ \eng{the previous year})
\end{enumerate}
\end{corrige}

\begin{corrige}[Exercice 5 — Traduction]
\begin{enumerate}
\item \eng{The doctor told me (that) I had to drink a lot of water.} — \eng{must} devient \eng{had to} au discours rapporté.
\item \eng{She asked me if/whether I had done my homework the day before.} — question fermée : \eng{if} + ordre sujet-verbe ; \eng{la veille} = \eng{the day before}.
\item \eng{The headmaster announced (that) the school would close the next day.} — futur français $\rightarrow$ \eng{would} + base verbale.
\item \eng{He advised me not to smoke.} — \eng{advise + complément + not to} + base verbale.
\item \eng{They said (that) they were cleaning the market.} — présent continu $\rightarrow$ prétérit continu (\eng{were cleaning}).
\item \eng{My father asked me where I had spent the night.} — question ouverte rapportée : \eng{where} + sujet + verbe, sans inversion.
\end{enumerate}
\textbf{Points de vigilance~:} ne pas traduire \eng{devoirs} par \eng{duties} (faux ami partiel) : \eng{homework} est invariable ; ne pas oublier le recul des temps même si le français garde l'imparfait ou le conditionnel.
\end{corrige}

\begin{corrige}[Exercice 6 — Transformation : du discours direct au discours indirect]
\begin{enumerate}
\item \eng{The boy admitted that he didn't eat enough fruit.} (présent $\rightarrow$ prétérit ; \eng{I} $\rightarrow$ \eng{he})
\item \eng{The nurse asked me if/whether I was feeling better that day.} (\eng{today} $\rightarrow$ \eng{that day} ; présent continu $\rightarrow$ prétérit continu)
\item \eng{The teacher wanted to know why we had missed the health talk the day before.} (prétérit $\rightarrow$ past perfect ; \eng{yesterday} $\rightarrow$ \eng{the day before})
\item \eng{My grandmother always tells me to wash my hands before eating.} (ordre rapporté : \eng{tell + me + to} ; pas de backshift car \eng{tells} est au présent)
\item \eng{The guide warned the tourists not to swim in that river because it was polluted.} (ordre négatif : \eng{warn + not to} ; \eng{this} $\rightarrow$ \eng{that})
\item \eng{The mayor declared that they had planted five hundred trees that year.} (present perfect $\rightarrow$ past perfect ; \eng{this year} $\rightarrow$ \eng{that year})
\item \eng{My aunt said that she would visit the new health centre the following week.} (\eng{will} $\rightarrow$ \eng{would} ; \eng{next week} $\rightarrow$ \eng{the following week})
\item \eng{She asked me if/whether I could help her carry those boxes.} (\eng{can} $\rightarrow$ \eng{could} ; \eng{these} $\rightarrow$ \eng{those})
\end{enumerate}
\end{corrige}

\begin{corrige}[Exercice 7 — Texte à trous : le compte rendu d'une réunion]
Texte complété :

\eng{Last Friday, our class teacher called a meeting about students' health. She said that many students \textbf{came} to school without eating breakfast. She explained that this \textbf{affected} their concentration in class. The school nurse added that she \textbf{had treated} several cases of malaria the previous month and that parents \textbf{had to} buy mosquito nets. One parent asked whether the school \textbf{could} organise a vaccination campaign. The headmaster replied that he \textbf{had already written} to the district health officer and that a team \textbf{would visit} the school the following term. Finally, the teacher told the parents \textbf{not to give} their children too much pocket money for snacks, and she promised that she \textbf{would send} them a written report the next day.}

\textbf{Justifications~:}
\begin{itemize}
\item \eng{came}, \eng{affected} : présent simple $\rightarrow$ prétérit (vérités observées à un moment passé).
\item \eng{had treated} : present perfect $\rightarrow$ past perfect (action antérieure au moment de la parole).
\item \eng{had to} : \eng{must} n'a pas de prétérit ; on le remplace par \eng{had to}.
\item \eng{could} : \eng{can} $\rightarrow$ \eng{could}.
\item \eng{had already written} : present perfect $\rightarrow$ past perfect.
\item \eng{would visit}, \eng{would send} : \eng{will} $\rightarrow$ \eng{would}.
\item \eng{not to give} : ordre négatif rapporté, \eng{tell + complément + not to} + base verbale.
\end{itemize}
\end{corrige}

\begin{corrige}[Exercice 8 — Repérage et correction d'erreurs]
\begin{enumerate}
\item Erreur : \eng{said me}. Correction : \eng{He \textbf{told me} that he was tired.} (ou \eng{He said to me that\ldots}) — \eng{say} ne prend pas de complément de personne sans \eng{to}.
\item Erreur : \eng{where did I live}. Correction : \eng{She asked me where \textbf{I lived}.} — pas d'inversion dans une question rapportée.
\item Erreur : \eng{to not smoke}. Correction : \eng{The doctor told me \textbf{not to smoke}.} — la négation précède \eng{to}.
\item Erreur : \eng{will come tomorrow}. Correction : \eng{My friend said he \textbf{would come the next day}.} — backshift obligatoire après \eng{said}.
\item Erreur : \eng{have seen}. Correction : \eng{They asked me if I \textbf{had seen} the film.} — present perfect $\rightarrow$ past perfect.
\item Erreur : \eng{said us}. Correction : \eng{The teacher \textbf{told us} to open our books.} — devant un complément + infinitive, seul \eng{tell} convient.
\item Erreur : \eng{told that}. Correction : \eng{She \textbf{said} that she was happy} (ou \eng{She told \textbf{me/us} that\ldots}) — \eng{tell} exige un complément de personne.
\item Erreur : \eng{am doing there now}. Correction : \eng{He asked me what I \textbf{was doing there then}.} — double recul : présent continu $\rightarrow$ prétérit continu et \eng{now} $\rightarrow$ \eng{then}.
\end{enumerate}
\end{corrige}

\begin{corrige}[Exercice 9 — Type Bac : compréhension et langue]
\textbf{Comprehension questions — réponses attendues~:}
\begin{enumerate}
\item \eng{He launched the campaign because the town was becoming too dirty and diseases were spreading.} (réponse avec ses propres mots ; toute formulation équivalente est acceptée)
\item \eng{She advised families to boil their drinking water.}
\item \eng{He promised that the council would provide free dustbins to every neighbourhood.}
\item \eng{He wanted to know whether the market would get public toilets.}
\end{enumerate}

\textbf{Language work~:}
\begin{enumerate}
\item Phrases rapportées~:
\begin{itemize}
\item \eng{The mayor declared that their town was becoming too dirty.} (présent continu $\rightarrow$ prétérit continu)
\item \eng{He said that they had to act then.} (\eng{must} $\rightarrow$ \eng{had to} ; \eng{now} $\rightarrow$ \eng{then})
\item \eng{The nurse explained that cholera cases had increased the previous year.} (prétérit $\rightarrow$ past perfect ; \eng{last year} $\rightarrow$ \eng{the previous year})
\item \eng{She warned them not to throw rubbish into streams.} (ordre négatif : \eng{warn + not to})
\item \eng{The trader asked whether the market would get public toilets.} (question fermée : \eng{whether} + ordre sujet-verbe ; \eng{will} $\rightarrow$ \eng{would})
\end{itemize}
\item Deux expressions qui changeraient : \eng{last month} $\rightarrow$ \eng{the previous month} ; \eng{last year} $\rightarrow$ \eng{the previous year}. (Autres possibilités : \eng{now} $\rightarrow$ \eng{then} ; \eng{the following month} est déjà une forme rapportée de \eng{next month}.)
\end{enumerate}

\textbf{Barème indicatif (sur 20)~:} compréhension 8 pts (2 pts par réponse correcte et bien formulée) ; langue 12 pts (5 $\times$ 2 pts pour les phrases rapportées ; 2 pts pour les expressions de temps). On sanctionne l'absence de backshift, l'inversion dans les questions rapportées et l'oubli de la transformation des expressions de temps.
\end{corrige}

\begin{corrige}[Exercice 10 — Type Bac : traduction]
\textbf{Proposition de traduction~:}

\eng{Yesterday, the hospital director announced to journalists that the hospital would receive new equipment the following month. He explained that the government had already sent the funds. A nurse asked whether the hospital would recruit more staff. The director replied that he could not answer that question that day, but he promised that he would give more information the following week. Finally, he advised the population to get vaccinated and not to be afraid of vaccines.}

\textbf{Points de vigilance~:}
\begin{itemize}
\item \eng{recevrait} (conditionnel français) $\rightarrow$ \eng{would receive} ; ne jamais écrire \eng{will} après un verbe introducteur au passé.
\item \eng{avait envoyé} (plus-que-parfait) $\rightarrow$ past perfect \eng{had sent}.
\item \eng{ce jour-là} $\rightarrow$ \eng{that day} ; \eng{le mois suivant} $\rightarrow$ \eng{the following month} ; \eng{la semaine suivante} $\rightarrow$ \eng{the following week}.
\item \eng{personnel} $\rightarrow$ \eng{staff} (invariable ; jamais \eng{staffs}).
\item \eng{se faire vacciner} $\rightarrow$ \eng{get vaccinated} ou \eng{be vaccinated}.
\item \eng{il a conseillé à la population de\ldots et de ne pas\ldots} $\rightarrow$ \eng{he advised the population to\ldots and not to\ldots} (un seul verbe introducteur pour les deux infinitives).
\end{itemize}
\textbf{Barème indicatif (sur 10)~:} 1 pt par segment correctement traduit (verbe introducteur, backshift, expressions de temps) ; $-0,5$ pt par faute de grammaire ou de vocabulaire.
\end{corrige}

\begin{corrige}[Exercice 11 — Type Bac : rédaction guidée]
\textbf{Proposition de rédaction modèle (environ 110 mots)~:}

\eng{Yesterday, I met the guidance counsellor and I asked him about my future studies. He told me that my results in science subjects were excellent and that I could apply to the faculty of medicine. I asked him whether the entrance examination was difficult, and he answered that it required serious preparation. He also asked me what I wanted to do after my studies. When I said that I dreamed of becoming a doctor, he advised me to work harder in biology and chemistry. He explained that the registration forms would be available the following week and told me not to miss the deadline. Finally, he promised that he would give me a recommendation letter the next day. I left his office feeling much more confident about my future.}

\textbf{Check-list de vérification (à faire relire à l'élève)~:}
\begin{itemize}
\item Trois déclarations rapportées avec backshift : \eng{were}, \eng{could apply}, \eng{would be available}, \eng{would give}.
\item Une question fermée : \eng{I asked him whether the entrance examination was difficult}.
\item Une question ouverte : \eng{He asked me what I wanted to do}.
\item Un conseil et un ordre : \eng{advised me to work harder} ; \eng{told me not to miss the deadline}.
\item Deux expressions de temps transformées : \eng{the following week}, \eng{the next day}.
\end{itemize}
\textbf{Barème indicatif (sur 20)~:} respect des consignes et des contraintes 6 pts ; correction grammaticale (backshift, ordre des mots, verbes introducteurs) 8 pts ; vocabulaire et cohérence 4 pts ; longueur respectée et présentation 2 pts.
\textbf{Points de vigilance~:} éviter l'inversion dans les questions rapportées ; ne pas oublier le complément de personne après \eng{tell} et \eng{advise} ; placer \eng{not} avant \eng{to} dans les ordres négatifs ; vérifier la concordance des temps dans tout le paragraphe.
\end{corrige}

\newpage

\coursec{Fiche bilan}

\begin{popfigure}
\begin{tikzpicture}[
  mindmap,
  grow cyclic,
  every node/.style={concept, font=\small},
  root concept/.append style={concept color=popblue, font=\bfseries\small, minimum size=3.2cm},
  level 1 concept/.append style={level distance=4.2cm, sibling angle=60, font=\scriptsize},
  text=white
]
\node[root concept]{Direct \& Indirect Speech}
  child[concept color=poporange]{ node[align=center]{Observation\\ “I am tired,” he said.\\ He said (that) he was tired.} }
  child[concept color=popgreen]{ node[align=center]{Backshift\\ present $\rightarrow$ past\\ past $\rightarrow$ past perfect\\ will $\rightarrow$ would} }
  child[concept color=poppurple]{ node[align=center]{Pronouns \& adverbs\\ I $\rightarrow$ he/she\\ now $\rightarrow$ then\\ here $\rightarrow$ there} }
  child[concept color=poppink]{ node[align=center]{Questions\\ ask + if/whether\\ ask + wh-word\\ ordre affirmatif} }
  child[concept color=popgold]{ node[align=center]{Orders \& advice\\ tell/order/ask/advise\\ + object + to-infinitive} }
  child[concept color=popblue!70]{ node[align=center]{Reporting verbs\\ say, tell, promise,\\ suggest, admit, deny} };
\end{tikzpicture}
\legende{Carte mentale de la leçon : rapporter les paroles en anglais}
\end{popfigure}

\begin{bilan}
\textbf{1. Lexique clé (English -- Français)}
\begin{itemize}
  \item \eng{direct speech} : discours direct ; \eng{indirect/reported speech} : discours indirect
  \item \eng{reporting verb} : verbe introducteur (\eng{say, tell, ask, advise, promise})
  \item \eng{backshift} : recul des temps ; \eng{statement} : déclaration
  \item \eng{to quote} : citer ; \eng{quotation marks} : guillemets ; \eng{whether} : si (interrogation indirecte)
\end{itemize}

\textbf{2. Règles à connaître par cœur}
\begin{center}
\begin{tabularx}{\linewidth}{|X|X|}\hline
\rowcolor{popblueL} \textbf{Discours direct} & \textbf{Discours indirect} \\ \hline
present simple : \eng{“I work hard.”} & past simple : \eng{He said he worked hard.} \\ \hline
present continuous : \eng{“I am reading.”} & past continuous : \eng{She said she was reading.} \\ \hline
past simple / present perfect & past perfect : \eng{He said he had left / had finished.} \\ \hline
\eng{will} & \eng{would} ; \eng{can} $\rightarrow$ \eng{could} ; \eng{may} $\rightarrow$ \eng{might} \\ \hline
\eng{must} (obligation) & \eng{had to} \\ \hline
\eng{now, today, here, this, ago} & \eng{then, that day, there, that, before} \\ \hline
\eng{tomorrow, yesterday, next week} & \eng{the next day, the day before, the following week} \\ \hline
Question : \eng{“Where do you live?”} & \eng{He asked me where I lived.} (pas d'inversion, pas de \eng{?}) \\ \hline
Question oui/non : \eng{“Are you ill?”} & \eng{She asked if/whether I was ill.} \\ \hline
Ordre : \eng{“Sit down!”} & \eng{He told me to sit down.} / négatif : \eng{not to go} \\ \hline
\end{tabularx}
\end{center}

\textbf{3. Méthode express en 4 étapes}
\begin{enumerate}
  \item Choisir le verbe introducteur (\eng{say} sans objet, \eng{tell} + objet).
  \item Adapter les pronoms et les possessifs à la situation.
  \item Faire reculer le temps (backshift) et changer les adverbes de temps et de lieu.
  \item Pour une question : supprimer l'inversion et le point d'interrogation ; pour un ordre : \eng{tell/ask + objet + to}-infinitif.
\end{enumerate}

\textbf{4. Pas de backshift si} : vérité générale ou fait encore vrai (\eng{He said the earth is round.}), ou verbe introducteur au présent (\eng{She says she is tired.}).
\end{bilan}

\begin{aretenir}[Checklist avant le Bac]
$\square$ Je distingue \eng{say} (sans objet) de \eng{tell} (+ objet : \eng{tell me}).\\
$\square$ Je sais que \eng{that} est facultatif après \eng{said (that)}.\\
$\square$ J'applique le backshift : present $\rightarrow$ past, past $\rightarrow$ past perfect.\\
$\square$ Je transforme \eng{will} en \eng{would}, \eng{can} en \eng{could}, \eng{must} en \eng{had to}.\\
$\square$ Je change les adverbes : \eng{now} $\rightarrow$ \eng{then}, \eng{here} $\rightarrow$ \eng{there}, \eng{ago} $\rightarrow$ \eng{before}.\\
$\square$ Je n'inverse pas le sujet et le verbe dans une question rapportée.\\
$\square$ J'utilise \eng{if} ou \eng{whether} pour les questions oui/non rapportées.\\
$\square$ Je construis les ordres rapportés avec \eng{tell/order/ask + objet + to}-infinitif.\\
$\square$ Je place \eng{not} avant l'infinitif : \eng{He told me not to be late.}\\
$\square$ Je ne fais pas de backshift pour une vérité générale ou un verbe introducteur au présent.\\
$\square$ Je connais des verbes introducteurs variés : \eng{promise, advise, suggest, admit, deny, warn}.\\
$\square$ Je relis ma phrase rapportée pour vérifier la concordance des pronoms.
\end{aretenir}

\findecours

\end{document}
